% La liberté et la responsabilité — Philosophie Tle Tech — Cours signé OnBuch+
% Programme officiel MINESEC. Compiler avec : tectonic N08-liberte-responsabilite.tex
\newcommand{\DOCMATIERE}{Philosophie}
\newcommand{\DOCNIVEAU}{Tle Tech}
\newcommand{\DOCMODULE}{Philosophie – classe de Terminale technique}
\newcommand{\DOCLECON}{N08}
\newcommand{\DOCTITRE}{La liberté et la responsabilité}
% OnBuch+ — préambule « cours signés » ENRICHI (compilé avec tectonic / XeLaTeX)
% Système visuel « Pop » d'OnBuch+ : fond crème (jamais blanc), bordures encre
% épaisses, ombres « dures » décalées, titres Archivo Black en capitales.
% Version MATHÉMATIQUES : graphes (pgfplots/tikz), boîtes pédagogiques
% (définition, propriété, méthode, attention, le savais-tu, exercice résolu,
% exercice, TP, bilan), page de garde et signature OnBuch+ ancrée en bas.
%
% Macros attendues dans main.tex AVANT \input{preamble} :
%   \DOCMATIERE  \DOCNIVEAU  \DOCMODULE  \DOCLECON (numéro)  \DOCTITRE
\documentclass[11pt]{article}

\usepackage{fontspec}
\usepackage[french]{babel}
\usepackage[a4paper,margin=2cm,top=3.4cm,bottom=2.8cm,headheight=1.4cm,footskip=1.3cm]{geometry}
\usepackage{amsmath,amssymb,amsfonts}
\usepackage{mathtools} % \xmapsto, \coloneqq... souvent utilisés par les modèles
\usepackage{cancel} % simplifications barrées (bilans de forces)
% \abs{z} (module, valeur absolue) et \norm{u} (norme), utilisés par les modèles
\providecommand{\abs}{}\let\abs\relax\DeclarePairedDelimiter{\abs}{\lvert}{\rvert}
\providecommand{\norm}{}\let\norm\relax\DeclarePairedDelimiter{\norm}{\lVert}{\rVert}
\usepackage[table]{xcolor} % [table] : fournit aussi \rowcolors (alternance de lignes)
\usepackage{pagecolor}
\usepackage{tikz}
\usetikzlibrary{arrows.meta,positioning,calc,shapes.geometric,shapes.misc,shapes.arrows,decorations.pathmorphing,decorations.pathreplacing,decorations.markings,patterns,shadows,mindmap,trees,fit,backgrounds,matrix,intersections,angles,quotes}
\usepackage{pgfplots}
\usepackage[version=4]{mhchem} % équations nucléaires \ce{^{235}_{92}U}
\usepackage[european,siunitx]{circuitikz} % schémas électriques
\pgfplotsset{compat=1.18}
\usepgfplotslibrary{fillbetween,groupplots} % aires entre courbes (calcul intégral)
\usepackage{enumitem}
\usepackage{fancyhdr}
% [most] charge toutes les bibliothèques tcolorbox (listings, minted, raster,
% documentation...) et gonfle inutilement la mémoire TeX sur un document long
% (~300 boîtes) : on ne charge que ce qui est réellement utilisé.
\usepackage[skins,breakable]{tcolorbox}
\usepackage{graphicx}
\usepackage{colortbl}
\usepackage{booktabs}
\usepackage{tabularx}
\usepackage{multirow}
\usepackage{diagbox} % case d'en-tête diagonale des tableaux à double entrée
% Colonnes extensibles centrée / gauche / droite (C, L, R), souvent utilisées par les modèles
\newcolumntype{C}{>{\centering\arraybackslash}X}
\newcolumntype{L}{>{\raggedright\arraybackslash}X}
\newcolumntype{R}{>{\raggedleft\arraybackslash}X}
\usepackage{array}
\usepackage{multicol}
\usepackage{siunitx}
% parse-numbers=false : accepte aussi \SI{2,5 \times 10^{-3}}{...}
\sisetup{locale=FR,output-decimal-marker={,},group-minimum-digits=4,parse-numbers=false}
\DeclareSIUnit\ohm{\ensuremath{\symup{\Omega}}} % U+2126 absent de la police
\DeclareSIUnit\division{div} % réglages d'oscilloscope (ms/div, V/div)
\DeclareSIUnit\tr{tr} % tours (vitesse de rotation, ex. \SI{1500}{\tr\per\minute})
\DeclareSIUnit\molar{M} % concentration molaire (ex. \SI{3}{\molar}, \SI{1}{\micro\molar})
\DeclareSIUnit\an{an} % année (le modèle confond parfois avec \year, un compteur TeX) — ex. \SI{5}{\kilo\meter\per\an}
\DeclareSIUnit\FCFA{FCFA} % franc CFA (ex. \SI{500}{\FCFA\per\kilo\gram})
\DeclareSIUnit\degreeBrix{\textdegree Brix} % teneur en sucre d'un sirop/jus (réfractométrie)
\DeclareSIUnit\month{mois} % le modèle utilise parfois \month, un compteur TeX, comme unité de durée
\DeclareSIUnit\microliter{\micro\liter} % le modèle écrit parfois \microliter en un seul token
\DeclareSIUnit\million{millions} % dénombrement (ex. \SI{15}{\million\per\mL} spermatozoïdes)
\usepackage{qrcode}
% \degree : commande fréquemment utilisée par les modèles pour un angle ou une
% température en dehors de \SI{}{} (non fournie par les paquets chargés).
\providecommand{\degree}{^{\circ}}
% \qed (fin de démonstration), utilisé par les modèles sans amsthm
\providecommand{\qed}{\hfill$\square$}
% \barre : double barre des valeurs interdites dans un tableau de variations (tkz-tab)
\providecommand{\barre}{\ensuremath{\|}}
% environnement proof (amsthm non chargé) utilisé par les modèles
\ifdefined\proof\else\newenvironment{proof}[1][Démonstration]{\par\noindent\textit{#1.}\ }{\qed\par}\fi
\usepackage{lastpage}
\usepackage{hyperref}
\hypersetup{hidelinks}

% ============================================================
% Palette « Pop » OnBuch+ — mêmes hex que la classe P (plus_ui.dart)
% ============================================================
\definecolor{poporange}{HTML}{F59321}
\definecolor{popdark}{HTML}{E07A0C}
\definecolor{popcream}{HTML}{FFF6E8}
\definecolor{popink}{HTML}{1C1714}
\definecolor{poppurple}{HTML}{7A5AE0}
\definecolor{popgreen}{HTML}{2E9E6A}
\definecolor{poppink}{HTML}{E0457B}
\definecolor{popblue}{HTML}{2F7DE1}
\definecolor{popgold}{HTML}{C98A12}
\definecolor{popmuted}{HTML}{8A7F76}
\definecolor{popline}{HTML}{EADFD2}
\definecolor{popgray}{HTML}{8A7F76}
\colorlet{popgrey}{popgray}
% Alias tolérants (le modèle nomme parfois les couleurs différemment)
\colorlet{popwhite}{white}
\colorlet{popblack}{popink}
\colorlet{popred}{poppink}
\colorlet{popyellow}{popgold}
% Teintes claires (fonds de tableaux/encadrés) — le modèle nomme parfois
% les couleurs avec un suffixe L (ex. popblueL) sans les avoir définies.
\colorlet{poporangeL}{poporange!20}
\colorlet{popdarkL}{popdark!20}
\colorlet{poppurpleL}{poppurple!20}
\colorlet{popgreenL}{popgreen!20}
\colorlet{poppinkL}{poppink!20}
\colorlet{popblueL}{popblue!20}
\colorlet{popgoldL}{popgold!20}
\colorlet{popredL}{popred!20}
\colorlet{popgrayL}{popgray!20}
% Teintes claires pour fonds de tableaux / zones
\colorlet{popblueL}{popblue!10!white}
\colorlet{poporangeL}{poporange!14!white}
\colorlet{popgreenL}{popgreen!12!white}
\colorlet{poppurpleL}{poppurple!12!white}
\colorlet{poppinkL}{poppink!12!white}
\colorlet{popgoldL}{popgold!14!white}

\pagecolor{popcream}

% ============================================================
% Typographie
% ============================================================
\defaultfontfeatures{Ligatures=TeX}
\setmainfont{PlusJakartaSans-Medium.ttf}[Path=../../../fonts/,BoldFont=PlusJakartaSans-Bold.ttf]
% Maths en sans-serif (Fira Math) avec chiffres et lettres droites de Plus
% Jakarta, pour que les formules se fondent dans le texte.
\usepackage[math-style=ISO,bold-style=ISO]{unicode-math}
\setmathfont{FiraMath-Regular.otf}[Path=../../../fonts/]
\setmathfont{PlusJakartaSans-Medium.ttf}[Path=../../../fonts/,range={up/num,up/{Latin,latin},\mathup}]
\usepackage{icomma} % virgule décimale sans espace en mode math
% Caractères mathématiques Unicode absents des polices de texte (Plus Jakarta,
% Archivo Black) : rendus par leur équivalent mathématique, en texte comme en
% formule — sinon ils s'affichent en carré vide (ex. « Arithmétique dans ℤ »).
\usepackage{newunicodechar}
\newunicodechar{ℝ}{\ensuremath{\mathbb{R}}}
\newunicodechar{ℤ}{\ensuremath{\mathbb{Z}}}
\newunicodechar{ℂ}{\ensuremath{\mathbb{C}}}
\newunicodechar{ℕ}{\ensuremath{\mathbb{N}}}
\newunicodechar{ℚ}{\ensuremath{\mathbb{Q}}}
\newunicodechar{𝔻}{\ensuremath{\mathbb{D}}}
\newunicodechar{α}{\ensuremath{\alpha}}
\newunicodechar{β}{\ensuremath{\beta}}
\newunicodechar{γ}{\ensuremath{\gamma}}
\newunicodechar{δ}{\ensuremath{\delta}}
\newunicodechar{ε}{\ensuremath{\varepsilon}}
\newunicodechar{θ}{\ensuremath{\theta}}
\newunicodechar{λ}{\ensuremath{\lambda}}
\newunicodechar{μ}{\ensuremath{\mu}}
\newunicodechar{σ}{\ensuremath{\sigma}}
\newunicodechar{τ}{\ensuremath{\tau}}
\newunicodechar{φ}{\ensuremath{\varphi}}
\newunicodechar{ψ}{\ensuremath{\psi}}
\newunicodechar{ω}{\ensuremath{\omega}}
\newunicodechar{Σ}{\ensuremath{\Sigma}}
% majuscules produites par \MakeUppercase dans les titres (α-aminés -> Α)
\newunicodechar{Α}{\ensuremath{\alpha}}
\newunicodechar{Β}{\ensuremath{\beta}}
\newunicodechar{Γ}{\ensuremath{\Gamma}}
\newunicodechar{Δ}{\ensuremath{\Delta}}
\newunicodechar{Ε}{\ensuremath{\varepsilon}}
\newunicodechar{Θ}{\ensuremath{\Theta}}
\newunicodechar{Λ}{\ensuremath{\Lambda}}
\newunicodechar{Μ}{\ensuremath{\mu}}
\newunicodechar{Τ}{\ensuremath{\tau}}
\newunicodechar{Φ}{\ensuremath{\Phi}}
\newunicodechar{Ψ}{\ensuremath{\Psi}}
\newunicodechar{Ω}{\ensuremath{\Omega}}
\newunicodechar{↦}{\ensuremath{\mapsto}}
\newunicodechar{∈}{\ensuremath{\in}}
\newunicodechar{∉}{\ensuremath{\notin}}
\newunicodechar{∧}{\ensuremath{\wedge}}
\newunicodechar{⇔}{\ensuremath{\Leftrightarrow}}
\newunicodechar{⟺}{\ensuremath{\Longleftrightarrow}}
\newunicodechar{⇒}{\ensuremath{\Rightarrow}}
\newunicodechar{⟹}{\ensuremath{\Longrightarrow}}
\newunicodechar{∘}{\ensuremath{\circ}}
\newunicodechar{∪}{\ensuremath{\cup}}
\newunicodechar{∩}{\ensuremath{\cap}}
\newunicodechar{≡}{\ensuremath{\equiv}}
\newunicodechar{⊂}{\ensuremath{\subset}}
\newunicodechar{⊥}{\ensuremath{\perp}}
\newunicodechar{∃}{\ensuremath{\exists}}
\newunicodechar{∀}{\ensuremath{\forall}}
\newunicodechar{∛}{\ensuremath{\sqrt[3]{\,}}}
\newunicodechar{∜}{\ensuremath{\sqrt[4]{\,}}}
\newunicodechar{⃗}{\ensuremath{{}^{\to}}}
\newunicodechar{ⁿ}{\ifmmode^{n}\else\textsuperscript{\ensuremath{n}}\fi}
\newunicodechar{ˣ}{\ifmmode^{x}\else\textsuperscript{\ensuremath{x}}\fi}
\newunicodechar{ᵇ}{\ifmmode^{b}\else\textsuperscript{\ensuremath{b}}\fi}
\newunicodechar{ʳ}{\ifmmode^{r}\else\textsuperscript{\ensuremath{r}}\fi}
\newunicodechar{ᵘ}{\ifmmode^{u}\else\textsuperscript{\ensuremath{u}}\fi}
\newunicodechar{ʸ}{\ifmmode^{y}\else\textsuperscript{\ensuremath{y}}\fi}
\newunicodechar{ᵃ}{\ifmmode^{a}\else\textsuperscript{\ensuremath{a}}\fi}
\newunicodechar{ᵉ}{\ifmmode^{e}\else\textsuperscript{\ensuremath{e}}\fi}
\newunicodechar{ᵏ}{\ifmmode^{k}\else\textsuperscript{\ensuremath{k}}\fi}
\newunicodechar{ᵖ}{\ifmmode^{p}\else\textsuperscript{\ensuremath{p}}\fi}
\newunicodechar{ᴺ}{\ifmmode^{N}\else\textsuperscript{\ensuremath{N}}\fi}
\newunicodechar{ᶜ}{\ifmmode^{c}\else\textsuperscript{\ensuremath{c}}\fi}
\newunicodechar{ʰ}{\ifmmode^{h}\else\textsuperscript{\ensuremath{h}}\fi}
\newunicodechar{ᵠ}{\ifmmode^{q}\else\textsuperscript{\ensuremath{q}}\fi}
\newunicodechar{ₙ}{\ifmmode_{n}\else\textsubscript{\ensuremath{n}}\fi}
\newunicodechar{ₐ}{\ifmmode_{a}\else\textsubscript{\ensuremath{a}}\fi}
\newunicodechar{ᵢ}{\ifmmode_{i}\else\textsubscript{\ensuremath{i}}\fi}
\newunicodechar{ᵣ}{\ifmmode_{r}\else\textsubscript{\ensuremath{r}}\fi}
\newunicodechar{ₖ}{\ifmmode_{k}\else\textsubscript{\ensuremath{k}}\fi}
\newunicodechar{ᵦ}{\ifmmode_{\beta}\else\textsubscript{\ensuremath{\beta}}\fi}
\newunicodechar{ₓ}{\ifmmode_{x}\else\textsubscript{\ensuremath{x}}\fi}
\newfontfamily\popdisplay{ArchivoBlack-Regular.ttf}[Path=../../../fonts/]
\newfontfamily\poplogo{SpaceGrotesk-Bold.ttf}[Path=../../../fonts/]
\newfontfamily\popsemi{PlusJakartaSans-SemiBold.ttf}[Path=../../../fonts/]
\newfontfamily\popxbold{PlusJakartaSans-ExtraBold.ttf}[Path=../../../fonts/]
\newfontfamily\popxboldls{PlusJakartaSans-ExtraBold.ttf}[Path=../../../fonts/,LetterSpace=3.0]

\setlist[enumerate]{leftmargin=1.7em,itemsep=5pt,topsep=4pt}
\setlist[itemize]{leftmargin=1.7em,itemsep=4pt,topsep=4pt,label={\color{poporange}\textbullet}}
\setlength{\parindent}{0pt}
\setlength{\parskip}{5pt}
\renewcommand{\arraystretch}{1.25}

% pgfplots : style Pop par défaut
\pgfplotsset{
  every axis/.append style={axis line style={popink,line width=1pt},tick style={popink},
    grid style={popline,line width=0.5pt},label style={font=\small},tick label style={font=\footnotesize}},
  popcurve/.style={poporange,line width=1.8pt,no markers,smooth},
}
% Même style disponible dans un \draw TikZ simple (hors pgfplots)
\tikzset{popcurve/.style={poporange,line width=1.8pt}}
\tikzset{popgrid/.style={popline,very thin}}
\colorlet{popcurve}{poporange} % le nom du style est parfois utilisé comme couleur

% ============================================================
% Pill / squiggle
% ============================================================
\newcommand{\poppill}[3]{%
  \tikz[baseline=-0.55ex]\node[draw=popink,line width=1.2pt,rounded corners=2.6pt,%
    fill=#2,inner xsep=6.5pt,inner ysep=3pt]{%
    \popxboldls\fontsize{7}{7}\selectfont\color{#3}\MakeUppercase{#1}};%
}
\newcommand{\popsquiggle}[1]{%
  \tikz[baseline=-0.4ex]{
    \draw[#1,line width=2.6pt,line cap=round]
      plot [smooth] coordinates {(0,0.09) (0.34,0.01) (0.68,0.21) (1.02,0.01) (1.36,0.10)};
  }%
}
% Mot-clé surligné (marqueur orange sous le texte)
\newcommand{\cle}[1]{{\popxbold\color{popdark}#1}}

% ============================================================
% En-tête / pied
% ============================================================
\pagestyle{fancy}
\fancyhf{}
\renewcommand{\headrulewidth}{0pt}
\renewcommand{\footrulewidth}{0pt}
\lhead{\raisebox{-0.15ex}{\poplogo\fontsize{13}{13}\selectfont\color{popink}OnBuch\color{poporange}+}}
\rhead{\poppill{\DOCMATIERE\ \textbullet\ \DOCNIVEAU\ \textbullet\ Leçon \DOCLECON}{white}{popink}}
\lfoot{\popxbold\fontsize{7.6}{7.6}\selectfont\color{popmuted}\MakeUppercase{Cours signé OnBuch+}\ \textbullet\ {\popsemi\DOCTITRE}}
\rfoot{\tikz[baseline=-0.6ex]\node[rounded corners=8.5pt,fill=popink,minimum height=17pt,minimum width=17pt,inner xsep=5pt,inner sep=0]{\popxbold\fontsize{8}{8}\selectfont\color{popcream}\thepage\,/\,\pageref*{LastPage}};}

% ============================================================
% Titres
% ============================================================
\newcommand{\chaptitre}[2]{%
  \begin{tcolorbox}[enhanced,colback=poporange,colframe=popink,boxrule=2.6pt,arc=11pt,
      drop shadow=popink,left=15pt,right=15pt,top=12pt,bottom=11pt,before skip=3pt,after skip=18pt]
    \poppill{#2}{popink}{popcream}\par\vspace{9pt}
    {\raggedright\hyphenpenalty=10000\exhyphenpenalty=10000\popdisplay\fontsize{22}{24}\selectfont\color{popcream}\MakeUppercase{#1}\par}
  \end{tcolorbox}
}
% Section numérotée : pastille orange avec le numéro + titre Archivo
\newcounter{popsec}
\newcommand{\coursec}[1]{%
  \stepcounter{popsec}\setcounter{popsub}{0}%
  \par\vspace{16pt}\noindent
  \begin{minipage}{\linewidth}
  \tikz[baseline=-0.7ex]\node[draw=popink,line width=1.6pt,fill=poporange,rounded corners=4pt,minimum size=22pt,inner sep=0]{\popdisplay\fontsize{12}{12}\selectfont\color{popcream}\thepopsec};%
  \hspace{8pt}{\popdisplay\fontsize{15}{17}\selectfont\color{popink}\MakeUppercase{#1}}\par
  \vspace{4pt}\noindent{\color{poporange}\rule{36pt}{3.6pt}}
  \end{minipage}\par\nopagebreak\vspace{8pt}
}
\newcounter{popsub}
\newcommand{\courssub}[1]{%
  \stepcounter{popsub}%
  \par\vspace{9pt}\noindent{\popxbold\fontsize{11.5}{13}\selectfont\color{popink}{\color{poporange}\thepopsec.\thepopsub}\hspace{6pt}#1}\par\nopagebreak\vspace{4pt}
}

% ============================================================
% Boîtes pédagogiques — même recette : fond blanc, bordure épaisse colorée,
% ombre dure encre, badge plein. Titre optionnel [#1].
% ============================================================
\tcbset{popbase/.style={breakable,enhanced,colback=white,boxrule=2.3pt,arc=9pt,
  shadow={3pt}{-3pt}{0pt}{popink},left=14pt,right=14pt,top=11pt,bottom=11pt,before skip=11pt,after skip=15pt,
  fontupper=\popsemi\fontsize{10.6}{14.5}\selectfont\color{popink},
  fontlower=\popsemi\fontsize{10.6}{14.5}\selectfont\color{popink}}}
% Coche dessinée (le glyphe ✓ n'existe pas dans les polices du document)
\newcommand{\popcheck}[1]{\tikz[baseline=-0.5ex]\draw[#1,line width=1.6pt,line cap=round,line join=round](-0.13,0.0)--(-0.04,-0.09)--(0.13,0.1);}
\providecommand{\coche}{\popcheck{popgreen}}
\providecommand{\resultat}[1]{\par\smallskip\noindent\fcolorbox{popgreen}{popgreenL}{\parbox{\dimexpr\linewidth-2\fboxsep-2\fboxrule\relax}{\textbf{Résultat.} #1}}\par\smallskip}
\newcommand{\popboxhead}[3]{\poppill{#1}{#2}{white}\ifx\relax#3\relax\else\hspace{6pt}{\popxbold\fontsize{10.6}{12}\selectfont\color{popink}#3}\fi\par\vspace{7pt}}

\newtcolorbox{aretenir}[1][]{popbase,colframe=popgreen,before upper={\popboxhead{À retenir}{popgreen}{#1}}}
\newtcolorbox{exemplebox}[1][]{popbase,colframe=poporange,before upper={\popboxhead{Exemple}{poporange}{#1}}}
\newtcolorbox{definition}[1][]{popbase,colframe=popblue,before upper={\popboxhead{Définition}{popblue}{#1}}}
\newtcolorbox{propriete}[1][]{popbase,colframe=poppurple,before upper={\popboxhead{Propriété}{poppurple}{#1}}}
\newtcolorbox{methode}[1][]{popbase,colframe=popgold,before upper={\popboxhead{Méthode}{popgold}{#1}}}
\newtcolorbox{attention}[1][]{popbase,colframe=poppink,before upper={\popboxhead{Attention}{poppink}{#1}}}
\newtcolorbox{experience}[1][]{popbase,colframe=popblue,colback=popblueL,before upper={\popboxhead{Expérience}{popblue}{#1}}}
% « Le savais-tu ? » : carte violette pleine (contexte, culture, Cameroun)
\newtcolorbox{savaistu}[1][]{popbase,colback=poppurple,colframe=popink,
  fontupper=\popsemi\fontsize{10.4}{14.2}\selectfont\color{white},
  before upper={\poppill{Le savais-tu\,?}{white}{poppurple}\ifx\relax#1\relax\else\hspace{6pt}{\popxbold\color{white}#1}\fi\par\vspace{7pt}}}
% Exercice résolu : énoncé en haut, \tcblower puis la solution sur fond crème
\newcounter{popexo}\newcounter{popexr}
\newtcolorbox{exoresolu}[1][]{popbase,colframe=poporange,colbacklower=popcream,
  segmentation style={popink,line width=1.2pt,dashed},
  before upper={\stepcounter{popexr}\popboxhead{Exercice résolu \thepopexr}{poporange}{#1}},
  before lower={\poppill{Solution}{popink}{popcream}\par\vspace{6pt}}}
% Exercice (non résolu en place) — la correction est dans la partie Corrigés
\newtcolorbox{exercice}[1][]{popbase,colframe=popink,
  before upper={\stepcounter{popexo}\popboxhead{Exercice \thepopexo}{popink}{#1}}}
% Corrigé
\newtcolorbox{corrige}[1][]{popbase,colframe=popgreen,colback=popgreenL,
  before upper={\popboxhead{Corrigé}{popgreen}{#1}}}
% Objectifs / prérequis / bilan
\newtcolorbox{objectifs}{popbase,colframe=popink,colback=poporangeL,
  before upper={\popboxhead{Objectifs de la leçon}{popink}{}}}
\newtcolorbox{prerequis}{popbase,colframe=popink,colback=popblueL,
  before upper={\popboxhead{Prérequis}{popblue}{}}}
\newtcolorbox{bilan}{popbase,colframe=popink,colback=white,boxrule=2.8pt,
  before upper={\popboxhead{Fiche bilan}{poporange}{L'essentiel à connaître pour l'examen}}}
% Figure encadrée avec légende
\newtcolorbox{popfigure}[1][]{enhanced,breakable=false,colback=white,colframe=popline,boxrule=1.4pt,arc=8pt,
  left=10pt,right=10pt,top=10pt,bottom=8pt,before skip=10pt,after skip=12pt,halign=center,
  lower separated=false,halign lower=center,
  fontlower=\popsemi\fontsize{9}{11.5}\selectfont\color{popmuted},
  #1}
\newcommand{\legende}[1]{\par\vspace{6pt}{\popsemi\fontsize{9}{11.5}\selectfont\color{popmuted}#1}}

% ============================================================
% Signature OnBuch+
% ============================================================
\newcommand{\poplogoinline}[1]{{\poplogo\fontsize{#1}{#1}\selectfont\color{popink}OnBuch\color{poporange}+}}
\newcommand{\signatureonbuch}{%
  \begin{tcolorbox}[enhanced,colback=popink,colframe=popink,boxrule=2.2pt,arc=10pt,
      drop shadow=poporange,left=14pt,right=14pt,top=10pt,bottom=10pt,before skip=0pt,after skip=0pt]
    \tikz[baseline=-0.7ex]\node[circle,fill=popgreen,draw=popcream,line width=1.3pt,minimum size=22pt,inner sep=0]{\popcheck{white}};%
    \hspace{9pt}\begin{minipage}[c]{0.62\linewidth}
      {\popxbold\fontsize{12}{13}\selectfont\color{popcream}Signé\ \textbullet\ {\poplogo OnBuch\color{poporange}+}}\par\vspace{2pt}
      {\raggedright\popsemi\fontsize{8}{10.5}\selectfont\color{popline}\MakeUppercase{Équipe pédagogique OnBuch+}\\\MakeUppercase{Conforme au programme officiel MINESEC}\par}
    \end{minipage}\hfill
    {\poplogo\fontsize{22}{22}\selectfont\color{popcream}OnBuch\color{poporange}+}
  \end{tcolorbox}%
}

% ============================================================
% Page de garde
% #1 titre · #2 matière · #3 niveau · #4 module · #5 n° leçon · #6 durée ·
% #7 description / objectifs (texte ou itemize)
% ============================================================
\newcommand{\pagegarde}[7]{%
  \thispagestyle{empty}
  \newgeometry{a4paper,margin=1.7cm,top=1.6cm,bottom=1.4cm}%
  \begin{tikzpicture}[remember picture,overlay]
    \fill[poporange!18] ([xshift=-3.2cm,yshift=-3.6cm]current page.north east) circle (4.6cm);
    \fill[poppurple!14] ([xshift=2.4cm,yshift=7.6cm]current page.south west) circle (3.8cm);
  \end{tikzpicture}%
  {\poplogo\fontsize{30}{30}\selectfont\color{popink}OnBuch\color{poporange}+}\hfill
  \raisebox{0.6ex}{\poppill{Cours signé \textbullet\ Premium}{popink}{popcream}}\par
  \vspace{1pt}\popsquiggle{poppurple}\par
  \vspace{8pt}
  \begin{tcolorbox}[enhanced,colback=poporange,colframe=popink,boxrule=2.8pt,arc=13pt,
      drop shadow=popink,left=18pt,right=18pt,top=12pt,bottom=13pt,after skip=10pt]
    \poppill{#2 \textbullet\ #3}{popink}{popcream}\hspace{5pt}\poppill{#4}{white}{popink}\par\vspace{8pt}
    {\popdisplay\fontsize{14}{14}\selectfont\color{popink}LEÇON #5}\par\vspace{4pt}
    {\raggedright\hyphenpenalty=10000\exhyphenpenalty=10000\popdisplay\fontsize{25}{27}\selectfont\color{popcream}\MakeUppercase{#1}\par}
    \vspace{8pt}
    {\popxbold\fontsize{9.5}{9.5}\selectfont\color{popink}\MakeUppercase{Durée conseillée : #6}}
  \end{tcolorbox}
  \begin{tcolorbox}[enhanced,colback=white,colframe=popink,boxrule=2.2pt,arc=10pt,
      drop shadow=popink,left=16pt,right=16pt,top=10pt,bottom=10pt,after skip=8pt]
    \poppill{Dans cette leçon}{poppurple}{white}\par\vspace{5pt}
    \popsemi\fontsize{9.3}{12.6}\selectfont\color{popink}#7
  \end{tcolorbox}
  \noindent\begin{minipage}[t]{0.487\linewidth}
    \begin{tcolorbox}[enhanced,colback=popgreen,colframe=popink,boxrule=2.2pt,arc=10pt,
        drop shadow=popink,left=11pt,right=11pt,top=10pt,bottom=10pt,height=3.1cm,valign=center]
      \tcbox[colback=white,colframe=popink,boxrule=1.3pt,arc=5pt,boxsep=0pt,nobeforeafter,
        left=3pt,right=3pt,top=3pt,bottom=3pt,valign=center]{\qrcode[height=1.9cm]{https://play.google.com/store/apps/details?id=com.luvvix.onbuch&hl=fr}}%
      \hspace{8pt}\begin{minipage}[c]{0.5\linewidth}
        {\popdisplay\fontsize{11}{12.5}\selectfont\color{white}\MakeUppercase{Télécharge OnBuch}}\par\vspace{4pt}
        {\raggedright\popsemi\fontsize{8.4}{11}\selectfont\color{white}Gratuit \textbullet\ Google Play\\Tuteur IA, annales, concours, résultats\par}
      \end{minipage}
    \end{tcolorbox}
  \end{minipage}\hfill
  \begin{minipage}[t]{0.487\linewidth}
    \begin{tcolorbox}[enhanced,colback=poppurple,colframe=popink,boxrule=2.2pt,arc=10pt,
        drop shadow=popink,left=11pt,right=11pt,top=10pt,bottom=10pt,height=3.1cm,valign=center]
      \tikz[baseline=-0.7ex]\node[circle,fill=white,draw=popink,line width=1.3pt,minimum size=1.6cm,inner sep=0]{\popdisplay\fontsize{24}{24}\selectfont\color{poppurple}+};%
      \hspace{8pt}\begin{minipage}[c]{0.6\linewidth}
        {\popdisplay\fontsize{11}{12.5}\selectfont\color{white}\MakeUppercase{Passe à OnBuch+}}\par\vspace{4pt}
        {\raggedright\popsemi\fontsize{8.4}{11}\selectfont\color{white}Tous les cours signés, Tuteur IA illimité, fiches PDF — onglet OnBuch+ de l'app\par}
      \end{minipage}
    \end{tcolorbox}
  \end{minipage}
  \vspace{8pt}
  \popcontenu
  \vfill
  \signatureonbuch
  \restoregeometry
  \newpage
}

% Rangée « Ce cours contient » (page de garde)
\newcommand{\poptile}[3]{% #1 couleur #2 gros texte #3 légende
  \begin{tcolorbox}[enhanced,colback=#1,colframe=popink,boxrule=2pt,arc=9pt,shadow={2.5pt}{-2.5pt}{0pt}{popink},
      left=6pt,right=6pt,top=9pt,bottom=9pt,halign=center,valign=center,height=1.75cm,width=\linewidth,nobeforeafter]
    {\popdisplay\fontsize{12.5}{13}\selectfont\color{white}\MakeUppercase{#2}}\par\vspace{3pt}
    {\centering\popsemi\fontsize{7.8}{9.5}\selectfont\color{white}#3\par}
  \end{tcolorbox}}
\newcommand{\popcontenu}{%
  \par\noindent{\popxboldls\fontsize{7.5}{7.5}\selectfont\color{popmuted}CE COURS SIGNÉ CONTIENT}\par\vspace{7pt}
  \noindent\begin{minipage}[t]{0.235\linewidth}\poptile{popblue}{Cours}{complet, illustré, pas à pas}\end{minipage}\hfill
  \begin{minipage}[t]{0.235\linewidth}\poptile{poporange}{Méthodes}{et exercices résolus}\end{minipage}\hfill
  \begin{minipage}[t]{0.235\linewidth}\poptile{poppink}{Exercices}{type examen + corrigés}\end{minipage}\hfill
  \begin{minipage}[t]{0.235\linewidth}\poptile{popgreen}{Fiche bilan}{l'essentiel à retenir}\end{minipage}\par}

% Sommaire
\newtcolorbox{sommaire}{enhanced,colback=white,colframe=popink,boxrule=2.2pt,arc=10pt,
  drop shadow=popink,left=18pt,right=18pt,top=13pt,bottom=13pt,before skip=6pt,after skip=20pt,
  before upper={\poppill{Sommaire}{poppurple}{white}\par\vspace{9pt}\popsemi\fontsize{10.6}{14.5}\selectfont\color{popink}}}

% Signature de fin de cours — poussée tout en bas de la dernière page
\newcommand{\findecours}{%
  \par\vfill\nopagebreak
  \begin{center}\popsquiggle{poporange}\end{center}\vspace{4pt}
  \signatureonbuch
}
\AtBeginDocument{\def\llbracket{\mathopen{[\mkern-3.5mu[}}\def\rrbracket{\mathclose{]\mkern-3.5mu]}}} % intervalles d'entiers (glyphe ⟦ absent de la police)
% Glyphes absents de Fira Math : redéfinis à partir de glyphes disponibles
\AtBeginDocument{%
  \def\setminus{\mathbin{\backslash}}\let\smallsetminus\setminus
  \def\vdots{\mathord{\vbox{\baselineskip3.2pt\lineskiplimit0pt\kern4pt\hbox{.}\hbox{.}\hbox{.}}}}%
  \def\ddots{\mathinner{\mkern1mu\raise6.5pt\hbox{.}\mkern2mu\raise3.6pt\hbox{.}\mkern2mu\raise0.7pt\hbox{.}\mkern1mu}}%
  \def\triangle{\mathord{\scalebox{0.9}{$\symup{\Delta}$}}}%
  \def\nabla{\mathord{\raisebox{\depth}{\scalebox{1}[-1]{$\symup{\Delta}$}}}}%
  \def\checkmark{\popcheck{popdark}}%
  \def\star{\mathbin{\ast}}\def\bigstar{\mathord{\ast}}%
  \def\diamond{\mathbin{\scalebox{0.6}{$\Diamond$}}}%
  \def\bigcup{\mathop{\mathchoice{\raisebox{-0.3em}{\scalebox{1.7}{$\cup$}}}{\scalebox{1.25}{$\cup$}}{\cup}{\cup}}\displaylimits}%
  \def\bigcap{\mathop{\mathchoice{\raisebox{-0.3em}{\scalebox{1.7}{$\cap$}}}{\scalebox{1.25}{$\cap$}}{\cap}{\cap}}\displaylimits}%
  \def\lesseqqgtr{\mathrel{\lessgtr}}\def\bigoplus{\mathop{\oplus}\displaylimits}%
}
\providecommand{\ssi}{si et seulement si} % abréviation fréquente
\AtBeginDocument{\providecommand{\wideparen}{\overparen}} % arc de cercle

%% FIGFIX-BEGIN v3 — correctifs globaux des figures (docs/onbuch-generation/tools/figfix)
\usepackage{adjustbox}\usepackage{etoolbox}
% toute figure TikZ/pgfplots plus large que la ligne est réduite à la largeur disponible
\BeforeBeginEnvironment{tikzpicture}{\begin{adjustbox}{max width=\linewidth}}
\AfterEndEnvironment{tikzpicture}{\end{adjustbox}}
% légendes pgfplots lisibles par-dessus les courbes
\pgfplotsset{every axis/.append style={legend style={fill=white,fill opacity=0.92,text opacity=1,draw=black!15,inner sep=3pt}}}
% étiquettes d'arêtes (syntaxe edge["..."]) avec fond blanc
\tikzset{every edge quotes/.append style={fill=white,inner sep=1.2pt}}
% bulles de cartes mentales : texte à la ligne dans la bulle, bulle un peu plus grande
\tikzset{every concept/.append style={text width=2.0cm,align=center,minimum size=2.3cm,inner sep=2pt}}
%% FIGFIX-END
\begin{document}

\pagegarde{La liberté et la responsabilité}{Philosophie}{Tle Tech}{Philosophie – classe de Terminale technique}{N08}{2 séances (fiche de progression harmonisée)}{Cette leçon interroge deux notions centrales de la philosophie morale et politique : la liberté et la responsabilité. À partir de situations concrètes camerounaises et africaines, elle montre que la liberté, loin d'être un pouvoir absolu de tout faire, est ce qui fonde notre responsabilité devant soi, devant autrui et devant la loi.
\vspace{4pt}
\begin{multicols}{2}\small
\begin{itemize}
\item À la fin de cette leçon, je sais définir les concepts de liberté et de responsabilité avec précision
\item Distinguer la liberté physique, la liberté morale et la liberté politique
\item Expliquer en quoi la liberté est le fondement de la responsabilité
\item Identifier les différents types de responsabilité (morale, juridique, sociale)
\item Analyser une situation concrète de la vie courante à l'aide des notions de liberté et de responsabilité
\item Construire un schéma conceptuel (diagramme) reliant liberté, choix, acte et responsabilité
\end{itemize}\end{multicols}}

\begin{sommaire}
\begin{multicols}{2}
\begin{enumerate}
\item Introduction : problématisation et définitions
\item Les formes de la liberté
\item Les formes de la responsabilité
\item La liberté, fondement de la responsabilité
\item Applications concrètes au Cameroun
\item Méthode : rédiger et justifier une démarche philosophique
\item Activité d'intégration
\item Méthodes et exercices résolus
\item Exercices
\item Corrigés des exercices
\item Fiche bilan
\end{enumerate}
\end{multicols}
\end{sommaire}

\begin{objectifs}
\begin{itemize}
\item À la fin de cette leçon, je sais définir les concepts de liberté et de responsabilité avec précision
\item Distinguer la liberté physique, la liberté morale et la liberté politique
\item Expliquer en quoi la liberté est le fondement de la responsabilité
\item Identifier les différents types de responsabilité (morale, juridique, sociale)
\item Analyser une situation concrète de la vie courante à l'aide des notions de liberté et de responsabilité
\item Construire un schéma conceptuel (diagramme) reliant liberté, choix, acte et responsabilité
\item Rédiger et justifier une démarche argumentée, une réponse ou une conclusion en dissertation
\item Mobiliser des références philosophiques (Sartre, Kant, Spinoza) à bon escient dans un devoir
\end{itemize}
\end{objectifs}

\begin{prerequis}
\begin{itemize}
\item La notion de devoir et d'obligation (éducation civique, collège)
\item La distinction entre le permis et le défendu (vie en société)
\item La notion de conscience morale (leçons antérieures de Terminale)
\item Les notions de droit et de justice (programme de Première)
\item La méthode de la dissertation philosophique (introduction, problématique)
\end{itemize}
\end{prerequis}

\newpage

\coursec{Introduction : problématisation et définitions}

\courssub{Situation d'accroche : le moto-taxi de Yaoundé}

Il est 17 heures, avenue Kennedy, à Yaoundé. La circulation est dense, les klaxons résonnent, les feux tricolores rythment péniblement la sortie des bureaux. Un jeune chauffeur de moto-taxi, que nous appellerons Arnaud, transporte un client pressé de rejoindre la gare d'Obili avant la fermeture. Devant un feu rouge, Arnaud hésite une fraction de seconde : le client le presse, lui promet un pourboire, lui crie de passer. Arnaud accélère, brûle le feu... et percute un piéton qui traversait sur le passage protégé. Le piéton est blessé, la moto endommagée, la foule se rassemble.

Devant le commissaire de police, Arnaud se défend avec véhémence : \og Ce n'est pas ma faute, c'est le client qui me pressait ! \fg{}

Cette scène, banale dans nos villes camerounaises, soulève immédiatement une série de questions philosophiques profondes :

\begin{itemize}
\item A-t-on le \cle{droit} de rejeter ainsi la faute sur autrui ? La pression du client efface-t-elle la décision d'Arnaud ?
\item Arnaud était-il \cle{libre} au moment où il a accéléré ? Pouvait-il faire autrement ?
\item Jusqu'où doit-il \cle{répondre} de son acte : devant la loi, devant le blessé, devant sa propre conscience ?
\end{itemize}

\begin{attention}[Erreur fréquente dès l'analyse de la situation]
Beaucoup d'élèves répondent trop vite : \og c'est la faute du client, il a payé pour qu'on roule vite \fg{}, ou au contraire \og c'est entièrement la faute du chauffeur \fg{}. La philosophie ne tranche pas d'emblée : elle \emph{analyse}. Il faut distinguer la \textbf{cause} (la pression du client est une circonstance) de la \textbf{décision} (c'est Arnaud qui a tourné la poignée d'accélérateur). Confondre circonstance et décision, c'est déjà confondre contrainte et liberté.
\end{attention}

Cette situation quotidienne pose le problème central de cette leçon : \textbf{la liberté est-elle le fondement de la responsabilité ?} Autrement dit : faut-il être libre pour être tenu responsable de ses actes, et sommes-nous réellement libres ?

\courssub{Étymologie des deux concepts}

Avant de définir rigoureusement, interrogeons les mots eux-mêmes. L'étymologie n'est pas un ornement : elle révèle l'intuition première que les hommes ont eue de ces notions.

\textbf{Liberté.} Le mot vient du latin \emph{libertas}, qui désigne l'\textbf{état de celui qui n'est pas esclave}. Dans la Rome antique, le \emph{liber} (l'homme libre) s'oppose à l'\emph{servus} (l'esclave) : l'esclave appartient à un maître qui décide à sa place, tandis que l'homme libre dispose de lui-même, de ses actes et de ses biens. La liberté est donc, à l'origine, un \emph{statut} : celui de n'être la chose de personne. Cette intuition demeure : être libre, c'est d'abord ne pas dépendre de la volonté d'un autre.

\textbf{Responsabilité.} Le mot vient du latin \emph{respondere}, qui signifie \og répondre en retour \fg{}, puis \og répondre de quelque chose \fg{}, c'est-à-dire \textbf{en rendre compte}. Dans le droit romain, \emph{respondere} désigne l'engagement de celui qui garantit une dette ou qui comparaît pour justifier ses actes. Être responsable, c'est donc étymologiquement pouvoir être \emph{appelé à répondre} : devant un juge, devant autrui, devant sa conscience.

\begin{savaistu}[D'où vient le mot « responsabilité » ?]
Le mot \og responsabilité \fg{} vient du latin \emph{respondere}, \og répondre en retour \fg{}. Être responsable, c'est littéralement \textbf{pouvoir répondre de soi}. Remarquons que dans beaucoup de nos langues camerounaises, l'idée est semblable : celui qui a causé un dommage doit \og venir s'expliquer \fg{} devant la famille ou le chef traditionnel. La palabre villageoise est une forme ancienne de mise en responsabilité : on ne fuit pas, on répond.
\end{savaistu}

\courssub{Définition de la liberté}

Partons de l'intuition. Quand je dis \og je suis libre \fg{}, je veux dire : personne ne me force, je choisis moi-même. L'enfant qui sort de classe sans autorisation se sent libre ; le prisonnier derrière les barreaux ne l'est pas. L'intuition commune identifie donc la liberté à l'\textbf{absence d'obstacle extérieur} et au \textbf{pouvoir de choisir}.

Mais cette intuition doit être précisée. Arnaud, notre moto-taximan, n'était enchaîné à rien : il pouvait accélérer ou s'arrêter. La pression du client était une incitation, non une chaîne. Il disposait donc du pouvoir de se déterminer lui-même.

\begin{definition}[La liberté]
La \cle{liberté} est le \textbf{pouvoir de se déterminer soi-même à agir ou à ne pas agir}, sans contrainte extérieure, selon sa propre volonté. Elle suppose deux conditions : l'\emph{absence de contrainte} (personne ne me force physiquement) et la \emph{capacité de choisir} (je peux délibérer entre plusieurs possibles).
\end{definition}

Analysons la structure de cette définition. Dire \og se déterminer soi-même \fg{}, c'est affirmer que la \textbf{cause de mon acte est en moi} et non hors de moi. Le philosophe Aristote distinguait déjà l'acte volontaire (dont le principe est dans l'agent) de l'acte involontaire (accompli par force ou par ignorance). Spinoza ajoutera qu'est libre ce qui \og existe par la seule nécessité de sa nature et se détermine seul à agir \fg{}. La liberté n'est donc pas n'importe quel mouvement : c'est un mouvement dont \emph{je} suis l'auteur.

\begin{attention}[La liberté n'est pas le caprice]
Erreur fréquente : confondre la liberté avec l'\textbf{absence totale de règles}, c'est-à-dire avec le \emph{caprice} ou la licence. Faire \og tout ce qui me passe par la tête \fg{} (insulter, voler, brûler les feux) n'est pas être libre : c'est obéir à chaque pulsion, donc être esclave de ses désirs immédiats. La liberté véritable suppose au contraire une \textbf{volonté réfléchie} : je choisis en connaissance de cause, et je peux même choisir de respecter une règle. Le feu rouge ne supprime pas la liberté d'Arnaud ; il lui donne au contraire l'occasion d'exercer un choix responsable.
\end{attention}

\courssub{Définition de la responsabilité}

Revenons au commissariat. Le policier demande à Arnaud : \og Qui conduisait la moto ? \fg{} Cette question simple institue Arnaud comme \emph{responsable} : on l'identifie comme l'auteur de l'acte, et on attend de lui qu'il en réponde.

\begin{definition}[La responsabilité]
La \cle{responsabilité} est le \textbf{fait de devoir répondre de ses actes} devant soi-même, devant autrui ou devant la loi, et d'\textbf{en assumer les conséquences} (réparation, sanction, remords).
\end{definition}

La responsabilité comporte trois dimensions qu'il faut distinguer dès maintenant :

\begin{enumerate}
\item \textbf{Devant la loi} (responsabilité juridique) : Arnaud devra répondre devant le tribunal, payer une amende, réparer le dommage, peut-être purger une peine.
\item \textbf{Devant autrui} (responsabilité morale envers les autres) : il doit quelque chose au piéton blessé — au minimum la reconnaissance de sa faute.
\item \textbf{Devant soi-même} (responsabilité de conscience) : même si personne ne l'avait vu, Arnaud saurait ce qu'il a fait ; le remords est la voix de cette responsabilité intérieure.
\end{enumerate}

\begin{exemplebox}[Répondre de ses actes dans la vie courante]
Un élève de Terminale technique triche à une évaluation et est découvert. Dire qu'il est \emph{responsable}, c'est dire trois choses : on peut lui \emph{imputer} l'acte (c'est bien lui qui a recopié), on peut lui en \emph{demander compte} (conseil de discipline), et il doit en \emph{assumer les conséquences} (zéro, sanction). S'il dit \og ce n'est pas ma faute, mon voisin m'a montré sa copie \fg{}, il fait exactement le raisonnement d'Arnaud : il dissout sa décision dans la circonstance.
\end{exemplebox}

\courssub{Distinction préliminaire : faire ce qu'on veut et assumer ce qu'on fait}

Nous pouvons maintenant formuler la distinction qui organise toute la leçon :

\begin{itemize}
\item La \textbf{liberté} pose la question : \emph{puis-je faire ce que je veux ?} Elle regarde le moment \emph{avant} et \emph{pendant} l'acte : le choix, la délibération, l'absence de contrainte.
\item La \textbf{responsabilité} pose la question : \emph{dois-je assumer ce que j'ai fait ?} Elle regarde le moment \emph{après} l'acte : l'imputation, le compte à rendre, la sanction ou la réparation.
\end{itemize}

Ces deux notions semblent distinctes, voire opposées : la liberté semble un droit (\og je fais ce que je veux \fg{}), la responsabilité une charge (\og tu paieras ce que tu as fait \fg{}). Pourtant, notre situation d'accroche suggère qu'elles sont liées : on ne demande de comptes à Arnaud que parce qu'on le suppose libre d'avoir agi autrement. Si sa moto avait été projetée par une explosion, personne ne l'accuserait.

\begin{popfigure}
\begin{center}
\begin{tikzpicture}
\draw[fill=popblueL,draw=popblue,thick] (0,0) circle (2.6);
\draw[fill=poporangeL,draw=poporange,thick] (0,0) circle (1.5);
\node[popdark,font=\bfseries] at (0,0.35) {Responsabilité};
\node[popmuted,font=\small] at (0,-0.35) {(conséquence)};
\node[popblue,font=\bfseries] at (0,-2.05) {Liberté (fondement)};
\draw[-{Stealth},thick,popdark] (0,-1.5) -- (0,-1.15);
\node[align=center,font=\small,popmuted] at (4.6,0.6) {Parce que je suis libre,\\ on me tient pour\\ responsable.};
\draw[-{Stealth},popmuted] (2.6,0.2) -- (3.4,0.5);
\end{tikzpicture}
\end{center}
\legende{Schéma 1 — La responsabilité s'enracine dans la liberté : la liberté est le cercle fondamental, la responsabilité en est la conséquence. Sans liberté supposée, aucune imputation n'est possible.}
\end{popfigure}

\begin{popfigure}
\begin{center}
\begin{tikzpicture}[mindmap, grow cyclic,
  concept color=popblueL,
  level 1 concept/.append style={sibling angle=72, minimum size=2.1cm, font=\small\bfseries},
  root concept/.append style={minimum size=2.6cm, font=\bfseries}]
\node[concept] {Liberté et responsabilité}
  child[concept color=poporangeL] { node {choix} }
  child[concept color=popgreenL] { node {volonté} }
  child[concept color=poppurpleL] { node {contrainte} }
  child[concept color=poppinkL] { node {sanction} }
  child[concept color=popgoldL] { node {culpabilité} };
\end{tikzpicture}
\end{center}
\legende{Carte mentale — Les cinq mots-clés de la leçon. Le \emph{choix} et la \emph{volonté} relèvent de la liberté ; la \emph{sanction} et la \emph{culpabilité} relèvent de la responsabilité ; la \emph{contrainte} est le concept-frontière qui menace la liberté et peut atténuer la responsabilité.}
\end{popfigure}

\courssub{Problématique et annonce du plan}

Rassemblons nos acquis. Nous savons désormais ce que signifient la liberté (se déterminer soi-même) et la responsabilité (répondre de ses actes). Nous avons pressenti qu'elles sont liées : la responsabilité semble \emph{présupposer} la liberté. Mais ce lien est-il nécessaire ? Ne sommes-nous pas, comme Arnaud le prétend, déterminés par mille causes (la pression sociale, la misère, l'éducation, le tempérament) ? Et si nous ne sommes pas libres, peut-on encore nous tenir responsables ?

\textbf{Problématique de la leçon :} \emph{la liberté est-elle la condition qui rend l'homme responsable de ses actes ?}

Pour répondre à cette question, nous suivrons le plan suivant :

\begin{enumerate}
\item \textbf{Les concepts} : les formes de la liberté (faire ce qu'on veut ? agir selon sa volonté ? obéir à la raison ?) et les formes de la responsabilité (juridique, morale, sociale).
\item \textbf{La liberté comme fondement de la responsabilité} : montrer que seul un être libre peut être tenu responsable, et examiner les objections (le déterminisme, les circonstances atténuantes).
\item \textbf{Applications} : appliquer ces notions à des situations concrètes de la vie camerounaise et apprendre à rédiger et justifier une démarche philosophique.
\end{enumerate}

\begin{exoresolu}[Première application : qui est responsable ?]
\textbf{Énoncé.} Reprenons la situation d'accroche. Le client d'Arnaud affirme au commissaire : \og Moi je n'ai rien fait, je ne conduisais pas ! \fg{} Arnaud affirme : \og Sans sa pression, je me serais arrêté. \fg{} En vous appuyant sur les définitions de la liberté et de la responsabilité, déterminez qui doit répondre de l'accident, et pourquoi.
\tcblower
\textbf{Solution détaillée.}
Étape 1 — Identifier l'auteur de l'acte. C'est Arnaud qui tenait le guidon et a tourné la poignée d'accélérateur. L'acte matériel lui est \emph{imputable}.

Étape 2 — Vérifier la condition de liberté. Arnaud n'était ni attaché, ni menacé d'une arme, ni privé de conscience. La pression du client est une \emph{incitation} (une circonstance), non une \emph{contrainte} au sens strict : il conservait le pouvoir de se déterminer à ne pas agir (ralentir, refuser, déposer le client). Il était donc libre au sens de notre définition.

Étape 3 — Conclure sur la responsabilité. Étant libre et auteur de l'acte, Arnaud est \textbf{responsable} : il doit répondre devant la loi (infraction au code de la route, dommages au piéton), devant autrui (réparation) et devant sa conscience.

Étape 4 — Nuancer. Le client n'est pas l'auteur de l'acte, mais sa pression peut être examinée comme \emph{circonstance} : elle n'efface pas la responsabilité d'Arnaud, elle peut tout au plus éclairer le contexte. La philosophie distingue ainsi \emph{expliquer} (comprendre les causes) et \emph{excuser} (effacer la faute) : expliquer n'est pas excuser.
\end{exoresolu}

\begin{aretenir}[L'essentiel de l'introduction]
\begin{itemize}
\item \textbf{Liberté} (latin \emph{libertas}, état de celui qui n'est pas esclave) : pouvoir de se déterminer soi-même à agir ou à ne pas agir, selon sa propre volonté. Elle n'est ni le caprice ni l'absence de règles.
\item \textbf{Responsabilité} (latin \emph{respondere}, répondre en retour) : fait de devoir répondre de ses actes devant soi-même, autrui ou la loi, et d'en assumer les conséquences.
\item Les deux notions sont liées : on ne tient responsable que celui qu'on suppose \textbf{libre}. La liberté apparaît ainsi comme le \textbf{fondement} de la responsabilité — c'est la thèse que la leçon devra éprouver.
\item Problématique : la liberté est-elle la condition qui rend l'homme responsable de ses actes ?
\end{itemize}
\end{aretenir}

\coursec{Les formes de la liberté}

Nous avons vu dans l'introduction que la liberté se définit d'abord comme le pouvoir d'agir sans contrainte et selon sa volonté, et que cette définition appelle aussitôt une question : la liberté est-elle une réalité unique, ou se décline-t-elle en plusieurs formes ? L'observation de la vie quotidienne suffit à montrer qu'on ne parle pas de la même chose lorsqu'on dit qu'un détenu de la prison centrale de Yaoundé n'est « pas libre », qu'un citoyen est « libre » de voter, ou qu'un élève est « libre » de choisir entre travailler et tricher. Il faut donc distinguer les \cle{formes de la liberté} : la liberté physique, la liberté morale, la liberté politique, et enfin la liberté la plus profonde, le libre arbitre. Chacune répond à un type d'obstacle différent.

\courssub{La liberté physique : l'absence d'obstacle extérieur}

Commençons par l'intuition la plus simple. Un oiseau en cage ne peut pas voler ; un homme enchaîné ne peut pas marcher. Dans les deux cas, il existe un \cle{obstacle extérieur} — les barreaux, la chaîne — qui empêche le corps d'accomplir le mouvement voulu. La liberté physique est précisément l'absence d'un tel obstacle.

\begin{definition}[Liberté physique]
La \cle{liberté physique} est le pouvoir de mouvoir son corps et de se déplacer sans rencontrer d'obstacle matériel extérieur. Elle est réalisée lorsque rien, dans l'espace, ne s'oppose à l'accomplissement du mouvement voulu.
\end{definition}

L'exemple paradigmatique est celui du \cle{prisonnier} : il n'est pas libre de ses mouvements, non parce que sa volonté serait affaiblie, mais parce que les murs de sa cellule s'opposent physiquement à son déplacement. De même, un voyageur bloqué par un pont effondré sur la route Douala--Bafoussam n'est pas libre de continuer son trajet : l'obstacle est matériel, indépendant de sa volonté.

Mais une observation attentive révèle une limite de cette conception. Le prisonnier, même enchaîné, reste capable de \emph{vouloir} s'évader, de choisir intérieurement son attitude face à sa captivité. L'obstacle physique n'atteint donc que le corps, jamais la pensée. C'est ce que les philosophes stoïciens avaient déjà remarqué : Épictète, lui-même ancien esclave, soutenait que l'on peut enchaîner le corps mais non la volonté. La liberté physique est ainsi la forme la plus \emph{extérieure} et la plus fragile de la liberté : elle dépend des circonstances matérielles.

\begin{attention}[Piège à éviter]
Ne pas confondre la liberté physique avec la liberté tout court. Dire « je suis libre » parce qu'aucune chaîne ne me retient ne prouve pas que je suis libre de choisir : un homme libre de ses mouvements peut être esclave de ses passions, de la mode ou de l'opinion des autres. L'absence d'obstacle extérieur est une condition nécessaire, mais non suffisante, de la liberté complète.
\end{attention}

\courssub{La liberté morale : choisir entre le bien et le mal}

Passons à un niveau plus intérieur. Un élève qui trouve le sujet d'examen sur la table d'un camarade absent peut recopier ou s'abstenir. Aucun obstacle physique ne l'empêche de tricher ; pourtant, il éprouve que le choix lui appartient et qu'il engage sa valeur morale. C'est le domaine de la \cle{liberté morale}.

\begin{definition}[Liberté morale]
La \cle{liberté morale} est la capacité de choisir entre le bien et le mal en suivant le jugement de sa propre conscience, et non sous la pression d'une contrainte extérieure ou d'une passion aveugle.
\end{definition}

La \cle{conscience morale} joue ici le rôle de juge intérieur : elle me fait connaître ce qui est bien (dire la vérité, respecter autrui) et ce qui est mal (mentir, voler), et elle me laisse la responsabilité du choix. La liberté morale suppose donc deux choses : d'une part la connaissance du bien et du mal, d'autre part le pouvoir effectif d'opter pour l'un ou pour l'autre. Un enfant très jeune ou une personne privée de discernement ne possède pas pleinement cette liberté, car le discernement en est la condition.

Remarquons que la liberté morale est plus exigeante que la liberté physique : elle implique un \emph{effort}. Céder à la colère, à la peur ou à l'appât du gain est facile ; résister demande que la volonté commande aux désirs. C'est pourquoi les moralistes anciens comparaient l'âme à un cocher conduisant un attelage de chevaux fougueux : être moralement libre, c'est tenir les rênes de ses passions.

\courssub{La liberté politique : les droits du citoyen}

Ni la liberté physique ni la liberté morale ne suffisent à garantir une vie libre en société. Un homme peut n'avoir ni chaîne ni remords, et pourtant être empêché de s'exprimer, de circuler ou de participer à la vie de la cité par un pouvoir arbitraire. D'où la troisième forme : la \cle{liberté politique}.

\begin{definition}[Liberté politique]
La \cle{liberté politique} est l'ensemble des droits civiques reconnus et garantis à chaque citoyen par l'État de droit : liberté d'expression, liberté de circulation, liberté de réunion, droit de vote, liberté de la presse.
\end{definition}

Ces libertés ont un caractère particulier : elles ne sont pas seulement des pouvoirs de fait, mais des \cle{droits}, c'est-à-dire des pouvoirs que la loi protège. Au Cameroun, la Constitution garantit la liberté d'expression, de la presse, de circulation et le droit de participer à la vie politique par le vote. Le citoyen qui se rend au bureau de vote le jour d'une élection exerce sa liberté politique ; le journaliste qui publie une enquête aussi.

Montesquieu, dans \emph{De l'esprit des lois} (1748), a donné de cette liberté une définition célèbre : « La liberté est le droit de faire tout ce que les lois permettent. » La liberté politique n'est donc pas la licence — faire tout ce qu'on veut — mais la possibilité garantie d'agir dans le cadre de la loi commune. Elle suppose l'\cle{État de droit}, c'est-à-dire un État où le pouvoir lui-même est soumis à la loi et où les droits des citoyens sont protégés contre l'arbitraire.

\begin{savaistu}[La source des libertés modernes]
L'article 10 de la Déclaration des droits de l'homme et du citoyen de 1789 affirme : « Nul ne doit être inquiété pour ses opinions, même religieuses, pourvu que leur manifestation ne trouble pas l'ordre public établi par la loi. » Ce texte, né de la Révolution française, est la source des constitutions modernes, y compris africaines et camerounaises, qui reconnaissent la liberté d'opinion et d'expression comme un droit fondamental.
\end{savaistu}

\begin{popfigure}
\begin{center}
\begin{tabularx}{\linewidth}{|l|X|X|}
\hline
\rowcolor{popblueL} \textbf{Forme de liberté} & \textbf{Définition} & \textbf{Exemple camerounais} \\
\hline
\textbf{Physique} & Absence d'obstacle matériel extérieur au mouvement du corps. & Le détenu de la prison de New-Bell à Douala n'est pas libre de ses mouvements ; le voyageur libre circule sur l'axe Yaoundé--Bamenda. \\
\hline
\textbf{Morale} & Capacité de choisir entre le bien et le mal selon sa conscience. & L'élève de Terminale à Garoua qui refuse de recopier le devoir de son voisin, malgré la tentation. \\
\hline
\textbf{Politique} & Droits civiques (vote, expression, circulation, presse) garantis par l'État de droit. & L'électeur qui vote à Yaoundé ; le journaliste qui publie un article dans le respect de la loi. \\
\hline
\end{tabularx}
\end{center}
\legende{Les trois formes fondamentales de la liberté : à chaque forme correspond un type d'obstacle différent (matériel, passionnel, arbitraire politique).}
\end{popfigure}

\courssub{Le libre arbitre : la liberté comme pouvoir de choisir}

Reste la question la plus profonde : lorsque je choisis, mon choix est-il vraiment \emph{à moi}, ou est-il le produit de causes qui m'échappent ? C'est le problème du \cle{libre arbitre}.

\begin{definition}[Libre arbitre]
Le \cle{libre arbitre} est la faculté de se déterminer par sa seule volonté, indépendamment de toute contrainte ou cause extérieure. Celui qui agit par libre arbitre est la source première de son acte : il pouvait choisir autrement.
\end{definition}

L'expérience intérieure semble attester le libre arbitre : lorsque je délibère entre deux conduites — dire ou taire la vérité —, j'éprouve que la décision finale dépend de moi et que je pourrais choisir l'inverse. Cette expérience de l'hésitation et de la décision est le principal argument en faveur du libre arbitre : nul ne délibère sur ce qu'il croit nécessaire ou impossible.

Mais cette expérience est-elle une preuve suffisante ? C'est ici qu'intervient la thèse adverse.

\begin{definition}[Déterminisme]
Le \cle{déterminisme} est la doctrine selon laquelle tous nos actes sont causés par des facteurs déterminants : la nature (hérédité, tempérament), l'éducation, le milieu social, les habitudes. Rien n'arrive sans cause ; le choix humain ne ferait pas exception.
\end{definition}

\textbf{La thèse de Spinoza.} Spinoza, dans l'\emph{Éthique} (1677), soutient que le libre arbitre est une \emph{illusion}. Sa formule est célèbre : les hommes « se croient libres parce qu'ils ont conscience de leurs désirs et ignorent les causes qui les déterminent ». Il compare l'homme à une pierre lancée dans les airs qui, si elle était consciente, croirait voler librement, alors que sa trajectoire est entièrement déterminée par la force du lancer et la pesanteur. De même, l'enfant désire le lait, l'adolescent désire la fête, et chacun croit choisir librement ce que son corps et son histoire le poussent à vouloir. La conscience de nos désirs sans la connaissance de leurs causes produit l'illusion de la liberté.

\textbf{La réponse de Sartre.} À l'opposé, Jean-Paul Sartre, dans \emph{L'existentialisme est un humanisme} (1946), affirme que « l'homme est \cle{condamné à être libre} ». Condamné, parce que nous n'avons pas choisi d'exister, mais qu'une fois jetés dans l'existence, nous sommes responsables de tout ce que nous faisons. Pour Sartre, aucune nature humaine, aucune éducation, aucun milieu ne décide à notre place : c'est nous qui donnons sens et valeur à ces facteurs en choisissant. Même refuser de choisir est encore un choix : l'élève qui dit « je n'y peux rien, tout le monde triche » a choisi de suivre les autres. On ne peut pas fuir la liberté, car fuir est encore un acte libre.

\begin{exemplebox}[Acte libre ou acte déterminé ?]
Un élève de Ngaoundéré choisit de tricher à l'examen du Probatoire. Deux lectures sont possibles. Lecture \emph{déterministe} (Spinoza) : la pression familiale, la peur de l'échec, l'exemple des camarades, la mauvaise préparation ont causé son geste ; il croit choisir, mais il ignore les causes qui le poussent. Lecture \emph{existentialiste} (Sartre) : malgré toutes ces pressions, il pouvait refuser ; d'autres élèves, dans les mêmes conditions, n'ont pas triché ; invoquer la pression sociale est une excuse pour se décharger de sa responsabilité. Le débat reste ouvert : c'est toute la question du libre arbitre.
\end{exemplebox}

\courssub{Les limites de la liberté : la liberté d'autrui}

La liberté, sous toutes ses formes, rencontre une limite intérieure : celle des autres. Ma liberté d'expression ne me permet pas d'insulter ; ma liberté de circulation ne me permet pas d'envahir la propriété d'autrui. La formule classique résume cette exigence : « \cle{ma liberté s'arrête où commence celle d'autrui} ».

Cette limite n'est pas une négation de la liberté, mais sa condition. Si chacun pouvait tout faire, le plus fort imposerait sa volonté au plus faible, et la liberté de tous serait détruite. La loi, en limitant chacun, protège donc la liberté de tous : c'est le paradoxe fécond de la vie en société.

\begin{attention}[Ne pas opposer liberté et loi]
Erreur fréquente dans les dissertations : présenter la loi comme l'ennemie de la liberté. Pour Kant, c'est le contraire : la vraie liberté est l'\cle{autonomie}, c'est-à-dire l'obéissance à la loi que l'on se donne soi-même par la raison. Celui qui suit ses seuls désirs n'est pas libre mais esclave de ses penchants ; celui qui agit selon une règle qu'il a rationnellement voulue pour tous est seul véritablement libre. La loi juste n'enchaîne pas la liberté : elle la réalise.
\end{attention}

\begin{popfigure}
\begin{center}
\begin{tikzpicture}[scale=1]
% poutre centrale
\draw[line width=2pt,popdark] (-3.2,0) -- (3.2,0);
\draw[line width=2pt,popdark] (0,0) -- (0,-0.5);
\draw[fill=popdark] (0,0.15) circle (0.12);
\draw[line width=2pt,popdark] (0,-0.5) -- (0,-1.6);
\draw[line width=2pt,popdark] (-0.9,-1.6) -- (0.9,-1.6);
% plateaux
\draw[line width=1.2pt,popblue] (-3.2,0) -- (-4.0,-1.1) (-3.2,0) -- (-2.4,-1.1);
\draw[line width=1.2pt,popblue] (-4.35,-1.1) arc (180:360:0.95);
\draw[line width=1.2pt,poporange] (3.2,0) -- (4.0,-1.1) (3.2,0) -- (2.4,-1.1);
\draw[line width=1.2pt,poporange] (2.45,-1.1) arc (180:360:0.95);
% étiquettes
\node[popblue,font=\bfseries,align=center] at (-3.2,-2.3) {Ma liberté\\individuelle};
\node[poporange,font=\bfseries,align=center] at (3.2,-2.3) {La liberté\\d'autrui};
\node[popmuted,align=center,font=\itshape] at (0,-2.3) {Équilibre garanti\\par la loi};
% flèche d'équilibre
\draw[<->,line width=1.4pt,popgreen] (-2.2,-3.1) -- (2.2,-3.1);
\node[popgreen,font=\itshape] at (0,-3.55) {« Ma liberté s'arrête où commence celle d'autrui »};
\end{tikzpicture}
\end{center}
\legende{La balance des libertés : la loi joue le rôle de pivot qui maintient l'équilibre entre ma liberté et celle d'autrui. Sans ce point d'appui commun, l'une écraserait l'autre.}
\end{popfigure}

\courssub{Étude de cas : la liberté de la presse au Cameroun}

Appliquons ces distinctions à un cas concret. La \cle{liberté de la presse} est une liberté politique : elle permet aux journalistes d'informer les citoyens, de critiquer la gestion publique et de contribuer au débat démocratique. Sans elle, les citoyens voteraient sans information, et les abus de pouvoir resteraient cachés.

Mais cette liberté a des \cle{limites légales}. La loi camerounaise, comme toutes les lois des États de droit, punit notamment :
\begin{itemize}
\item la \cle{diffamation} : publier de fausses accusations qui portent atteinte à l'honneur d'une personne ;
\item l'atteinte à la vie privée et la publication de fausses nouvelles ;
\item le \cle{trouble à l'ordre public} : les appels à la haine, à la violence ou à l'insurrection.
\end{itemize}

Un journaliste de Douala qui publie une enquête documentée sur la mauvaise gestion d'un marché public exerce légitimement sa liberté ; celui qui invente une accusation contre un adversaire politique commet une diffamation et engage sa responsabilité. On voit ici se rejoindre toutes les formes étudiées : la liberté politique (droit de publier) est limitée par la liberté d'autrui (droit à l'honneur), et l'exercice de cette liberté engage la responsabilité morale et juridique du journaliste — ce qui prépare la suite de notre leçon.

\begin{exoresolu}[Liberté ou licence ?]
\textbf{Énoncé.} Un étudiant déclare : « Puisque je suis libre, je peux publier sur les réseaux sociaux tout ce que je pense de mes enseignants, même des insultes. » Cette affirmation est-elle philosophiquement correcte ? Rédigez une réponse argumentée en cinq à huit lignes.
\tcblower
\textbf{Solution détaillée.} L'affirmation est incorrecte, car elle confond la \emph{liberté} avec la \emph{licence}. La liberté d'expression est bien un droit politique garanti par l'État de droit, héritier de l'article 10 de la Déclaration de 1789. Mais ce même article précise que la manifestation des opinions ne doit pas « troubler l'ordre public établi par la loi ». Or l'insulte et la diffamation portent atteinte à la liberté et à la dignité d'autrui : ma liberté s'arrête où commence celle d'autrui. De plus, selon Kant, la vraie liberté n'est pas de faire ce qu'on veut, mais d'obéir à la loi qu'on se donne rationnellement à soi-même ; or nul ne peut rationnellement vouloir une société où chacun insulte chacun. L'étudiant exerce donc non sa liberté, mais un caprice qui le rend responsable devant la loi et devant sa conscience.
\end{exoresolu}

\begin{aretenir}[L'essentiel de la section]
\begin{itemize}
\item La \cle{liberté physique} est l'absence d'obstacle matériel extérieur (le prisonnier n'est pas libre de ses mouvements).
\item La \cle{liberté morale} est la capacité de choisir entre le bien et le mal selon sa conscience.
\item La \cle{liberté politique} est l'ensemble des droits civiques garantis par l'État de droit (vote, expression, presse, circulation).
\item Le \cle{libre arbitre} est la faculté de se déterminer par sa seule volonté ; le \cle{déterminisme} (Spinoza) le nie, Sartre l'affirme radicalement : « l'homme est condamné à être libre ».
\item La liberté a pour limite la liberté d'autrui ; pour Kant, la loi n'est pas l'ennemie de la liberté mais sa réalisation (autonomie).
\end{itemize}
\end{aretenir}

\begin{exercice}[Pour s'entraîner]
1. Définissez le libre arbitre et le déterminisme, puis opposez les thèses de Spinoza et de Sartre en quelques lignes.

2. À l'aide d'un exemple camerounais de votre choix, montrez comment la liberté politique d'un citoyen rencontre la limite de la liberté d'autrui.

3. Sujet de dissertation : « La loi est-elle l'ennemie de la liberté ? » Construisez un plan en deux parties à partir des notions étudiées.
\end{exercice}

\begin{corrige}[Exercice 1]
Le libre arbitre est la faculté de se déterminer par sa seule volonté, indépendamment de toute cause extérieure : celui qui choisit pouvait choisir autrement. Le déterminisme est la doctrine selon laquelle tous nos actes sont causés par des facteurs (nature, éducation, milieu). Spinoza nie le libre arbitre : nous nous croyons libres parce que nous avons conscience de nos désirs tout en ignorant les causes qui nous déterminent, comme la pierre lancée qui croirait voler librement. Sartre affirme au contraire que l'homme est « condamné à être libre » : aucune cause ne décide à notre place, et même refuser de choisir est un choix. Le débat oppose donc l'illusion de la liberté (Spinoza) à la liberté inévitable (Sartre).
\end{corrige}

Ainsi, la liberté n'est pas une mais multiple : physique, morale, politique, et peut-être radicale si le libre arbitre existe. Mais qu'il s'agisse de l'élève qui triche, du journaliste qui publie ou du citoyen qui vote, une constante apparaît : là où il y a liberté, il y a choix, et là où il y a choix, il y a quelqu'un qui doit répondre de son acte. C'est précisément ce lien qu'il nous reste à établir : la liberté comme fondement de la responsabilité.

\coursec{Les formes de la responsabilité}

Nous avons vu, dans la section précédente, que la liberté n'est pas un simple pouvoir de faire ce que l'on veut : elle prend des formes diverses (liberté physique, liberté morale, liberté civique). Mais une question surgit immédiatement : si je suis libre, que dois-je à autrui et à moi-même ? Un agent libre peut-il agir sans jamais avoir à \emph{répondre} de ses actes ? C'est ici qu'intervient la notion de \cle{responsabilité}. Être libre, c'est être l'auteur de ses actes ; être responsable, c'est devoir en assumer les conséquences. La responsabilité est donc l'envers nécessaire de la liberté. Mais elle ne prend pas un visage unique : on peut répondre de ses actes devant sa conscience, devant la loi, ou devant la société. C'est ce que nous allons examiner.

\courssub{La responsabilité morale : répondre de ses actes devant sa conscience}

Commençons par une expérience intérieure que chacun a vécue. Un élève triche à une évaluation et obtient une bonne note. Personne ne l'a vu, aucune sanction ne le menace. Pourtant, le soir venu, il se sent mal à l'aise : une voix intérieure lui reproche son acte. Ce sentiment, c'est le \cle{remords} ; cette voix, c'est la \cle{conscience morale}.

\begin{definition}[Responsabilité morale]
La \cle{responsabilité morale} est l'obligation intérieure de répondre de ses actes devant sa propre \cle{conscience}. Elle se manifeste par le \cle{remords} (regret douloureux d'un acte passé) ou la \cle{culpabilité} (sentiment d'avoir mal agi), et par le sentiment de satisfaction lorsque l'on a bien agi.
\end{definition}

La responsabilité morale possède trois traits essentiels. D'abord, elle est \emph{intérieure} : son tribunal est la conscience, non un juge extérieur. Ensuite, elle est \emph{universelle} : même celui qui échappe à toute sanction sociale peut être rattrapé par sa conscience. Enfin, elle est \emph{exigeante} : elle ne juge pas seulement les actes, mais aussi les intentions. Kant insiste sur ce point : la valeur morale d'une action réside dans l'intention de faire son devoir, non dans ses résultats. Un commerçant de Douala qui est honnête par peur du contrôle fiscal agit conformément au devoir, mais non \emph{par} devoir : sa conscience morale n'est pas véritablement engagée.

\begin{exemplebox}[Le remords du mensonge]
Un jeune de Yaoundé accuse faussement son camarade d'avoir cassé le matériel de l'atelier. Le camarade est puni. Aucun témoin ne peut démentir le menteur : juridiquement, il ne risque rien. Moralement, pourtant, il est responsable : sa conscience lui reproche une injustice, et ce remords peut durer des années. La responsabilité morale atteint là où la loi ne peut pas atteindre.
\end{exemplebox}

\courssub{La responsabilité juridique : répondre de ses actes devant la loi}

La conscience ne suffit pas à organiser la vie en société. Il faut des règles communes, sanctionnées publiquement : c'est le domaine du droit.

\begin{definition}[Responsabilité juridique]
La \cle{responsabilité juridique} est l'obligation légale, imposée par la loi, de \emph{réparer un dommage} causé à autrui ou de \emph{subir une sanction} prévue par la loi. Elle se divise en deux grandes branches : la responsabilité civile et la responsabilité pénale.
\end{definition}

\begin{itemize}
\item La \cle{responsabilité civile} oblige à \emph{réparer} le dommage causé à autrui, le plus souvent par le versement de dommages et intérêts. Exemple : un mototaximan qui renverse un piéton à Bafoussam devra indemniser la victime. Ici, il ne s'agit pas de punir, mais de \emph{compenser} le préjudice.
\item La \cle{responsabilité pénale} oblige à \emph{subir une peine} (amende, emprisonnement) pour une infraction prévue par le Code pénal (vol, escroquerie, coups et blessures, détournement de deniers publics). Ici, il s'agit de \emph{sanctionner} une atteinte à l'ordre social.
\end{itemize}

\begin{attention}[Ne pas confondre responsabilité et culpabilité]
Erreur fréquente à l'examen : employer « responsabilité » et « culpabilité » comme des synonymes. La \cle{culpabilité} est un \emph{sentiment} éprouvé par l'agent (ou une faute reconnue par le juge) ; la \cle{responsabilité} est un \emph{statut} : le fait d'être tenu de répondre d'un acte. On peut être juridiquement responsable sans se sentir coupable (un employeur responsable des actes de son employé), et se sentir coupable sans être juridiquement responsable (un malade mental).
\end{attention}

\courssub{La responsabilité sociale et professionnelle}

Entre la conscience et la loi, il existe un troisième niveau : celui des \emph{rôles sociaux}. Occuper une fonction, c'est accepter des devoirs précis.

\begin{definition}[Responsabilité sociale et professionnelle]
La \cle{responsabilité sociale} est l'ensemble des devoirs attachés à un rôle ou à une fonction dans la société : le parent doit nourrir et éduquer ses enfants, l'enseignant doit transmettre les savoirs et protéger ses élèves, le fonctionnaire doit gérer honnêtement les deniers publics, l'élu doit servir l'intérêt général, l'entrepreneur doit respecter ses employés et l'environnement.
\end{definition}

Cette responsabilité a une particularité : elle peut être engagée \emph{sans faute personnelle directe}. Le chef d'établissement répond de la sécurité de son lycée ; le ministre répond de la gestion de son ministère. On parle alors de responsabilité \emph{du fait de la fonction}. C'est pourquoi, dans de nombreux pays, un ministre démissionne après un scandale dans son administration, même s'il n'y a pas participé : il assume symboliquement sa responsabilité de chef.

\courssub{Les conditions de la responsabilité et les causes d'irresponsabilité}

On ne peut pas rendre n'importe qui responsable de n'importe quoi. Pour qu'un agent soit légitimement tenu pour responsable, trois conditions doivent être réunies :

\begin{enumerate}
\item La \cle{conscience de l'acte} : l'agent doit savoir ce qu'il fait au moment où il agit.
\item La \cle{liberté du choix} : l'agent doit avoir pu agir autrement, sans contrainte extérieure irrésistible.
\item La \cle{connaissance des conséquences} : l'agent doit pouvoir prévoir, au moins approximativement, les effets de son acte.
\end{enumerate}

Lorsque l'une de ces conditions fait défaut, la responsabilité est \emph{atténuée} (diminuée) ou \emph{abolie} (supprimée). C'est ce que l'on appelle les causes d'\cle{irresponsabilité}.

\begin{center}
\begin{tabularx}{\linewidth}{|l|X|X|}
\hline
\rowcolor{popblueL} \textbf{Condition} & \textbf{Cause d'irresponsabilité correspondante} & \textbf{Exemple} \\
\hline
Conscience de l'acte & \textbf{La folie} (trouble psychique) : la responsabilité pénale est abolie & Un malade mental commet un acte violent sans comprendre ce qu'il fait \\
\hline
Liberté du choix & \textbf{La contrainte physique} ou morale irrésistible : responsabilité abolie & Un chauffeur forcé sous la menace d'une arme à transporter des marchandises volées \\
\hline
Connaissance des conséquences & \textbf{La minorité} : responsabilité atténuée ; \textbf{l'ignorance invincible} (impossible à surmonter) : responsabilité atténuée ou abolie & Un enfant de 10 ans met le feu par jeu, sans mesurer le danger \\
\hline
\end{tabularx}
\end{center}

\begin{exemplebox}[Un exemple chiffré : la majorité pénale au Cameroun]
Au Cameroun, la \cle{majorité pénale} est fixée à \textbf{18 ans}. En dessous de cet âge, le mineur est présumé ne pas avoir une pleine capacité de discernement : sa responsabilité pénale est \emph{atténuée}, et il relève d'une justice spéciale orientée vers l'éducation plutôt que vers la punition. De même, l'article du Code pénal camerounais relatif au trouble psychique dispose que n'est pas pénalement responsable celui qui était atteint, au moment des faits, d'un trouble psychique ayant aboli son discernement : la responsabilité est alors \emph{abolie}. Le droit applique donc concrètement les conditions philosophiques de la responsabilité.
\end{exemplebox}

\begin{savaistu}[L'opération « Épervier »]
Au Cameroun, l'opération anticorruption dite « \cle{Épervier} », lancée dans les années 2000, a conduit devant les tribunaux de hauts responsables accusés de détournement de deniers publics. Elle illustre concrètement la responsabilité pénale des agents publics : nul, quelle que soit sa fonction, n'est au-dessus de la loi. Elle pose aussi une question philosophique : la sanction suffit-elle à réparer le tort causé à la communauté ?
\end{savaistu}

\courssub{Schéma récapitulatif : les formes de la responsabilité}

\begin{popfigure}
\begin{tikzpicture}[
  grow=right,
  level 1/.style={sibling distance=3.6cm, level distance=3.2cm},
  level 2/.style={sibling distance=1.7cm, level distance=3.4cm},
  every node/.style={draw, rounded corners, align=center, font=\small, inner sep=3pt},
  edge from parent/.style={draw=popdark, thick, -latex}
]
\node[fill=popblue!70, text=white, font=\bfseries] {Responsabilité}
  child { node[fill=popgreenL, align=center]{Sociale et\\professionnelle}
    child { node[fill=popgreen!15, align=center]{Devoirs du rôle :\\parent, enseignant,\\fonctionnaire, élu} } }
  child { node[fill=poporangeL] {Juridique}
    child { node[fill=poporange!15, align=center]{Pénale : subir\\une sanction\\(peine, amende)} }
    child { node[fill=poporange!15, align=center]{Civile : réparer\\le dommage\\(indemnisation)} } }
  child { node[fill=poppurpleL] {Morale}
    child { node[fill=poppurple!15, align=center]{Répondre devant\\sa conscience :\\remords, culpabilité} } };
\end{tikzpicture}
\legende{Arbre des formes de la responsabilité. La responsabilité morale relève de la conscience ; la responsabilité juridique se divise en civile (réparer) et pénale (être puni) ; la responsabilité sociale découle des rôles et fonctions occupés dans la société.}
\end{popfigure}

\courssub{Étude de cas : le détournement des fonds d'une école rurale}

\begin{methode}[Analyser un cas de responsabilité en quatre étapes]
Pour traiter philosophiquement un cas concret (méthode exigible en dissertation et en étude de cas) :
\begin{enumerate}
\item \textbf{Identifier l'acte} : de quoi s'agit-il exactement ? Qui est l'agent ? Quel dommage est causé ?
\item \textbf{Vérifier la liberté de l'agent} : l'agent agissait-il en connaissance de cause ? Était-il contraint ? Les trois conditions (conscience, liberté, prévisibilité) sont-elles réunies ?
\item \textbf{Évaluer les conséquences} : quel préjudice matériel, moral et social l'acte a-t-il produit ?
\item \textbf{Conclure} : quelles formes de responsabilité (morale, civile, pénale, sociale) sont engagées, et à quel degré ?
\end{enumerate}
\end{methode}

\begin{exoresolu}[Le fonctionnaire et l'école rurale]
\textbf{Énoncé.} Un fonctionnaire chargé de gérer les fonds destinés à la construction d'une école rurale dans la région de l'Extrême-Nord détourne une partie de l'argent à son profit. L'école n'est jamais achevée ; les enfants du village continuent d'étudier sous un hangar. Analysez les responsabilités engagées.
\tcblower
\textbf{Solution détaillée.} Appliquons la méthode en quatre étapes.
\begin{enumerate}
\item \emph{Identification de l'acte.} L'acte est un détournement de deniers publics : l'agent, un fonctionnaire, a utilisé à des fins personnelles des fonds publics. Le dommage est double : matériel (perte financière de l'État) et social (privation d'éducation pour des centaines d'enfants).
\item \emph{Vérification de la liberté.} Le fonctionnaire est adulte, sain d'esprit, non contraint : il connaissait la destination des fonds et pouvait prévoir les conséquences de son acte. Les trois conditions de la responsabilité sont pleinement réunies. Aucune cause d'irresponsabilité ne peut être invoquée.
\item \emph{Évaluation des conséquences.} Le préjudice dépasse la somme détournée : il compromet l'avenir d'une génération d'élèves et détruit la confiance des citoyens dans les institutions.
\item \emph{Conclusion.} Quatre responsabilités sont engagées. (a) \emph{Responsabilité pénale} : le détournement de deniers publics est une infraction punie par le Code pénal camerounais (prison et amendes, comme le montre l'opération « Épervier »). (b) \emph{Responsabilité civile} : le fonctionnaire devra restituer les fonds et réparer le préjudice. (c) \emph{Responsabilité professionnelle} : il a trahi les devoirs de sa fonction et encourt des sanctions disciplinaires (révocation). (d) \emph{Responsabilité morale} : devant sa conscience, il doit répondre d'une injustice faite aux enfants de son propre pays.
\end{enumerate}
\end{exoresolu}

\courssub{La responsabilité collective : un groupe peut-il être responsable ?}

Reste une question difficile : peut-on rendre responsable non pas un individu, mais un \emph{groupe} — une entreprise, un État, une communauté ?

D'un côté, seuls les individus ont une conscience et une volonté : dire « la société est responsable » risque de diluer la faute et de déresponsabiliser chacun (« ce n'est pas moi, c'est le système »). De l'autre côté, certains dommages sont produits par des organisations entières : une entreprise qui pollue le fleuve Wouri, un État qui viole les droits de ses citoyens. Le droit moderne reconnaît donc la \cle{responsabilité des personnes morales} : une entreprise ou une collectivité peut être condamnée à réparer ou à payer une amende. L'État lui-même peut être tenu de réparer les dommages causés par ses services (responsabilité administrative).

\begin{aretenir}[L'essentiel de la section]
\begin{itemize}
\item La responsabilité morale fait répondre l'agent devant sa conscience (remords, culpabilité).
\item La responsabilité juridique est l'obligation légale de réparer un dommage (civile) ou de subir une sanction (pénale).
\item La responsabilité sociale attache des devoirs aux rôles et fonctions (parent, enseignant, fonctionnaire, élu).
\item Trois conditions fondent la responsabilité : conscience de l'acte, liberté du choix, connaissance des conséquences. Leur absence (folie, contrainte, minorité, ignorance invincible) atténue ou abolit la responsabilité — distinction appliquée par le Code pénal camerounais (majorité pénale à 18 ans, trouble psychique).
\item Le droit reconnaît aussi une responsabilité collective des personnes morales (entreprises, État).
\end{itemize}
\end{aretenir}

\begin{exercice}[Pour s'entraîner]
Un chauffeur de bus, fatigué après une nuit sans sommeil, provoque un accident sur l'axe Yaoundé--Douala. (1) Identifiez les formes de responsabilité susceptibles d'être engagées. (2) La fatigue peut-elle constituer une cause d'atténuation de la responsabilité ? Justifiez à l'aide des trois conditions de la responsabilité.
\end{exercice}

\begin{corrige}[Exercice 1]
(1) Trois formes de responsabilité peuvent être engagées : la \emph{responsabilité pénale} (blessures ou homicides involontaires), la \emph{responsabilité civile} (indemnisation des victimes, éventuellement partagée avec l'employeur), et la \emph{responsabilité morale} (remords du chauffeur). La responsabilité professionnelle de l'entreprise de transport peut aussi être mise en cause si elle a imposé des horaires épuisants.
(2) La fatigue n'abolit pas la responsabilité : le chauffeur restait conscient de son acte et libre de refuser de conduire. Elle peut toutefois \emph{atténuer} la responsabilité si la contrainte venait de l'employeur (pression professionnelle), car la condition de liberté du choix est alors partiellement diminuée. L'analyse doit donc distinguer la part de responsabilité du chauffeur et celle de l'entreprise : c'est un bon exemple de responsabilité individuelle et collective combinées.
\end{corrige}

\coursec{La liberté, fondement de la responsabilité}

Nous avons distingué, dans la section précédente, les grandes formes de la responsabilité : responsabilité morale, juridique, civile, pénale, politique. Mais une question demeure en suspens : à quelle condition peut-on légitimement déclarer quelqu'un responsable ? Un juge de Yaoundé ne condamne pas n'importe qui n'importe comment ; un père ne gronde pas son enfant de la même manière selon que celui-ci a cassé un verre en trébuchant ou en le jetant par colère. Derrière ces pratiques quotidiennes se cache une intuition profonde : on ne répond que de ce qu'on a fait librement. C'est cette intuition qu'il faut maintenant transformer en thèse rigoureusement établie.

\courssub{La thèse centrale : pas de responsabilité sans liberté}

\begin{propriete}[Thèse fondamentale de la leçon]
La \cle{liberté} est la \textbf{condition nécessaire} de la \cle{responsabilité}. Autrement dit : on ne peut être tenu responsable que de ce que l'on a accompli librement. Sans liberté, il n'y a pas d'\cle{imputation} possible : un acte accompli sous la contrainte ou dans l'ignorance ne peut pas être moralement ni juridiquement attribué à son auteur apparent.
\end{propriete}

L'intuition de départ est simple. Si quelqu'un me pousse dans l'escalier et que je bouscule un passant, personne ne m'accusera : mon corps a causé le choc, mais je n'en suis pas l'\emph{auteur}. La responsabilité suppose donc davantage que le simple contact causal entre un corps et un événement : elle suppose que l'acte vienne de moi, de ma décision, de ma volonté. C'est précisément ce que nomme le concept d'imputation.

\begin{definition}[Imputation]
L'\cle{imputation} est le fait d'attribuer un acte à un agent reconnu comme en étant l'auteur \textbf{libre et conscient}. Imputer une faute à quelqu'un, c'est affirmer deux choses : premièrement, qu'il est bien la cause de l'acte ; deuxièmement, qu'il aurait pu agir autrement. L'imputation est donc le pont conceptuel entre la liberté et la responsabilité.
\end{definition}

La formule « il aurait pu agir autrement » est décisive. Dire d'un agent qu'il est responsable, c'est dire que l'acte n'était pas inévitable pour lui : il disposait d'une marge de choix. Retirer cette marge, c'est retirer du même coup la responsabilité. Voyons maintenant les trois arguments qui établissent cette thèse.

\courssub{Premier argument : la contrainte exclut la responsabilité}

Celui qui agit sous la \cle{contrainte} n'est pas responsable de son acte, parce que sa volonté n'a pas participé à ce qu'il a fait. La contrainte peut être physique (on me force la main) ou morale (on me menace de mort).

\begin{exemplebox}[L'otage et l'homme menacé]
Un braqueur, à Douala, oblige un employé de banque, un pistolet sur la tempe, à ouvrir le coffre. L'employé ouvre : a-t-il commis une complicité de vol ? Évidemment non. Son geste matériel a servi le braquage, mais sa volonté n'y était pour rien : le choix n'existait plus pour lui, ou plutôt le « choix » qui restait (obéir ou mourir) n'est pas un choix libre. Le droit, ici comme partout, reconnaît l'état de \textbf{contrainte} comme une cause d'irresponsabilité : on ne peut imputer un acte à celui dont la volonté a été annulée par la force ou par la menace.
\end{exemplebox}

Le raisonnement est le suivant : la responsabilité suppose un choix ; or la contrainte supprime le choix ; donc la contrainte supprime la responsabilité. Ce raisonnement confirme par l'absurde notre thèse : là où la liberté disparaît, la responsabilité disparaît avec elle — preuve que la première est bien la condition de la seconde.

\courssub{Deuxième argument : l'ignorance diminue la responsabilité}

Celui qui \textbf{ignore} ce qu'il fait n'est pas pleinement responsable. La liberté n'est pas seulement absence de contrainte extérieure : elle suppose aussi la \cle{lucidité}, c'est-à-dire la connaissance de ce que l'on fait et de ce que cela va produire.

\begin{exemplebox}[L'enfant et le malade mental]
Un enfant de cinq ans, à Bafoussam, joue avec des allumettes et met le feu à la maison. On réparera, on surveillera mieux, mais on ne le jugera pas comme un adulte : il ne pouvait pas mesurer la portée de son geste. De même, un malade mental en crise qui blesse un voisin n'est pas condamné comme un criminel ordinaire : le droit pénal reconnaît que son discernement était aboli. Dans les deux cas, la justice adapte la réponse (éducation, soins) au lieu de punir : la sanction perd son sens là où la conscience de l'acte fait défaut.
\end{exemplebox}

Remarquons la nuance : l'ignorance ne supprime pas toujours totalement la responsabilité, elle la \emph{diminue}. Un adolescent de seize ans est moins responsable qu'un adulte, mais plus qu'un enfant de cinq ans. La responsabilité croît avec la lucidité, comme elle croît avec la liberté : les deux progressent ensemble, ce qui confirme encore leur lien.

\courssub{Troisième argument : choisir librement, c'est s'exposer à répondre}

Le troisième argument renverse la perspective : si l'absence de liberté excuse, alors la présence de liberté \emph{engage}. Choisir librement, c'est par là même accepter d'être l'auteur de son choix, donc d'en répondre.

Jean-Paul \cle{Sartre} radicalise cette idée : l'homme est « condamné à être libre », c'est-à-dire qu'il ne peut pas ne pas choisir — même refuser de choisir est encore un choix. D'où sa formule célèbre : « je suis responsable de tout, sauf de ma \textbf{responsabilité même} ». Je n'ai pas choisi d'être responsable : je le suis du seul fait que je suis libre. Ma responsabilité n'est pas un fardeau que j'aurais accepté par contrat ; elle est collée à ma liberté comme l'ombre au corps. L'étudiant de Terminale qui choisit de ne pas réviser ne peut pas ensuite accuser le hasard de son échec : son choix le précède et l'attend.

\begin{aretenir}[Le lien logique en trois temps]
\begin{enumerate}
\item La contrainte supprime le choix, donc supprime la responsabilité.
\item L'ignorance diminue la lucidité, donc diminue la responsabilité.
\item La liberté pleine implique le choix, donc implique la pleine responsabilité.
\end{enumerate}
Conclusion : la responsabilité varie exactement comme la liberté — preuve que la liberté en est la condition.
\end{aretenir}

\courssub{La chaîne conceptuelle : de la liberté à la responsabilité}

Nous pouvons maintenant reconstituer la chaîne complète qui relie les deux notions. Tout part de la liberté, capacité de se déterminer soi-même ; la liberté s'exerce dans un choix entre plusieurs possibles ; le choix se traduit en acte ; l'acte produit des conséquences dans le monde ; et c'est parce que la chaîne entière remonte jusqu'à ma liberté que ces conséquences me sont imputables : je suis responsable.

\begin{popfigure}
\begin{center}
\begin{tikzpicture}[
  box/.style={draw=popblue, fill=popblueL, rounded corners=3pt, minimum height=1.1cm, minimum width=2.4cm, align=center, font=\small\bfseries},
  lastbox/.style={draw=poporange, fill=poporangeL, rounded corners=3pt, minimum height=1.1cm, minimum width=2.6cm, align=center, font=\small\bfseries},
  arr/.style={-{Stealth[length=3mm]}, very thick, popdark}
]
\node[box] (lib) at (0,0) {Liberté};
\node[box] (choix) at (3.2,0) {Choix};
\node[box] (acte) at (6.4,0) {Acte};
\node[box] (cons) at (9.6,0) {Conséquences};
\node[lastbox] (resp) at (13.2,0) {Responsabilité};
\draw[arr] (lib) -- (choix);
\draw[arr] (choix) -- (acte);
\draw[arr] (acte) -- (cons);
\draw[arr] (cons) -- (resp);
\node[font=\footnotesize\itshape, text=popmuted, align=center] at (1.6,-0.9) {je peux\\agir autrement};
\node[font=\footnotesize\itshape, text=popmuted, align=center] at (4.8,-0.9) {je décide\\en connaissance de cause};
\node[font=\footnotesize\itshape, text=popmuted, align=center] at (8.0,-0.9) {je suis la cause\\de l'effet};
\node[font=\footnotesize\itshape, text=popmuted, align=center] at (11.4,-0.9) {l'acte m'est\\imputé};
\end{tikzpicture}
\end{center}
\legende{La chaîne conceptuelle de la responsabilité : chaque maillon transmet l'imputabilité au suivant. Si le premier maillon (la liberté) est brisé — par la contrainte ou l'ignorance —, toute la chaîne s'effondre et la responsabilité disparaît.}
\end{popfigure}

Ce schéma permet de lire d'un coup d'œil toute la leçon : la responsabilité n'est pas un maillon premier, c'est le maillon \emph{final} d'une chaîne dont la liberté est le premier anneau. Brisez la liberté, et tout s'écroule.

\begin{exemplebox}[Le chauffeur de taxi-brousse]
Deux accidents sur la route Yaoundé--Bafoussam. \textbf{Cas 1} : un chauffeur surcharge volontairement son véhicule de passagers et de bagages, roule trop vite pour faire un voyage de plus, et provoque un accident. Il est pleinement responsable : il connaissait les risques, il a librement choisi le gain contre la sécurité. \textbf{Cas 2} : un autre chauffeur, dont le véhicule est entretenu régulièrement, subit une rupture soudaine de freins due à un défaut de fabrication invisible. Il n'est pas responsable de la même manière : il n'a ni choisi le défaut ni pu le connaître. Même type d'accident, imputations opposées : la différence ne tient pas au résultat, mais à la liberté et à la lucidité de l'agent.
\end{exemplebox}

\courssub{L'objection déterministe : et si personne n'était libre ?}

Mais une objection redoutable menace tout l'édifice. Le \cle{déterminisme} affirme que tout événement, y compris nos décisions, est l'effet nécessaire de causes antérieures : hérédité, éducation, milieu social, état du cerveau. Si mes « choix » ne font que dérouler mécaniquement ce que mon passé a programmé, alors personne n'agit jamais librement — et, par notre propre thèse, \textbf{personne n'est jamais responsable}. Punir serait aussi absurde que punir une pierre qui tombe.

\begin{attention}[Piège de la dissertation]
Ne concluez jamais brutalement que « le déterminisme supprime toute responsabilité, donc la responsabilité n'existe pas ». C'est oublier que le déterminisme est une \emph{thèse discutable}, non un fait établi : il est impossible de prouver que chacune de mes décisions était totalement inévitable. Inversement, ne rejetez pas le déterminisme d'un revers de main : il rappelle utilement que notre liberté est réelle mais \emph{limitée} (par le corps, le milieu, l'histoire). La bonne copie montre les limites de chaque thèse : le libre arbitre absolu ignore les conditionnements ; le déterminisme absolu rend incompréhensibles le remords, la délibération et toute la vie morale.
\end{attention}

La réponse décisive à l'objection est d'ordre pratique. Toute société humaine — et la société camerounaise ne fait pas exception — \textbf{présuppose} la liberté pour fonctionner : le droit distingue le coupable de l'irresponsable, la sanction n'a de sens que si le sanctionné aurait pu agir autrement, la récompense ne valorise que le mérite d'un effort librement accompli, et le remords que chacun éprouve témoigne que nous nous croyons nous-mêmes auteurs de nos actes. Même le déterministe, dans la vie quotidienne, gronde son enfant et remercie son bienfaiteur : il vit donc comme s'il croyait à la liberté. La responsabilité est ainsi un présupposé incontournable de la vie commune, avant d'être une thèse métaphysique démontrée.

\begin{center}
\begin{tabularx}{\linewidth}{|l|X|X|X|}
\hline
\rowcolor{popblueL} & \textbf{Thèse : le libre arbitre} & \textbf{Antithèse : le déterminisme} & \textbf{Synthèse : la responsabilité assumée} \\
\hline
\textbf{Idée} & L'homme choisit librement, il est cause première de ses actes. & Tout acte est l'effet nécessaire de causes antérieures ; la liberté serait une illusion. & La liberté est réelle mais située et limitée ; elle suffit à fonder l'imputation. \\
\hline
\textbf{Conséquence} & Responsabilité pleine et entière. & Plus de responsabilité : ni faute, ni mérite. & On répond de ses actes dans la mesure de sa liberté et de sa lucidité. \\
\hline
\textbf{Limite} & Ignore les conditionnements réels (milieu, éducation, corps). & Rend absurdes le droit, la sanction, le remords, la délibération. & Reconnaît les conditionnements sans dissoudre le sujet responsable. \\
\hline
\end{tabularx}
\end{center}

\courssub{Synthèse : la responsabilité, prix de la liberté}

Nous pouvons conclure. La responsabilité n'est pas une punition infligée de l'extérieur aux êtres libres : elle est le \textbf{prix même de la liberté}. Être libre, ce n'est pas faire ce qu'on veut sans comptes à rendre ; c'est au contraire accepter d'assumer ce que l'on fait, parce que c'est bien \emph{nous} qui le faisons. La liberté sans responsabilité ne serait que caprice ou toute-puissance enfantine ; la responsabilité sans liberté ne serait qu'injustice. Les deux notions se tiennent ou tombent ensemble.

\begin{exoresolu}[Analyse d'une situation concrète]
\textbf{Énoncé.} Un élève de Terminale technique, à Garoua, triche à l'examen blanc en recopiant le téléphone de son voisin. Surpris, il déclare : « Ce n'est pas ma faute, tout le monde triche, et puis mon père m'a dit que si je rate, il me chasse de la maison : j'étais obligé. » Discutez cette défense en mobilisant le lien entre liberté et responsabilité.
\tcblower
\textbf{Solution détaillée.} L'élève invoque deux excuses qui correspondent exactement aux deux grandes causes d'irresponsabilité étudiées. (1) « Tout le monde triche » relève du déterminisme social : le milieu me pousserait à tricher. Mais une pression d'entourage n'est pas une contrainte qui abolit le choix : d'autres élèves, dans la même salle, ne trichent pas. L'excuse est donc faible. (2) La menace du père ressemble à la contrainte morale de l'otage ; mais l'analogie ne tient pas : l'otage risque la mort immédiate, l'élève risque une sanction familiale future et discutable. Il conservait des possibles réels (réviser davantage, parler à ses parents, assumer un échec). Sa liberté n'était donc pas abolie, seulement rendue coûteuse. Conclusion : l'acte lui est imputable ; il est responsable de sa tricherie, car il a choisi — sous pression, certes, mais choisi. La pression peut \emph{atténuer} la faute, elle ne l'efface pas.
\end{exoresolu}

\begin{methode}[Articuler liberté et responsabilité dans une dissertation]
\begin{enumerate}
\item Ne traitez jamais les deux notions séparément : annoncez dès l'introduction que le problème est leur \emph{lien} (la liberté fonde-t-elle la responsabilité ?).
\item Définissez chaque notion, puis le concept-charnière : l'\cle{imputation}.
\item Établissez le lien par les deux épreuves négatives : la contrainte et l'ignorance suppriment ou diminuent la responsabilité.
\item Confrontez l'objection déterministe, puis répondez par le présupposé pratique de la vie sociale (droit, sanction, récompense).
\item Concluez par la synthèse : la responsabilité est le prix de la liberté ; être libre, c'est accepter d'assumer.
\end{enumerate}
\end{methode}

\begin{savaistu}[Complément Bac : Kant et l'autonomie]
Emmanuel \cle{Kant} offre la réconciliation la plus profonde des deux notions. Pour lui, la vraie liberté n'est pas de faire n'importe quoi, mais l'\cle{autonomie} : se donner à soi-même (\emph{autos}, soi ; \emph{nomos}, loi) sa propre loi morale, par la raison, au lieu de la recevoir de l'extérieur (passions, coutumes, autorités). Or celui qui se donne sa loi ne peut pas ensuite s'en désavouer : il en est l'auteur, il en répond. L'autonomie fait ainsi coïncider parfaitement liberté et responsabilité : je suis responsable de mes actes précisément parce que la loi qu'ils suivent est la mienne. L'être libre, chez Kant, n'est pas celui qui n'obéit à rien, mais celui qui n'obéit qu'à la loi qu'il s'est rationnellement prescrite — et qui, pour cela même, en assume toutes les conséquences.
\end{savaistu}

\begin{aretenir}[L'essentiel de la section]
La liberté est la condition nécessaire de la responsabilité : sans elle, pas d'imputation possible. La contrainte et l'ignorance, en supprimant ou en diminuant la liberté, suppriment ou diminuent la responsabilité — ce qui prouve le lien par l'absurde. L'objection déterministe échoue car toute vie sociale présuppose la liberté. Enfin, avec Kant, l'autonomie réconcilie les deux notions : être libre, c'est se donner sa loi et en répondre. La responsabilité est le prix de la liberté.
\end{aretenir}

\coursec{Introduction : problématisation et définitions}

\courssub{De l'expérience commune à la question philosophique}

\begin{exemplebox}[Situation déclenchante : l'accident de la route]
À un carrefour de Yaoundé, un conducteur de moto-taxi (\emph{benskin}) grille un feu rouge, percute un piéton et prend la fuite. Arrêté plus tard, il déclare : « Je n'avais pas le choix, j'étais pressé, le client hurlait. » Le piéton, lui, réclame justice : « Il a choisi de foncer, il doit payer. »
\end{exemplebox}

Cette scène banale condense le cœur de notre leçon. Le conducteur invoque la \cle{contrainte} (pression du client, urgence) pour nier sa \cle{liberté} et échapper à sa \cle{responsabilité}. La victime, elle, postule que l'agent a \emph{choisi} (liberté) et doit \emph{répondre} (responsabilité). La philosophie ne tranche pas le procès — c'est le rôle du juge — mais elle éclaire les concepts qui le fondent : \textbf{qu'est-ce qu'être libre ? Qu'est-ce qu'être responsable ? Et pourquoi la première conditionne la seconde ?}

\begin{definition}[Liberté (sens philosophique)]
Capacité pour un agent de s'autodéterminer, c'est-à-dire d'être la \emph{cause première} de ses actes, par choix réfléchi, sans contrainte extérieure insurmontable ni nécessité intérieure aveugle. Elle suppose l'\cle{autonomie} (se donner sa propre loi) et le \cle{libre arbitre} (pouvoir choisir entre plusieurs possibles).
\end{definition}

\begin{definition}[Responsabilité (sens philosophique)]
Obligation pour un agent moral de \emph{répondre} de ses actes devant autrui (tribunal, conscience, société) et d'en assumer les conséquences (réparation, sanction, blame ou louange). Elle repose sur l'\cle{imputabilité} : l'acte doit pouvoir être \emph{attribué} à son auteur comme source libre et consciente.
\end{definition}

\begin{propriete}[Le lien ontologique : liberté $\rightarrow$ responsabilité]
\begin{itemize}
    \item \textbf{Principe :} Nul n'est responsable que de ce qu'il a fait \emph{librement}. \emph{Nulla poena sine lege}, mais surtout \emph{nulla responsabilitas sine libertate}.
    \item \textbf{Conditions de la responsabilité (droit et morale) :}
    \begin{enumerate}
        \item \textbf{Liberté :} Absence de contrainte physique ou morale irrésistible.
        \item \textbf{Connaissance (discernement) :} Savoir ce que l'on fait et en mesurer la portée (exclusion : démence, minorité, ivresse involontaire).
        \item \textbf{Intentionnalité / Volonté :} Vouloir l'acte ou en accepter le risque (dol direct, dol éventuel, faute non intentionnelle).
    \end{enumerate}
    \item \textbf{Conséquence :} Si l'une de ces conditions fait défaut, la responsabilité s'atténue ou disparaît (irresponsabilité pénale, cause de non-responsabilité morale).
\end{itemize}
\end{propriete}

\begin{savaistu}[Étymologie : \emph{Respondere} et \emph{Liber}]
\begin{itemize}
    \item \textbf{Responsabilité} : du latin \emph{respondere} (« répondre, garantir »). Le responsable est celui qui \emph{répond} de ses actes, qui en est \emph{garant}.
    \item \textbf{Liberté} : du latin \emph{liber} (« libre, affranchi »), opposé à \emph{servus} (esclave). L'homme libre (\emph{libertas}) a la \emph{disponibilité} de soi-même ; il n'est pas la propriété d'un autre.
    \item \textbf{Autonomie} : grec \emph{autos} (soi) + \emph{nomos} (loi). Se donner sa loi soi-même (Kant). C'est la forme supérieure de la liberté, fondement de la responsabilité morale.
\end{itemize}
\end{savaistu}

\coursec{Les formes de la liberté}

\courssub{Liberté négative et liberté positive (Benjamin Constant, Isaiah Berlin)}

\begin{center}
\begin{tabularx}{\linewidth}{|l|X|X|}\hline
\textbf{Critère} & \textbf{Liberté négative (Liberté \emph{de})} & \textbf{Liberté positive (Liberté \emph{pour})} \\ \hline
\textbf{Définition} & Absence d'obstacles, de contraintes, d'ingérence extérieures. « Zone de non-ingérence. » & Capacité effective d'agir, de se réaliser, de se gouverner soi-même (autonomie). \\ \hline
\textbf{Question clé} & « Jusqu'où puis-je aller sans qu'on m'arrête ? » & « Suis-je capable de faire ce que je \emph{veux vraiment} faire ? » \\ \hline
\textbf{Auteurs} & Benjamin Constant, John Stuart Mill, Isaiah Berlin. & Rousseau, Kant, Hegel, Marx, T.H. Green. \\ \hline
\textbf{Exemple camerounais} & Liberté de circuler (pas de barrage illégal), liberté d'ouvrir un commerce (droit OHADA). & Avoir les \emph{moyens} d'étudier (bourse, école proche), de se soigner (assurance maladie), de voter \emph{éclairé} (éducation civique). \\ \hline
\textbf{Risque} & Laisser le fort écraser le faible (liberté du renard dans le poulailler). & Justifier l'ingérence de l'État « pour votre bien » (paternalisme, totalitarisme). \\ \hline
\end{tabularx}
\end{center}

\begin{aretenir}[Synthèse : deux faces d'une même pièce]
La liberté négative est la \textbf{condition nécessaire} (espace vide), la liberté positive en est la \textbf{réalisation effective} (plein pouvoir). Un élève a la liberté négative de ne pas venir en cours (personne ne l'attache) ; mais sans liberté positive (moyens de transport, santé, motivation, projet), cette liberté reste vide. L'État de droit doit garantir la première \emph{et} favoriser la seconde (éducation, santé, emploi).
\end{aretenir}

\courssub{Libre arbitre, déterminisme et compatibilisme : le débat métaphysique}

\begin{definition}[Libre arbitre (indéterminisme libertarien)]
Thèse selon laquelle la volonté humaine n'est pas soumise aux lois de la causalité physique universelle. L'agent est une \emph{cause non causée} (\emph{causa sui}), capable d'initier une action sans antécédent déterminant. (Descartes, Leibniz, Sartre).
\end{definition}

\begin{definition}[Déterminisme (dur / spinoziste)]
Thèse selon laquelle tout événement, y compris les décisions humaines, est l'effet nécessaire de causes antérieures (physiques, biologiques, psychologiques, sociales). La liberté n'est qu'une \emph{illusion} née de l'ignorance des causes qui nous déterminent (Spinoza, Holbach, Freud, neurosciences contemporaines).
\end{definition}

\begin{definition}[Compatibilisme (déterminisme mou / liberté conditionnelle)]
Thèse qui réconcilie déterminisme causal et responsabilité : la liberté ne s'oppose pas à la causalité, mais à la \emph{contrainte}. Je suis libre si j'agis \emph{conformément à mes désirs et ma raison}, même si ces désirs ont des causes. (Hume, Kant \emph{pratique}, Stuart Mill, juristes modernes).
\end{definition}

\begin{exemplebox}[Le voleur affamé : trois lectures]
\begin{itemize}
    \item \textbf{Libertarien :} Il \emph{a pu} ne pas voler. Il est pleinement responsable.
    \item \textbf{Deterministe dur :} La faim, la misère, son histoire \emph{l'ont déterminé} à voler. Il n'est pas responsable ; il faut soigner/éduquer, non punir.
    \item \textbf{Compatibiliste :} Il a volé \emph{volontairement} (pas de pistolet sur la tempe). Il est responsable \emph{juridiquement}, mais la peine doit tenir compte de la \emph{contrainte sociale} (circonstance atténuante). C'est la position du droit camerounais (Code pénal, causes d'atténuation).
\end{itemize}
\end{exemplebox}

\begin{attention}[Piège de l'examen : confondre « cause » et « contrainte »]
Dire « mes gènes / mon milieu / mon passé m'ont fait faire ça » invoque des \emph{causes}. La responsabilité exige l'absence de \emph{contrainte} (force majeure, menace de mort, hypnose, maladie mentale abolissant le discernement). Avoir des raisons (causes) d'agir, c'est être un agent rationnel ; être forcé, c'est ne pas être agent. Ne confondez pas \textbf{explication causale} et \textbf{exonération responsable}.
\end{attention}

\coursec{Les formes de la responsabilité}

\courssub{Responsabilité morale vs responsabilité juridique}

\begin{center}
\begin{tabularx}{\linewidth}{|l|X|X|}\hline
\textbf{Dimension} & \textbf{Responsabilité morale} & \textbf{Responsabilité juridique} \\ \hline
\textbf{Fondement} & Conscience, devoir-être, valeurs (honneur, dignité, justice). & Loi positive (Constitution, Codes, Traités), sanction étatique. \\ \hline
\textbf{Juge} & Soi-même (conscience), opinion publique, groupe d'appartenance. & Tribunaux (pénal, civil, administratif, constitutionnel). \\ \hline
\textbf{Sanction} & Remords, honte, blame social, exclusion symbolique, perte d'estime. & Peine (prison, amende), réparation (dommages-intérêts), nullité, destitution. \\ \hline
\textbf{Étendue} & Illimitée (je suis responsable de mes pensées, de mes omissions, du sort du monde). & Limitée par la loi (pas de peine sans texte, prescription, immunités). \\ \hline
\textbf{Exemple} & Ne pas visiter ses parents âgés (honte, rupture lien). & Ne pas payer la pension alimentaire (saisie sur salaire, prison). \\ \hline
\end{tabularx}
\end{center}

\begin{propriete}[Les régimes de responsabilité juridique (Droit camerounais / OHADA)]
\begin{itemize}
    \item \textbf{Responsabilité pénale :} Atteinte à l'ordre public (infractions : crimes, délits, contraventions). Principe : \emph{personnalité} (nul n'est responsable que de son fait), \emph{légalité} (pas de crime sans loi), \emph{non-rétroactivité} (sauf loi plus douce). Sanction : peine.
    \item \textbf{Responsabilité civile (délitique / extra-contractuelle) :} Art. 1382-1386 Code civil (inspiré). Faute + Dommage + Lien de causalité $\rightarrow$ Réparation intégrale (dommages-intérêts). Responsabilité du fait personnel, du fait d'autrui (parents/enfants, employeurs/employés), du fait des choses (garde).
    \item \textbf{Responsabilité contractuelle :} Non-exécution ou mauvaise exécution d'un contrat (force majeure, fait du prince, fait du créancier comme exonérateurs).
    \item \textbf{Responsabilité administrative :} Faute de service public (faute de service, faute personnelle détachable). Juridiction : Tribunal administratif.
    \item \textbf{Responsabilité politique :} Redevabilité des gouvernants (voir Section 5).
\end{itemize}
\end{propriete}

\courssub{Responsabilité individuelle et responsabilité collective}

\begin{definition}[Responsabilité individuelle]
L'agent unique répond de son acte personnel. C'est le modèle par défaut en droit moderne (Code pénal art. 2 : « Nul n'est responsable que de son fait personnel »).
\end{definition}

\begin{definition}[Responsabilité collective / solidaire]
Groupe tenu pour responsable (État, entreprise, association, famille élargie). En droit : responsabilité de l'État pour faits de ses agents (art. 1384 al. 1 Code civil), responsabilité solidaire des co-auteurs/complices (Code pénal art. 27-28). En morale : dette historique, devoir de mémoire, solidarité clanique.
\end{definition}

\begin{experience}[Débat : La responsabilité du groupe WhatsApp « Élèves Tle »]
Un message haineux est posté. L'auteur est identifié. L'admin n'a pas supprimé. Les 200 membres l'ont lu, 50 l'ont transféré.
\begin{itemize}
    \item \textbf{Auteur :} Responsabilité pénale pleine (Loi 2010/012 art. 9).
    \item \textbf{Admin :} Responsabilité de \emph{modération} (négligence). Peut être poursuivi comme complice (facilitation).
    \item \textbf{Transférateurs :} Responsabilité de \emph{propagation} (complicité par instruction ou recel). « Je n'ai fait que transférer » n'exonère pas.
    \item \textbf{Membres passifs :} Responsabilité \emph{morale} (devoir de signaler, devoir de ne pas valider par le silence). « Voir et ne rien dire, c'est cautionner. »
\end{itemize}
\end{experience}

\coursec{La liberté, fondement de la responsabilité}

\courssub{L'argument ontologique : pas de responsabilité sans liberté}

\begin{methode}[Démonstration guidée : Pourquoi la liberté est la condition \emph{sine qua non} de la responsabilité]
\begin{enumerate}
    \item \textbf{Prémisse 1 (Définition) :} La responsabilité = imputabilité = attribution d'un acte à un auteur comme \emph{source}.
    \item \textbf{Prémisse 2 (Analyse de l'acte) :} Un acte humain = mouvement corporel + intention (volonté). Sans intention (réflexe, contrainte, sommeil), c'est un \emph{événement} qui m'arrive, pas mon \emph{action}.
    \item \textbf{Prémisse 3 (Liberté = pouvoir faire autrement) :} Être source de son acte, c'est avoir pu \emph{ne pas le faire} (pouvoir choisir entre A et non-A). Si je \emph{dois} faire A (déterminisme absolu, contrainte physique), je ne suis pas la source, je suis un \emph{instrument}.
    \item \textbf{Conclusion :} \textbf{Si pas de liberté (pouvoir choisir), pas d'auteur véritable, donc pas d'imputabilité, donc pas de responsabilité.} $\square$
\end{enumerate}
\end{methode}

\courssub{Les grands auteurs : Kant, Sartre, Ricœur}

\begin{savaistu}[Trois figures de la liberté fondatrice]
\begin{itemize}
    \item \textbf{Kant (Fondement de la métaphysique des mœurs) :} \emph{Autonomie} = volonté se donnant la loi morale à elle-même (Impératif catégorique : « Agis uniquement d'après la maxime qui fait que tu puisses vouloir en même temps qu'elle devienne loi universelle »). \textbf{Liberté = Autonomie = Devoir.} Je suis responsable \emph{parce que} je me suis législateur.
    \item \textbf{Sartre (L'Être et le Néant, L'Existentialisme est un humanisme) :} \emph{L'homme est condamné à être libre.} Pas de nature humaine prédéfinie, pas de Dieu. Je \emph{choisis} mon essence par mes actes. En choisissant pour moi, je choisis pour \emph{tous les hommes} (universalisation). \textbf{Responsabilité absolue, angoissante, sans excuse.}
    \item \textbf{Paul Ricœur (La Mémoire, l'Histoire, l'Oubli) :} \emph{L'imputabilité} comme noyau de l'agir. Le sujet responsable est celui qui \emph{répond} de ses actes par sa parole (promesse, aveu, récit). La responsabilité est \emph{narrative} : je me tiens pour auteur de mon histoire.
\end{itemize}
\end{savaistu}

\begin{attention}[Nuance cruciale : Liberté $\neq$ Toute-puissance]
Être libre et responsable ne signifie pas « pouvoir tout faire » ni « tout contrôler ». Je suis responsable de mes \emph{choix} dans la \emph{exemplebox} qui est la mienne (facticité : corps, histoire, milieu, lois physiques). La responsabilité porte sur \emph{comment j'assume ma situation}, pas sur la situation elle-même. (Ex: Je ne suis pas responsable d'être né pauvre au Cameroun, mais je suis responsable de ce que je \emph{fais} de cette condition : résignation, travail, révolte, fraude, excellence).
\end{attention}

\coursec{Applications concrètes au Cameroun}

La section précédente a établi que la \cle{liberté} est le fondement ontologique de la \cle{responsabilité} : on ne peut être tenu pour responsable que des actes que l'on a accomplis librement, en connaissance de cause et sans contrainte insurmontable. Mais la philosophie ne saurait rester une spéculation abstraite ; elle doit éclairer le concret de l'existence. Au Cameroun, comme dans toute société, la rencontre entre la liberté individuelle et les exigences du vivre-ensemble se joue quotidiennement dans des situations précises : dans l'isoloir, sur la route, en classe, au bureau, dans la forêt, sur les réseaux sociaux. C'est à cette épreuve du réel que nous soumettrons maintenant nos concepts.

\courssub{La responsabilité du citoyen : entre droits constitutionnels et devoirs civiques}

\begin{savaistu}[La Constitution de 1996 : libertés et devoirs]
Le Préambule de la Constitution de la République du Cameroun (loi n°96-06 du 18 janvier 1996) proclame solennellement l'attachement du peuple camerounais aux \emph{libertés fondamentales} (liberté d'opinion, de presse, de réunion, d'association, de mouvement) et à la \emph{déclaration universelle des droits de l'homme}. Mais ce même Préambule affirme que « tout citoyen a des devoirs envers la Nation et la Société » : respect de la Constitution et des lois, défense de la patrie, protection de l'environnement, paiement de l'impôt, participation au développement national. La liberté constitutionnelle n'est donc pas un blanc-seing ; elle s'accompagne d'une responsabilité juridique et morale explicite.
\end{savaistu}

La responsabilité du citoyen camerounais s'exerce dans au moins cinq sphères interdépendantes. Le diagramme ci-dessous en donne une représentation synthétique.

\begin{popfigure}
\begin{tikzpicture}
    \def\radius{3.5cm}
    \def\spheres{
        Famille/90/popblue,
        École/162/popgreen,
        Travail/234/poporange,
        État/306/poppurple,
        Environnement/378/poppink}
    % Dessiner les parts du camembert
    \foreach \label/\angle/\color in \spheres {
        \pgfmathsetmacro{\start}{\angle}
        \pgfmathsetmacro{\end}{\angle+72}
        \draw[fill=\color!60, draw=white, thick] (0,0) -- (\start:\radius) arc (\start:\end:\radius) -- cycle;
        % Étiquette au centre de la part
        \pgfmathsetmacro{\mid}{(\start+\end)/2}
        \node[rotate=\mid, anchor=south, font=\bfseries\small, text=white] at (\mid:\radius*0.6) {\label};
    }
    % Cercle central
    \draw[fill=popcream, draw=popdark, thick] (0,0) circle (1.2cm);
    \node[font=\bfseries\small, align=center, text=popdark] at (0,0) {CITOYEN\\RESPONSABLE};
    % Légende externe
    \foreach \label/\angle/\color [count=\i] in \spheres {
        \pgfmathsetmacro{\mid}{(\angle+\angle+72)/2}
        \node[anchor=west, font=\small] at (5.5, 2.5-\i*1.2) {\textcolor{\color}{\rule{5mm}{3mm}} \label};
    }
\end{tikzpicture}
\legende{Les cinq sphères de la responsabilité citoyenne au Cameroun. Le citoyen libre est au centre ; ses devoirs rayonnent vers la famille, l'école, le travail, l'État et l'environnement.}
\end{popfigure}

\begin{itemize}
    \item \textbf{Famille :} La responsabilité commence par la solidarité communautaire. L'adulte libre assume l'éducation, la santé et la protection de ses enfants et de ses parents âgés. C'est l'application concrète du devoir de « protection de la famille » (Préambule, al. 12).
    \item \textbf{École :} L'élève (voir infra) et le parent d'élève portent la responsabilité de la scolarisation effective, de l'assiduité et de l'honnêteté intellectuelle.
    \item \textbf{Travail :} Le travailleur libre honore son contrat, produit un travail de qualité, respecte les normes de sécurité. L'employeur, libre d'entreprendre, assume la responsabilité sociale : salaires décents, conditions de travail décentes, non-discrimination.
    \item \textbf{État :} Voter (exercice de la liberté politique), payer l'impôt (devoir de contribution aux charges publiques), respecter le code de la route (sécurité d'autrui), dénoncer la corruption (devoir de probité).
    \item \textbf{Environnement :} Ne pas jeter les déchets n'importe où, participer aux opérations « Jeudi propre » (ou initiatives locales), signaler les pollutions.
\end{itemize}

\begin{exemplebox}[Voter, payer l'impôt, respecter le code de la route : trois figures de la responsabilité]
\begin{itemize}
    \item \textbf{Voter :} C'est l'acte par lequel le citoyen \emph{exerce} sa liberté pour \emph{choisir} ses représentants. Ne pas voter, c'est abandonner sa part de responsabilité dans la conduite de la cité à d'autres. Au Cameroun, l'inscription sur les listes électorales (ELECAM) est un acte volontaire ; le vote lui-même est un devoir civique (non sanctionné pénalement mais moralement exigé).
    \item \textbf{Payer l'impôt :} La liberté d'entreprendre et de gagner sa vie suppose la contrepartie du financement des services publics (routes, hôpitaux, écoles). La fraude fiscale est un vol envers la collectivité ; elle brise le pacte social.
    \item \textbf{Code de la route :} La liberté de circuler s'arrête où commence le danger pour autrui. Griller un feu rouge, rouler en état d'ivresse, refuser la priorité aux piétons : ce sont des atteintes à la responsabilité juridique (amendes, prison) et morale (mise en danger de la vie d'autrui).
\end{itemize}
\end{exemplebox}

\courssub{La responsabilité de l'élève : assiduité, honnêteté et l'épreuve du Baccalauréat}

L'élève en classe de Terminale technique est un agent moral en devenir. Sa liberté grandissante (choix de série, orientation, gestion de son temps) appelle une responsabilité proportionnée.

\begin{exemplebox}[La fraude à l'examen : double responsabilité, double sanction]
\textbf{Situation :} Un élève de Terminale TSE introduit une fiche de révision dans la salle d'épreuve de Mathématiques au Baccalauréat.
\textbf{Analyse philosophique :}
\begin{enumerate}
    \item \textbf{Liberté :} L'élève a \emph{choisi} de frauder. Il n'était pas contraint physiquement ; il a pesé le pour (réussir sans effort) et le contre (risque de sanction) et a décidé d'agir. Il est donc \emph{auteur} de l'acte.
    \item \textbf{Responsabilité morale :} Il trahit la confiance de l'institution, de sa famille, de lui-même. Il s'attribue une valeur (le diplôme) qu'il n'a pas méritée. C'est une faute contre l'\cle{honnêteté} et la \cle{justice}.
    \item \textbf{Responsabilité juridique/disciplinaire :} Les textes régissant les examens officiels (arrêtés MINESEC en vigueur) prévoient des sanctions lourdes : annulation de l'épreuve, exclusion temporaire ou définitive des examens officiels, radiation de l'établissement, poursuites pénales pour faux et usage de faux (Code pénal, art. 123 et s.).
    \item \textbf{Conséquence existentielle :} Le fraudeur se déresponsabilise ; il apprend que la tricherie paie (si non pris) ou que l'impunité n'existe pas (si pris). Dans les deux cas, il ne construit pas son autonomie.
\end{enumerate}
\end{exemplebox}

\begin{attention}[« Tout le monde le fait » n'est pas un argument]
La \cle{responsabilité collective} (le fait que la fraude soit massive dans un centre) ne \emph{jamais} efface la \cle{responsabilité individuelle}. Dire « les autres trichent aussi » est une \emph{rationalisation} (mécanisme de défense psychologique) pour fuir sa propre culpabilité. Le juge, comme la conscience morale, juge \emph{chaque} agent pour \emph{ses} actes. Au Bac, les sanctions sont individuelles : chaque fraudeur répond de \emph{sa} fiche, de \emph{son} téléphone, de \emph{sa} complicité.
\end{attention}

\begin{methode}[Appliquer les notions à une situation concrète : l'élève face à l'examen]
Pour analyser philosophiquement une situation (ex: fraude, absentéisme, violence à l'école), suivez ces quatre étapes :
\begin{enumerate}
    \item \textbf{Repérer l'agent :} Qui agit ? (L'élève, le surveillant, le proviseur). Quelle est sa marge de liberté ? (Contraintes : pression familiale, fatigue, peur de l'échec).
    \item \textbf{Identifier l'acte et sa nature :} Est-ce un acte libre, réfléchi, volontaire ? Y a-t-il intentionnalité ? (La fraude est préméditée ; l'oubli de copie non).
    \item \textbf{Déterminer les devoirs en jeu :} Devoir de sincérité, devoir de travail, devoir de respect des règles, devoir envers soi-même (construction de son avenir).
    \item \textbf{Évaluer les conséquences :} Pour l'agent (sanctions, conscience), pour les autres (injustice pour les honnêtes, dévalorisation du diplôme), pour l'institution (crédibilité du système éducatif).
    \item \textbf{Conclure sur la responsabilité :} L'agent est-il responsable ? Sur quel fondement (liberté, connaissance, absence de contrainte) ? Quelle réponse éthique et juridique apporter ?
\end{enumerate}
\end{methode}

\courssub{La responsabilité des gouvernants : la redevabilité (accountability)}

La liberté de gouverner (mandat électif, nomination) confère un pouvoir : décider de l'allocation des ressources publiques, orienter les politiques nationales. Ce pouvoir n'est pas une propriété privée ; c'est une \cle{délégation} de souveraineté. D'où le principe de \cle{redevabilité} : le gouvernant doit \emph{rendre compte} de sa gestion aux mandants (le peuple, le Parlement, la Cour Suprême, la CONAC, la CONSUPE).

\begin{aretenir}[Responsabilité politique vs responsabilité pénale]
\begin{itemize}
    \item \textbf{Responsabilité politique :} Sanction par le vote (non-réélection), motion de censure, démission morale. Elle porte sur l'\emph{incompétence}, la \emph{négligence}, le \emph{non-respect des promesses}.
    \item \textbf{Responsabilité pénale :} Sanction par la justice (prison, amendes, inéligibilité). Elle porte sur les \emph{détournements de fonds publics} (art. 134 Code pénal), \emph{corruption} (art. 138), \emph{concussion}, \emph{favoritisme} dans les marchés publics.
\end{itemize}
Au Cameroun, l'Opération \emph{Épervier} (depuis 2006) illustre la traque de la responsabilité pénale des gestionnaires publics. Mais la responsabilité politique (reddition de comptes annuelle devant l'Assemblée, transparence budgétaire) reste perfectible.
\end{aretenir}

\begin{exoresolu}[Analyse d'une situation : un maire et le marché public de voirie]
\textbf{Énoncé :} Un maire signe un marché de bitumage de 500 millions FCFA avec une entreprise qui n'a pas les capacités techniques. La route se dégrade après trois mois. L'enquête révèle que l'entreprise est détenue par le frère du maire.
\tcblower
\textbf{Solution détaillée :}
\begin{enumerate}
    \item \textbf{Liberté du maire :} Il a librement choisi l'attributaire (liberté d'appréciation dans le cadre du Code des marchés publics). Il connaissait le lien de parenté (connaissance).
    \item \textbf{Violation des devoirs :} Devoir de probité, d'impartialité, de gestion saine des deniers publics (loi sur les marchés publics, Constitution).
    \item \textbf{Responsabilité engagée :}
    \begin{itemize}
        \item \emph{Pénale :} Favoritisme (art. 139 Code pénal), détournement si surfacturation prouvée.
        \item \emph{Administrative :} Suspension, révocation par le Ministre de tutelle (Minatd).
        \item \emph{Politique :} Perte de légitimité, risque de non-réélection.
        \item \emph{Civile :} Obligation de réparer le préjudice financier de la commune.
    \end{itemize}
    \item \textbf{Conclusion :} La liberté de décision administrative n'est pas l'arbitraire. Elle est encadrée par des règles dont la violation engage la responsabilité personnelle du décideur. « J'ai suivi la procédure » ne suffit pas si le choix réel a été dicté par l'intérêt privé.
\end{enumerate}
\end{exoresolu}

\courssub{La responsabilité environnementale : déforestation et pollution du Wouri — qui répond ?}

L'environnement est le bien commun par excellence. Au Cameroun, deux crises majeures illustrent la difficulté d'assigner la responsabilité : la déforestation (Sud, Est, Adamaoua) et la pollution de l'estuaire du Wouri (Douala).

\begin{experience}[Enquête de terrain : la pollution du Wouri]
\textbf{Observation :} L'estuaire du Wouri reçoit quotidiennement : eaux usées non traitées (HYSACAM, CAMWATER, industries), déchets plastiques (ménages, marchés), hydrocarbures (stations-service, ports), boues de vidange.
\textbf{Acteurs libres et responsables :}
\begin{itemize}
    \item \textbf{État (Ministères de l'Environnement, de l'Eau, de l'Urbanisme) :} Liberté réglementaire (édicter normes, sanctionner). Responsabilité : défaillance du contrôle, absence de stations d'épuration fonctionnelles.
    \item \textbf{Industries (brasseries, savonneries, raffinerie) :} Liberté d'entreprendre. Responsabilité : rejet d'effluents non conformes (rejet direct en période de pluie, bypass des stations de traitement). \emph{Principe pollueur-payeur} (Loi n°96/12 du 5 août 1996, art. 17).
    \item \textbf{Collectivités (CUD, communes d'arrondissement) :} Liberté de gestion locale. Responsabilité : collecte des ordures défaillante, absence de décharges contrôlées, curage des caniveaux inexistant.
    \item \textbf{Citoyens/riverains :} Liberté de consommation et de mode de vie. Responsabilité : jet de sachets plastiques dans les caniveaux, raccordement sauvage aux égouts, vidange dans la nature.
\end{itemize}
\textbf{Problème philosophique :} La \cle{dilution de la responsabilité} (« c'est la faute de l'État / des entreprises / des autres ») paralyse l'action. La responsabilité environnementale exige une \emph{responsabilité partagée mais différenciée} : l'État régule et sanctionne, l'entreprise dépollue et paie, le citoyen trie et ne pollue pas.
\end{experience}

\begin{popfigure}
\begin{tikzpicture}[scale=0.9, every node/.style={font=\small}]
    % Titre
    \node[font=\bfseries\normalsize, text=popdark] at (0,5.5) {Chantier de construction d'une route : chaîne de responsabilités};
    % Acteurs
    \node[draw, rounded corners, fill=popblue!20, thick, align=center] (maitre) at (-5, 3.5){Maître d'ouvrage\\ (MINTP / Commune)};
    \node[draw, rounded corners, fill=popgreen!20, thick, align=center] (controle) at (0, 3.5){Bureau de contrôle\\ (Indépendant)};
    \node[draw, rounded corners, fill=poporange!20, thick, align=center] (entreprise) at (5, 3.5){Entreprise\\ (Attributaire)};
    \node[draw, rounded corners, fill=poppurple!20, thick, align=center] (labo) at (-2.5, 1.5){Laboratoire\\ (Géotechnique)};
    \node[draw, rounded corners, fill=poppink!20, thick, align=center] (population) at (2.5, 1.5){Populations\\ (Usagers / Riverains)};
    % Flux
    \draw[->, thick, popblue] (maitre) -- node[above, sloped] {Finance, définit cahier des charges} (entreprise);
    \draw[->, thick, popgreen] (controle) -- node[above, sloped] {Vérifie conformité, valide paiements} (entreprise);
    \draw[->, thick, poporange] (entreprise) -- node[above, sloped] {Exécute, engage sa responsabilité décennale} (maitre);
    \draw[->, thick, poppurple] (labo) -- node[above, sloped] {Teste matériaux (granulats, liant)} (controle);
    \draw[<->, thick, poppink] (population) -- node[above] {Signale malfaçons, utilise l'ouvrage} (entreprise);
    % Responsabilité finale
    \node[draw, rounded corners, fill=popgold!20, thick, align=center, font=\bfseries\small] at (0, -1) {SI LA ROUTE SE DÉGRADE EN 6 MOIS :\\ \textbf{Tous} répondent selon leur part de liberté et de devoir\\ (Contractuelle, Pénale, Civile, Morale)};
    % Légende visuelle "photo commentée"
    \draw[dashed, popmuted] (-6.5, 4.5) rectangle (6.5, -2.5);
    \node[align=left, text=popmuted, font=\tiny] at (-6.3, -2.2) {\textbullet\ Maître d'ouvrage : Choix de l'itinéraire, financement, réception des travaux.\\
    \textbullet\ Contrôleur : Indépendance vis-à-vis de l'entreprise ? Conflit d'intérêts ?\\
    \textbullet\ Entreprise : Qualité matériaux, respect normes (NF P 98-086), main-d'œuvre qualifiée.\\
    \textbullet\ Laboratoire : Indépendance ? Falsification de résultats ?\\
    \textbullet\ Population : Vigilance citoyenne, respect de l'ouvrage (pas de surcharge, pas de vol de glissières).};
\end{tikzpicture}
\legende{Photographie commentée (schématisée) : un chantier routier. Chaque acteur possède une marge de liberté technique et contractuelle qui fonde sa responsabilité spécifique. La qualité finale est le produit de l'interaction responsable de tous.}
\end{popfigure}

\courssub{La responsabilité dans la famille africaine : solidarité communautaire et devoirs envers les siens}

La conception occidentale de la responsabilité (individuelle, contractuelle, juridique) rencontre au Cameroun la conception communautaire : « Je suis parce que nous sommes » (\emph{Ubuntu}, philosophie bantoue). La liberté de l'individu s'inscrit dans un réseau de \cle{solidarité} obligée.

\begin{definition}[Responsabilité communautaire (conception africaine)]
C'est l'obligation morale, non écrite mais socialement contraignante, pour tout membre capable du lignage/clan, de subvenir aux besoins des membres vulnérables (enfants, vieillards, malades, chômeurs) et de contribuer aux événements collectifs (deuils, mariages, fêtes, cotisations). Elle repose sur la \cle{réciprocité} : j'aide aujourd'hui parce que j'aurai besoin d'aide demain.
\end{definition}

\begin{exemplebox}[Le « noir » (cotisation) de deuil : liberté ou contrainte ?]
Un jeune cadre à Douala reçoit l'appel de son village : « Il faut cotiser 50 000 FCFA pour les funérailles du patriarche. »
\begin{itemize}
    \item \textbf{Contrainte sociale :} Le refus entraîne l'exclusion symbolique (honte, « il a oublié d'où il vient »), le refus d'aide future, la rupture des liens.
    \item \textbf{Liberté réelle :} Le cadre \emph{peut} refuser (personne ne le met en prison). Mais le \emph{coût} de ce refus (perte d'identité, de protection, de réseau) est si élevé que la liberté de refuser est \emph{théorique}.
    \item \textbf{Responsabilité assumée :} Le cadre paie. Il transforme une contrainte extérieure en devoir intérieur : « Je le fais par amour, par respect, par devoir. » C'est là que naît la \emph{responsabilité authentique} (Sartre) : assumer librement ce que la tradition impose.
\end{itemize}
\end{exemplebox}

\begin{attention}[Ne pas confondre solidarité et irresponsabilité individuelle]
La solidarité familiale ne doit pas devenir un \emph{parasitisme} : « Mon oncle est fonctionnaire, il doit tout payer, je n'ai pas à travailler. » La responsabilité du jeune diplômé sans emploi est de \emph{chercher} activement, de \emph{créer} son activité, de ne pas user de sa situation de « neveu » comme d'une rente. La liberté de l'oncle (gagner sa vie) ne fonde pas le droit du neveu à l'assistanat.
\end{attention}

\courssub{La responsabilité numérique : propos haineux et fausses informations sur les réseaux sociaux}

L'espace numérique (Facebook, WhatsApp, X/Twitter, TikTok) offre une \cle{liberté d'expression} démultipliée : anonymat relatif, viralité instantanée, absence de frontières. Mais cette liberté technique ne supprime pas la responsabilité ; elle l'aggrave par la portée du dommage.

\begin{propriete}[Cadre juridique camerounais : Loi n°2010/012 du 21 décembre 2010 relative à la cybercriminalité]
\begin{itemize}
    \item \textbf{Art. 7 :} Atteinte à l'honneur et à la réputation par le biais des TIC (diffamation, injure) : peine de 6 mois à 2 ans de prison et/ou 500 000 à 5 000 000 FCFA d'amende.
    \item \textbf{Art. 8 :} Propagation de fausses nouvelles (fake news) de nature à troubler l'ordre public : peine de 3 mois à 3 ans de prison et/ou 1 000 000 à 10 000 000 FCFA.
    \item \textbf{Art. 9 :} Incitation à la haine tribale, religieuse, raciale : peine de 1 à 5 ans de prison et/ou 5 000 000 à 10 000 000 FCFA.
\end{itemize}
La responsabilité est ici \emph{personnelle} (l'auteur du post), mais peut être \emph{partagée} (administrateur de groupe WhatsApp qui ne modère pas, plateforme qui ne retire pas le contenu signalé).
\end{propriete}

\begin{exemplebox}[Cas concret : un message vocal WhatsApp tribaliste]
Un étudiant en Terminale enregistre un message vocal de 2 minutes : « Les gens de la région X sont des voleurs, il faut les chasser de notre quartier. » Il l'envoie sur un groupe « Étudiants Tle Douala » (200 membres). Le message est transféré sur 50 autres groupes en 1 heure.
\begin{enumerate}
    \item \textbf{Liberté :} Il a choisi le contenu, le ton, le moment, le destinataire. Nul ne l'a forcé.
    \item \textbf{Responsabilité pénale :} Incitation à la haine tribale (Art. 9 Loi 2010/012). Risque de prison ferme.
    \item \textbf{Responsabilité civile :} Dommages-intérêts aux victimes identifiées (atteinte honneur).
    \item \textbf{Responsabilité scolaire :} Exclusion définitive de l'établissement (Règlement intérieur, MINESEC).
    \item \textbf{Responsabilité morale :} Il a contribué à la division nationale, à la violence potentielle. Il a utilisé sa liberté (parole, technique) pour \emph{détruire} du lien social au lieu de le \emph{construire}.
\end{enumerate}
\end{exemplebox}

\courssub{Débat guidé : « La jeunesse camerounaise est-elle libre et responsable de son avenir ? »}

Ce débat structure la synthèse de la leçon. Il ne s'agit pas de donner un avis « pour » ou « contre », mais de \emph{problématiser} en croisant les analyses précédentes.

\begin{methode}[Conduite du débat philosophique (20 min)]
\begin{enumerate}
    \item \textbf{Clarification des termes :} « Jeunesse » (15-35 ans, définition MINJEC), « Libre » (absence de contrainte ? pouvoir d'agir ? autonomie ?), « Responsable » (répondre de ses actes ? construire son avenir ? répondre devant la nation ?), « Avenir » (individuel ? collectif ?).
    \item \textbf{Thèse 1 (Oui, elle est libre et responsable) :} Arguments : Accès à l'information (internet), droit de vote (18 ans), liberté d'entreprendre (statut auto-entrepreneur, programmes PIAASI, PAJER-U), mobilité (études à l'étranger, concours internationaux), engagement associatif (ONG, Junior Chamber). \emph{Exemple :} Jeunes agriculteurs de la Menoua qui transforment le manioc en farine certifiée.
    \item \textbf{Antithèse (Non, elle est aliénée et déresponsabilisée) :} Arguments : Chômage de masse (structurel), népotisme (accès à la fonction publique, marchés), éducation inadaptée (décalage formation/emploi), pauvreté (liberté formelle seulement), manipulation politique (jeunes « militants » instrumentalisés), fuite des cerveaux (désespoir). \emph{Exemple :} Diplômé Master 2 chauffeur de moto-taxi (benskin) à Yaoundé.
    \item \textbf{Synthèse (Dépassement dialectique) :} La jeunesse camerounaise possède une \emph{liberté réelle mais contrainte} (structurellement). Sa responsabilité n'est pas nulle, mais elle est \emph{partagée} avec l'État (créer les conditions de l'emploi), l'école (former à l'autonomie), la famille (soutenir sans étouffer), l'entreprise (recruter sur compétences). \textbf{La responsabilité de la jeunesse est de transformer ses contraintes en projets collectifs} (associations, coopératives, veille citoyenne, vote éclairé, innovation locale).
    \item \textbf{Conclusion provisoire :} Être libre et responsable, ce n'est pas « tout pouvoir » ; c'est « faire ce qu'on peut avec ce qu'on a, là où on est » (Th. Roosevelt), en refusant la victimisation comme l'arrogance.
\end{enumerate}
\end{methode}

\courssub{Bilan : liberté et responsabilité, conditions de la vie en société et du développement national}

\begin{aretenir}[Synthèse finale : Le couple indissociable]
\begin{center}
\begin{tabularx}{\linewidth}{|l|X|X|}\hline
\textbf{Dimension} & \textbf{Liberté (Condition)} & \textbf{Responsabilité (Exigence)} \\ \hline
\textbf{Individuelle} & Autonomie morale, choix de vie, conscience & Assumer ses actes, se construire, tenir sa parole \\ \hline
\textbf{Civique} & Vote, expression, association, entreprise & Payer l'impôt, respecter la loi, voter, contrôler l'élu \\ \hline
\textbf{Professionnelle} & Choisir son métier, innover, négocier & Compétence, éthique, respect contrats, sécurité \\ \hline
\textbf{Écologique} & User des ressources, habiter le territoire & Préserver, restaurer, transmettre (développement durable) \\ \hline
\textbf{Numérique} & Publier, communiquer, créer & Vérifier sources, respecter autrui, protéger données \\ \hline
\textbf{Familiale/Communautaire} & Recevoir solidarité, transmettre culture & Aider les siens, éduquer enfants, honorer aînés \\ \hline
\end{tabularx}
\end{center}
\end{aretenir}

\begin{propriete}[Liberté et responsabilité : le moteur du développement national]
Aucun pays ne s'est développé sans citoyens \emph{libres} d'innover, d'entreprendre, de critiquer, de choisir leurs dirigeants, ET \emph{responsables} dans la gestion des fonds publics, le respect de l'État de droit, la protection de l'environnement, l'honnêteté aux examens, la probité aux affaires. Le Cameroun de 2035 (Vision 2035 : « Pays émergent, démocratique et uni dans sa diversité ») ne se construira pas par décret, mais par l'exercice quotidien, difficile et courageux, de cette \cle{double exigence} : \textbf{Être libre, c'est être responsable. Être responsable, c'est honorer sa liberté.}
\end{propriete}

\coursec{Méthode : rédiger et justifier une démarche philosophique}

\courssub{La dissertation philosophique au Probatoire / Baccalauréat technique}

\begin{methode}[Structure type d'une dissertation (3h -- 4 parties)]
\begin{enumerate}
    \item \textbf{Introduction (20-25 min) :}
    \begin{itemize}
        \item \textbf{Accroche :} Citation, fait d'actualité camerounaise, paradoxe, définition étymologique.
        \item \textbf{Définition des termes du sujet :} Préciser le sens philosophique (pas le sens courant). Distinguer les acceptions.
        \item \textbf{Problématisation :} Transformer le sujet en question tensionnée (ex: « La liberté est-elle absence de limites ou autonomie responsable ? »).
        \item \textbf{Annonce du plan :} « Nous verrons d'abord que... (Thèse), puis que... (Antithèse), pour enfin montrer que... (Synthèse). »
    \end{itemize}
    \item \textbf{Développement (2h à 2h15) : 3 grandes parties (I, II, III), chacune en 2 ou 3 sous-parties (A, B, C).}
    \begin{itemize}
        \item \textbf{Thèse (I) :} Défendre l'affirmation du sujet avec arguments, auteurs, exemples concrets (Cameroun).
        \item \textbf{Antithèse (II) :} Montrer les limites, l'opposition, la complexité. Nuancer, contredire.
        \item \textbf{Synthèse (III) :} Dépasser l'opposition. Montrer la condition, la articulation, le niveau supérieur. \emph{Pas de simple résumé.}
    \end{itemize}
    \item \textbf{Conclusion (15-20 min) :}
    \begin{itemize}
        \item \textbf{Bilan synthétique :} Réponse directe à la problématique.
        \item \textbf{Ouverture :} Autre domaine (art, science, politique), autre auteur, perspective historique, question nouvelle.
    \end{itemize}
\end{enumerate}
\end{methode}

\begin{methode}[Le commentaire de texte (alternative au Probatoire)]
\begin{enumerate}
    \item \textbf{Présentation du texte :} Auteur, œuvre, date, contexte, thèse générale.
    \item \textbf{Analyse structurée (explication linéaire ou thématique) :} Suivre le fil du texte, expliquer \emph{chaque} argument, définir \emph{chaque} concept clé, montrer les enchaînements logiques (donc, mais, or, car).
    \item \textbf{Discussion critique (évaluation) :} Le texte est-il convaincant ? Quelles objections ? Quels contre-exemples (Cameroun) ? Quelles limites ?
    \item \textbf{Conclusion :} Portée et actualité de la thèse.
\end{enumerate}
\end{methode}

\begin{attention}[Critères d'évaluation au Bac (Grille type)]
\begin{itemize}
    \item \textbf{Connaissances (4-5 pts) :} Maîtrise des notions, auteurs, théories, vocabulaire précis.
    \item \textbf{Argumentation (6-8 pts) :} Pertinence, cohérence, enchaînement logique, exemples pertinents (Cameroun !), nuance.
    \item \textbf{Expression (4-5 pts) :} Français correct, style philosophique (connecteurs logiques : \emph{en effet, toutefois, par conséquent, c'est-à-dire, notamment}), lisibilité.
    \item \textbf{Respect des consignes :} Sujet traité, plan visible, introduction/conclusion soignées.
\end{itemize}
\end{attention}

\begin{exercice}[Entraînement au Probatoire : Sujet de dissertation type]
\textbf{Sujet :} « La liberté n'est pas le droit de faire n'importe quoi, c'est le devoir de faire ce qu'il faut. » Discutez cette affirmation en vous appuyant sur des exemples camerounais.
\textbf{Pistes de corrigé (structure) :}
\begin{enumerate}
    \item \textbf{Introduction :} Accroche (ex: code de la route), définition des termes (liberté = autonomie/absence contrainte ; devoir = obligation morale/juridique), problématique (la liberté se définit-elle par l'absence de limites ou par l'auto-limitation responsable ?), annonce du plan.
    \item \textbf{I. La liberté comme absence de contrainte (liberté négative) : droit de faire ce que l'on veut ?} Exemples : liberté d'expression (réseaux sociaux), liberté de circulation (moto-taxi), liberté de commerce (marchés). Limite : le « n'importe quoi » conduit au chaos (accidents, fraude, pollution, haine).
    \item \textbf{II. La liberté comme autonomie responsable (liberté positive) : devoir de faire ce qu'il faut.} Fondement : Kant (autonomie = se donner sa loi morale), Sartre (choisir pour tous les hommes). Exemples camerounais : payer l'impôt (devoir civique), ne pas frauder au Bac (devoir d'honnêteté), planter un arbre (devoir écologique), voter éclairé (devoir politique).
    \item \textbf{III. La tension concrète au Cameroun : éduquer à la liberté responsable.} Rôle de l'école (philosophie, éducation à la citoyenneté), de la famille, de l'État (justice, exemplarité des gouvernants), des médias. Le développement (Vision 2035) suppose ce passage de la licence à la liberté responsable.
    \item \textbf{Conclusion :} Résumé. Ouverture : la liberté responsable comme fondement de la \emph{fraternité} (3ème devise nationale).
\end{enumerate}
\end{exercice}

\begin{corrige}[Exercice n°1 : Éléments attendus pour la dissertation]
\textbf{Introduction :} « Votre liberté s'arrête où commence celle des autres » (Déclaration 1789, art. 4). Au Cameroun, le mot « liberté » (liberté de la presse, liberté de manifester, liberté d'entreprendre) est souvent brandi comme un droit absolu. Mais l'expérience du quotidien (embouteillages anarchiques, fraude aux examens, corruption, discours de haine tribale) montre qu'une liberté sans responsabilité détruit la société.
\textbf{Développement :}
\begin{itemize}
    \item \textbf{I.} Analyser la conception libérale classique (Locke, Constant) : liberté = non-ingérence. Montrer ses limites au Cameroun : anarchie routière (liberté de rouler = mort), « liberté » de polluer le Wouri = mort du fleuve, « liberté » de tricher = diplôme sans valeur.
    \item \textbf{II.} Analyser la conception républicaine/kantienne : liberté = obéissance à la loi qu'on s'est prescrite. Devoir = impératif catégorique. Exemples concrets : l'élève qui refuse de frauder \emph{même s'il peut} (autonomie), l'entrepreneur qui paie ses impôts \emph{même sans contrôle} (civisme), le jeune qui refuse le tribalisme \emph{même sous pression} (humanisme).
    \item \textbf{III.} Synthèse : La loi (Code de la route, Code pénal, Loi cybercriminalité, Code du travail) traduit le « devoir de faire ce qu'il faut » en contrainte extérieure pour ceux qui n'ont pas encore la contrainte intérieure (morale). L'éducation (cette leçon de philo) vise à internaliser la loi. Le développement national = somme des responsabilités individuelles librement assumées.
\end{itemize}
\textbf{Conclusion :} La citation (souvent attribuée à P. Claudel ou J. Maritain) résume l'éthique de la citoyenneté : la vraie liberté n'est pas l'indifférence aux règles, c'est la force morale de s'y soumettre librement pour le bien commun.
\end{corrige}

\coursec{Méthode : rédiger et justifier une démarche philosophique}

Après avoir appliqué les notions de liberté et de responsabilité à des situations concrètes du Cameroun — le conducteur de moto-taxi de Douala, l'électeur de Yaoundé, le chef de famille responsable de ses enfants —, il reste une dernière étape, décisive pour le Probatoire et le Baccalauréat technique : savoir \emph{rédiger} cette réflexion de manière ordonnée et \emph{justifier} chacune de ses affirmations. Car à l'examen, on ne vous demande pas seulement de penser juste : on vous demande de le \cle{prouver par écrit}. Cette section vous donne les outils de la dissertation philosophique, appliqués au sujet central de la leçon.

\courssub{La structure générale de la dissertation philosophique}

\begin{definition}[La dissertation philosophique]
La \cle{dissertation philosophique} est un exercice écrit qui consiste à discuter un problème posé sous forme de question, en définissant les termes du sujet, en confrontant des thèses argumentées et en aboutissant à une réponse personnelle et justifiée. Elle comporte trois moments : l'\cle{introduction}, le \cle{développement} (en deux ou trois parties) et la \cle{conclusion}.
\end{definition}

L'intuition est simple : une dissertation ressemble à un voyage. L'introduction indique d'où l'on part et où l'on va ; le développement est le trajet lui-même, avec ses étapes ; la conclusion est l'arrivée, où l'on fait le bilan du chemin parcouru. Un correcteur du Baccalauréat technique juge d'abord la clarté de ce voyage : si le lecteur se perd, la copie est perdue.

\begin{propriete}[Les trois moments obligatoires]
\begin{enumerate}
\item \textbf{Introduction} : accroche (amener le sujet), définition des termes, problématique (la question qui guide toute la réflexion), annonce du plan.
\item \textbf{Développement} : deux ou trois parties, chacune contenant une thèse, des arguments, des exemples et des références à des auteurs.
\item \textbf{Conclusion} : réponse claire et nuancée à la problématique, bilan des arguments décisifs, ouverture éventuelle.
\end{enumerate}
\end{propriete}

\begin{popfigure}
\begin{center}
\begin{tikzpicture}[font=\small, scale=0.9]
% Entonnoir : du général au particulier (introduction)
\fill[popblueL] (0,4.6) -- (5,4.6) -- (4.2,3.4) -- (0.8,3.4) -- cycle;
\fill[popblue!45] (0.8,3.2) -- (4.2,3.2) -- (3.6,2.0) -- (1.4,2.0) -- cycle;
\fill[popblue!70] (1.4,1.8) -- (3.6,1.8) -- (3.1,0.7) -- (1.9,0.7) -- cycle;
\node[align=center] at (2.5,4.0) {\textbf{Général : le sujet dans}\\\textbf{ l'expérience humaine}};
\node[align=center] at (2.5,2.6) {\textbf{Définitions + problématique}};
\node[align=center, text=white] at (2.5,1.25) {\textbf{Plan annoncé}};
% Développement
\fill[poporangeL] (1.9,0.3) rectangle (3.1,-2.6);
\node[rotate=90, align=center] at (2.5,-1.15) {\textbf{Développement : thèses, arguments, exemples}};
% Retour au général (conclusion)
\fill[poppurpleL] (1.9,-2.9) -- (3.1,-2.9) -- (4.4,-4.1) -- (0.6,-4.1) -- cycle;
\node[align=center] at (2.5,-3.5) {\textbf{Conclusion : réponse nuancée,}\\\textbf{ ouverture}};
% Annotations latérales
\node[align=left, text=popdark] at (7.6,3.3) {Introduction :\\ du général\\ vers le problème};
\draw[->, popmuted, thick] (5.9,3.3) -- (4.4,3.3);
\node[align=left, text=popdark] at (7.6,-1.15) {Chaque partie :\\ thèse $\to$ argument\\ $\to$ exemple $\to$ référence};
\draw[->, popmuted, thick] (5.9,-1.15) -- (3.2,-1.15);
\node[align=left, text=popdark] at (7.6,-3.5) {Retour au général :\\ on répond à la question};
\draw[->, popmuted, thick] (5.9,-3.5) -- (4.5,-3.5);
\end{tikzpicture}
\end{center}
\legende{Schéma de la structure d'une dissertation réussie : l'introduction resserre progressivement le propos (du général vers le problème précis), le développement déploie l'argumentation, la conclusion rouvre vers une réponse d'ensemble.}
\end{popfigure}

\courssub{Définir les termes du sujet dès l'introduction}

La première erreur du candidat pressé est de foncer dans la discussion sans définir les mots du sujet. Or, en philosophie, les mots ordinaires sont trompeurs : « libre » peut vouloir dire « faire ce qu'on veut » (caprice) ou « agir en connaissance de cause » (autonomie). Définir, c'est \cle{fixer le sens} dans lequel on emploiera chaque terme pendant toute la copie.

\begin{methode}[Rédiger une introduction en quatre étapes]
\begin{enumerate}
\item \textbf{Amener le sujet} : partir d'un constat large et vrai (expérience commune, situation de la vie courante) qui conduit naturellement au sujet. Exemple : « Chaque jour, nous tenons les hommes pour responsables de leurs actes : nous les félicitons, nous les blâmons, parfois nous les punissons. »
\item \textbf{Définir les termes} : donner une définition claire et simple de chaque mot important du sujet. La \cle{liberté} est le pouvoir d'agir par soi-même, sans contrainte extérieure, selon sa propre volonté. La \cle{responsabilité} est l'obligation de répondre de ses actes et d'en assumer les conséquences.
\item \textbf{Poser le problème} : montrer que le sujet cache une difficulté, une tension. Ici : si l'homme n'était pas libre, pourrait-on encore le tenir pour responsable ? La responsabilité suppose-t-elle nécessairement la liberté ?
\item \textbf{Annoncer le plan} : indiquer en une ou deux phrases les étapes de la réflexion (« Nous verrons d'abord que..., puis nous examinerons..., enfin nous montrerons que... »).
\end{enumerate}
\end{methode}

\begin{attention}[Erreur fréquente au Bac : paraphraser au lieu de problématiser]
Beaucoup de candidats se contentent de \emph{répéter le sujet avec d'autres mots} : « La liberté est-elle le fondement de la responsabilité ? C'est une question importante car la liberté et la responsabilité sont deux notions essentielles... » Ce n'est pas une problématique, c'est une \cle{paraphrase}. Problématiser, c'est dégager une \emph{difficulté réelle} : un doute, un conflit entre deux thèses, un paradoxe. Bon test : si votre « problème » ne peut pas recevoir deux réponses opposées et défendables, ce n'est pas un problème.
\end{attention}

\courssub{Construire un argument : thèse, argument, exemple, référence}

Le développement n'est pas un déversement de souvenirs de cours. Chaque partie obéit à une mécanique rigoureuse, que l'on peut résumer par le schéma \textbf{T-A-E-R} : \cle{Thèse}, \cle{Argument}, \cle{Exemple}, \cle{Référence}.

\begin{propriete}[La structure d'un paragraphe argumenté]
\begin{enumerate}
\item \textbf{Thèse} : la position défendue dans la partie, énoncée en une phrase claire (« La liberté est la condition de la responsabilité »).
\item \textbf{Argument} : le raisonnement qui établit la thèse (« On ne peut reprocher à quelqu'un que ce qu'il aurait pu ne pas faire : imputer un acte suppose que l'agent avait le choix »).
\item \textbf{Exemple} : un cas concret qui illustre l'argument (« Un tribunal camerounais relaxe l'auteur d'un acte commis sous la menace d'une arme : la contrainte a annulé la liberté, donc la responsabilité »).
\item \textbf{Référence} : un auteur qui donne de l'autorité à la thèse (« Kant : il n'y a d'imputation morale que là où il y a une volonté libre »).
\end{enumerate}
\end{propriete}

\begin{exemplebox}[Un paragraphe type]
\emph{Thèse} : la responsabilité suppose la liberté de choix. \emph{Argument} : tenir quelqu'un pour responsable, c'est estimer qu'il aurait pu agir autrement ; or celui qui est contraint n'a pas ce pouvoir. \emph{Exemple} : l'élève puni pour une bêtise qu'il a délibérément commise est justement sanctionné, mais on ne punit pas celui qui a été bousculé et a brisé un objet malgré lui. \emph{Référence} : comme le montre Sartre, être libre, c'est être « condamné » à assumer ce que l'on fait de soi.
\end{exemplebox}

\courssub{Articuler les parties : de la définition des concepts à leur relation}

Une dissertation sur « liberté et responsabilité » ne doit pas être deux monologues juxtaposés (une partie sur la liberté, une sur la responsabilité, sans lien). L'enjeu du sujet est la \emph{relation} entre les deux notions. Le plan doit donc progresser : partir des définitions, confronter les thèses sur leur rapport, puis trancher.

\begin{exemplebox}[Sujet type Bac : « Faut-il être libre pour être responsable ? »]
\begin{itemize}
\item \textbf{I. Oui, la liberté fonde la responsabilité.} On n'est responsable que de ce qu'on a choisi ; l'imputation morale et juridique suppose le libre arbitre (Kant ; le droit pénal qui distingue l'acte volontaire de l'acte contraint ou de la folie).
\item \textbf{II. Limites : déterminisme et contraintes.} Nos actes dépendent de causes que nous ne maîtrisons pas : éducation, milieu social, passions, misère (Spinoza : les hommes se croient libres parce qu'ils ignorent les causes qui les déterminent). La liberté absolue est peut-être une illusion.
\item \textbf{III. Synthèse : la responsabilité comme exercice de la liberté.} La liberté n'est pas un donné mais une conquête : on devient responsable en s'efforçant de comprendre ce qui nous détermine et en assumant ses choix (Sartre : nous sommes responsables de ce que nous faisons de ce qu'on a fait de nous).
\end{itemize}
\end{exemplebox}

Remarquez la logique d'ensemble : la partie I expose la thèse spontanée, la partie II la met en difficulté, la partie III dépasse l'opposition. C'est le mouvement \cle{dialectique} (thèse, antithèse, synthèse), très apprécié des correcteurs parce qu'il montre que vous savez \emph{discuter} et non seulement réciter.

\courssub{Rédiger une conclusion justifiée}

La conclusion n'est pas une formalité. C'est le moment où vous \emph{répondez} enfin à la question posée dans l'introduction. Une conclusion sans réponse claire est une copie inachevée.

\begin{methode}[Justifier une conclusion en trois gestes]
\begin{enumerate}
\item \textbf{Reprendre la problématique} : rappeler la question initiale, pour montrer que tout le développement y répond (« Nous demandions si la liberté est le fondement de la responsabilité »).
\item \textbf{Rappeler les arguments décisifs} : résumer en deux ou trois phrases le résultat de chaque partie, en insistant sur ce qui a fait pencher la balance.
\item \textbf{Trancher avec nuance} : donner une réponse nette mais pas simpliste (« Oui, la liberté fonde la responsabilité, mais d'une liberté à conquérir sur les déterminismes, non d'une liberté absolue donnée d'avance »). On peut terminer par une \cle{ouverture} : une question nouvelle qui prolonge la réflexion (« Reste à savoir comment l'éducation peut former des êtres capables d'assumer cette liberté »).
\end{enumerate}
\end{methode}

\begin{attention}[Pièges de la conclusion]
Évitez trois fautes : introduire un \emph{argument nouveau} dans la conclusion (tout argument doit être développé dans le corps de la copie) ; répondre « oui et non » sans trancher (la nuance n'est pas l'indécision) ; écrire une conclusion qui contredirait votre développement (relecture indispensable).
\end{attention}

\courssub{Entraînement : rédaction guidée d'une introduction}

Appliquons la méthode au sujet de la leçon : \emph{« La liberté est-elle le fondement de la responsabilité ? »}

\begin{exoresolu}[Introduction rédigée et annotée]
Rédiger l'introduction du sujet : « La liberté est-elle le fondement de la responsabilité ? »
\tcblower
\textbf{Étape 1 — Amener le sujet.} « Dans la vie quotidienne, nous tenons spontanément les hommes pour responsables de leurs actes : nous louons le travailleur honnête, nous blâmons le voleur, les tribunaux punissent le criminel. Mais cette pratique repose sur un présupposé que nous n'examinons jamais : celui qui agit aurait pu agir autrement. »

\textbf{Étape 2 — Définir les termes.} « La liberté désigne le pouvoir d'agir selon sa propre volonté, sans contrainte extérieure ; la responsabilité est l'obligation de répondre de ses actes et d'en assumer les conséquences devant soi-même et devant autrui. »

\textbf{Étape 3 — Poser le problème.} « Or ce lien va de soi seulement en apparence : si nos actes étaient entièrement déterminés par notre éducation, notre milieu ou nos passions, serait-il encore juste de nous en tenir responsables ? La responsabilité suppose-t-elle nécessairement la liberté, ou peut-on être responsable sans être libre ? »

\textbf{Étape 4 — Annoncer le plan.} « Nous montrerons d'abord que la liberté est bien la condition de toute imputation ; nous examinerons ensuite les objections du déterminisme ; enfin, nous verrons que la responsabilité est moins un état qu'un exercice de la liberté. »

\emph{Commentaire} : chaque étape est identifiable, les définitions sont simples, la problématique oppose deux réponses possibles (oui / non), et le plan annoncé correspond exactement aux trois parties du développement.
\end{exoresolu}

\courssub{Les critères d'évaluation au Baccalauréat technique}

Le correcteur du Baccalauréat technique évalue votre copie selon quatre critères qu'il faut connaître pour viser juste.

\begin{center}
\begin{tabularx}{\linewidth}{|l|X|}\hline
\rowcolor{popblueL} \textbf{Critère} & \textbf{Ce que le correcteur attend} \\ \hline
\cle{Pertinence} & Tout ce qui est écrit répond au sujet posé ; pas de hors-sujet, pas de récitation du cours sans lien avec la question. \\ \hline
\cle{Cohérence} & Les parties s'enchaînent logiquement ; les transitions sont explicites ; la conclusion répond à l'introduction. \\ \hline
\cle{Exemples} & Chaque thèse est illustrée par au moins un exemple concret et précis (vie courante, histoire, réalités camerounaises). \\ \hline
\cle{Expression} & Français correct, phrases claires, vocabulaire philosophique employé à bon escient, orthographe soignée. \\ \hline
\end{tabularx}
\end{center}

\begin{aretenir}[L'essentiel de la méthode]
Une dissertation réussie = une introduction en quatre étapes (amener, définir, problématiser, annoncer), un développement où chaque partie suit le schéma thèse--argument--exemple--référence, une conclusion qui répond clairement à la problématique. La règle d'or : \emph{chaque phrase de la copie doit servir la réponse à la question posée}.
\end{aretenir}

\begin{exercice}[À vous de jouer]
Rédigez l'introduction complète du sujet : « Peut-on être responsable sans être libre ? » en suivant les quatre étapes de la méthode. Puis rédigez, pour la première partie de votre plan, un paragraphe complet respectant le schéma thèse--argument--exemple--référence.
\end{exercice}

\begin{corrige}[Exercice 1]
\emph{Éléments de correction.} Introduction attendue : amener le sujet par un constat (la société punit et récompense) ; définir liberté (pouvoir d'agir selon sa volonté) et responsabilité (devoir de répondre de ses actes) ; problématiser en opposant deux réponses (la responsabilité exige la liberté / on peut être tenu responsable d'actes déterminés, par exemple devant la loi ou devant sa conscience) ; annoncer un plan en deux ou trois parties. Paragraphe type : thèse (la responsabilité suppose la liberté de choix) ; argument (on ne peut reprocher que ce qui aurait pu être évité) ; exemple (le conducteur qui freine trop tard parce qu'un enfant surgit n'est pas responsable de la même façon que celui qui roulait ivre) ; référence (Kant : « Tu dois, donc tu peux » — le devoir moral suppose le pouvoir d'agir). Toute copie respectant cette architecture, avec une expression correcte, obtient une bonne note.
\end{corrige}

\newpage

\coursec{Activité d'intégration}

\courssub{Objectifs de l'activité}

Cette activité d'intégration, intitulée \emph{Le tribunal des idées : qui est responsable ?}, vise à mobiliser l'ensemble des savoirs et savoir-faire de la leçon dans une démarche collective puis individuelle. À l'issue de l'activité, l'élève doit être capable de :

\begin{itemize}
\item définir avec exactitude les concepts de \cle{liberté} et de \cle{responsabilité} et les distinguer des notions voisines (libertinage, fatalité, contrainte) ;
\item identifier, dans une situation concrète, l'\cle{agent}, son degré de liberté et les \cle{formes de responsabilité} engagées (morale, juridique, sociale) ;
\item appliquer la thèse philosophique selon laquelle \emph{la liberté est le fondement de la responsabilité} à des cas authentiques de la vie camerounaise ;
\item construire un schéma argumentatif en chaîne : liberté $\rightarrow$ choix $\rightarrow$ acte $\rightarrow$ conséquences $\rightarrow$ responsabilité ;
\item rédiger et justifier une conclusion personnelle et argumentée, conformément aux exigences de la dissertation au Probatoire et au Baccalauréat technique.
\end{itemize}

\courssub{Situation}

La classe de Terminale technique du lycée technique de Nkongsamba est transformée, le temps d'une double séance, en un \emph{tribunal des idées}. Le professeur de philosophie, M. Etoundi, remet à chaque groupe de quatre à cinq élèves un dossier documentaire contenant trois cas authentiques, inspirés de faits réels de la vie camerounaise. Chaque groupe joue le rôle d'un collège de juges : il doit examiner les faits, entendre les arguments, puis rendre un \emph{verdict argumenté}.

\textbf{Dossier n° 1 : la fraude à l'examen.} Lors de la session 2025 du Baccalauréat technique, série F1, à Douala, un candidat de 19 ans, Armand, est surpris par un surveillant en train de recopier des réponses contenues dans un téléphone portable dissimulé dans sa trousse. Convoqué devant la commission de discipline, il se défend en déclarant : \og Tout le monde le fait. Si je ne fraude pas, je serai le seul à échouer. Je n'avais pas le choix. \fg{} Les statistiques de l'Office du Baccalauréat du Cameroun (OBC) indiquent que plusieurs centaines de cas de fraude sont relevés chaque année sur l'ensemble du territoire national.

\textbf{Dossier n° 2 : l'accident de la route.} Sur l'axe Yaoundé--Bafoussam, un chauffeur de bus interurbain de 70 places, employé par une agence de voyage, provoque un accident qui fait 12 blessés. L'enquête établit que le chauffeur roulait à plus de \SI{110}{km/h} sur un tronçon limité à \SI{80}{km/h}. Le chauffeur explique que son employeur menace de licencier tout chauffeur qui n'effectue pas au moins trois rotations par jour, ce qui oblige à rouler vite. L'employeur, de son côté, déclare : \og Je n'ai jamais tenu le volant à sa place. \fg{}

\textbf{Dossier n° 3 : la fausse information.} Dans un quartier de Bamenda, une étudiante de 20 ans, Larissa, reçoit sur WhatsApp un message affirmant qu'un commerçant du quartier a empoisonné des produits alimentaires. Sans vérifier l'information, elle la relaie dans cinq groupes totalisant environ 800 membres. Le message devient viral ; une foule en colère saccage la boutique du commerçant, qui est innocent. Interrogée, Larissa affirme : \og J'ai seulement partagé ce que j'ai reçu. Ce n'est pas moi qui ai cassé la boutique. \fg{}

La question centrale du tribunal des idées est la suivante : \emph{dans chacun de ces cas, qui est responsable, de quoi, et pourquoi ?} Et, plus fondamentalement : \emph{la liberté est-elle le fondement de la responsabilité ?}

\courssub{Consignes}

\textbf{Travail en groupe (première séance).} Pour chacun des trois dossiers, le groupe répond aux questions suivantes :

\begin{enumerate}
\item Identifiez l'\cle{agent} de l'acte : qui a agi ? L'acte est-il imputable à une personne déterminée ?
\item Déterminez le \cle{degré de liberté} de l'agent : a-t-il agi par choix délibéré, sous contrainte, par ignorance ? La contrainte invoquée supprime-t-elle totalement la liberté, ou la réduit-elle seulement ?
\item Distinguez les \cle{formes de responsabilité} engagées : responsabilité morale (faute devant la conscience), responsabilité juridique (faute sanctionnée par la loi : civile et pénale), responsabilité sociale (faute jugée par la communauté). Justifiez chaque réponse.
\item Recherchez les \emph{causes atténuantes ou aggravantes} : pression du groupe, menace de licenciement, ignorance, précipitation. Ces circonstances modifient-elles le verdict ?
\item Construisez, pour chaque dossier, le schéma de la chaîne : \emph{liberté $\rightarrow$ choix $\rightarrow$ acte $\rightarrow$ conséquences $\rightarrow$ responsabilité}.
\item Rédigez un \emph{verdict argumenté} de dix lignes environ par dossier, comportant : une thèse (qui est responsable), au moins deux arguments, une référence à un auteur ou à une définition du cours, et une conclusion claire.
\end{enumerate}

\textbf{Débat plénier (deuxième séance).} Chaque groupe présente ses verdicts. Les autres groupes peuvent objecter, à condition de formuler des arguments et non de simples opinions. Le professeur arbitre et note les désaccords au tableau.

\textbf{Travail individuel (fin de la deuxième séance).} Chaque élève rédige une conclusion personnelle de quinze à vingt lignes répondant à la question : \og La liberté est-elle le fondement de la responsabilité ? \fg{} La conclusion doit être \emph{justifiée} : elle doit s'appuyer sur au moins deux des trois dossiers et sur une thèse du cours.

\courssub{Ressources à mobiliser}

\begin{aretenir}[Notions et thèses utiles]
\begin{itemize}
\item \textbf{Définitions.} La \cle{liberté} est le pouvoir d'agir par soi-même, en l'absence de contrainte extérieure, selon un choix que l'on peut s'approprier. La \cle{responsabilité} est l'obligation de répondre de ses actes et d'en assumer les conséquences.
\item \textbf{Formes de la liberté.} Liberté physique (absence d'entrave), liberté morale (autonomie de la volonté), liberté politique (droits reconnus par la loi).
\item \textbf{Formes de la responsabilité.} Responsabilité morale, juridique (civile et pénale) et sociale.
\item \textbf{Thèse de Sartre.} L'homme est \og condamné à être libre \fg{} : même sous pression, il choisit la manière dont il répond à la situation ; invoquer la contrainte pour nier toute responsabilité est une \emph{mauvaise foi}.
\item \textbf{Thèse de Kant.} La liberté est l'autonomie : obéir à la loi morale que l'on se donne à soi-même. On ne peut imputer un acte qu'à un être capable de se déterminer par la raison.
\item \textbf{Point de méthode.} Une conclusion justifiée comporte : une réponse claire à la question, des arguments, des exemples tirés des dossiers, et une nuance éventuelle (la contrainte peut \emph{atténuer} la responsabilité sans la supprimer).
\end{itemize}
\end{aretenir}

\courssub{Démarche et corrigé commenté}

\begin{exoresolu}[Le tribunal des idées : résolution complète]
Résoudre les trois dossiers du tribunal des idées et rédiger une conclusion justifiée à la question : \og La liberté est-elle le fondement de la responsabilité ? \fg{}
\tcblower
\textbf{Étape 1 : analyse du dossier n° 1 (la fraude).} L'agent est Armand : c'est lui qui a introduit le téléphone et recopié les réponses. Son degré de liberté est élevé : rien ne le contraignait physiquement ; l'argument \og tout le monde le fait \fg{} décrit une pression sociale, non une contrainte irrésistible. La pression du groupe ne supprime pas le choix : d'autres candidats, dans la même salle, n'ont pas fraudé. Selon Sartre, dire \og je n'avais pas le choix \fg{} est une mauvaise foi : Armand a choisi de céder à la pression. Responsabilités engagées : morale (il trompe sa propre conscience et celle des autres), juridique (la fraude à l'examen est sanctionnée par la réglementation de l'OBC et par la loi pénale), sociale (il dévalorise le diplôme de tous). Chaîne : liberté (possibilité de travailler ou de frauder) $\rightarrow$ choix (frauder) $\rightarrow$ acte (recopier) $\rightarrow$ conséquences (annulation possible, injustice envers les candidats honnêtes) $\rightarrow$ responsabilité pleine et entière.

\textbf{Étape 2 : analyse du dossier n° 2 (l'accident).} L'agent immédiat est le chauffeur, qui tenait le volant et a choisi d'accélérer ; mais la menace de licenciement constitue une contrainte réelle qui \emph{réduit} sa liberté sans l'abolir : il pouvait refuser, dénoncer l'employeur, ou rouler moins vite en acceptant le risque professionnel. L'employeur, lui, est co-responsable : en imposant un objectif incompatible avec les limitations de vitesse, il a organisé les conditions de la faute. La responsabilité est donc \emph{partagée} : le chauffeur engage sa responsabilité pénale (excès de vitesse, blessures involontaires) et morale ; l'employeur engage sa responsabilité juridique (civile pour le dommage, voire pénale pour mise en danger de la vie d'autrui et conditions de travail illicites) et sociale. Ce dossier montre que la chaîne de la responsabilité peut remonter au-delà de l'auteur matériel de l'acte, jusqu'à celui qui a déterminé les conditions du choix.

\textbf{Étape 3 : analyse du dossier n° 3 (la fausse information).} Larissa n'est pas l'auteur du message, mais elle a \emph{choisi} de le relayer sans vérification : l'acte de partager est un acte libre, accompli en quelques secondes, dont elle pouvait prévoir les conséquences possibles (800 membres, rumeur accusatrice). Son excuse (\og j'ai seulement partagé \fg{}) illustre la dilution de la responsabilité dans la foule numérique : chacun se dit trop petit pour être responsable, alors que la violence collective résulte de l'addition des choix individuels. Responsabilités engagées : morale (négligence, manque de prudence), sociale (elle a contribué à la violence du quartier), et juridique (la diffusion de fausses informations est répréhensible par la loi sur la cybercriminalité). Le degré de responsabilité est moindre que celui de l'auteur du message, mais il n'est pas nul : la liberté du relais fonde une responsabilité du relais.

\textbf{Étape 4 : schéma général de la chaîne.}
\begin{center}
\begin{tabular}{|c|ccccc|}
\hline
\rowcolor{popblueL} Étape & Liberté & Choix & Acte & Conséquences & Responsabilité \\
\hline
Dossier 1 & pouvoir de ne pas frauder & frauder & recopie & injustice, sanction & pleine \\
\hline
Dossier 2 & liberté réduite par la menace & rouler vite & excès de vitesse & 12 blessés & partagée \\
\hline
Dossier 3 & pouvoir de vérifier ou de taire & relayer & diffusion virale & saccage & partielle mais réelle \\
\hline
\end{tabular}
\end{center}

\textbf{Étape 5 : conclusion individuelle rédigée (modèle).} \emph{La liberté est bien le fondement de la responsabilité : on ne peut demander à quelqu'un de répondre d'un acte que s'il a pu choisir de le faire ou de ne pas le faire. Le dossier de la fraude le montre clairement : Armand, qui pouvait ne pas frauder, est pleinement responsable, et son excuse (\og tout le monde le fait \fg{}) relève de ce que Sartre appelle la mauvaise foi. Cependant, les trois dossiers invitent à nuancer : lorsque la liberté est réduite par une contrainte réelle, comme la menace de licenciement du chauffeur, la responsabilité n'est pas abolie mais partagée avec celui qui a créé la contrainte. De même, Larissa n'a pas tout causé, mais son choix libre de relayer l'information engage sa responsabilité. Ainsi, la responsabilité est proportionnelle au degré de liberté : plus l'agent est libre, plus il est responsable ; moins il est libre, plus la responsabilité se déplace vers les auteurs de la contrainte. Dire \og je n'étais pas libre \fg{} n'est donc jamais une excuse automatique : il faut encore prouver que la contrainte était irrésistible.}
\end{exoresolu}

\courssub{Bilan de l'activité}

Cette activité a permis de mettre en évidence trois acquis essentiels de la leçon.

Premièrement, les élèves ont vérifié concrètement que la \cle{responsabilité} suppose la \cle{liberté} : dans chaque dossier, le verdict a été construit en mesurant d'abord le degré de liberté de l'agent. Là où la liberté est pleine (fraude), la responsabilité est pleine ; là où elle est réduite (chauffeur), la responsabilité est partagée ; là où elle est diluée dans la foule (relais numérique), elle demeure réelle mais proportionnée.

Deuxièmement, l'activité a montré que la responsabilité n'est pas un bloc unique : elle se décline en responsabilité \cle{morale}, \cle{juridique} (civile et pénale) et \cle{sociale}, qui peuvent se superposer ou se dissocier. Un même acte peut être moralement condamnable, juridiquement sanctionné et socialement stigmatisé — ou l'un sans l'autre.

Troisièmement, le débat plénier a exercé les élèves à un savoir-faire central de l'épreuve de philosophie : \emph{rédiger et justifier une conclusion}. Le verdict argumenté et la conclusion individuelle reproduisent, à petite échelle, la structure d'une dissertation : thèse, arguments, référence philosophique, nuance, réponse claire à la question posée. Les élèves ont enfin compris qu'invoquer la contrainte (\og tout le monde le fait \fg{}, \og c'est la faute du patron \fg{}, \og j'ai seulement partagé \fg{}) n'efface pas la responsabilité : elle oblige seulement à examiner, avec rigueur, jusqu'à quel point l'agent était libre.

\newpage

\coursec{Méthodes et exercices résolus}

\coursec{La liberté et la responsabilité}

\courssub{Introduction : problématisation et définitions}

\begin{definition}[Liberté]
Au sens courant, la liberté est l'état de celui qui n'est pas esclave, qui n'est pas enfermé, qui peut faire ce qu'il veut. Philosophiquement, la liberté est la \cle{puissance de se déterminer soi-même} par un choix conscient, sans contrainte extérieure ni nécessité intérieure aveugle. Elle suppose l'existence de \emph{possibles} entre lesquels l'agent tranche.
\end{definition}

\begin{definition}[Responsabilité]
Au sens courant, être responsable, c'est devoir répondre de ses actes, en assumer les conséquences. Philosophiquement, la responsabilité est l'\cle{obligation pour un agent moral de répondre de ses actes} devant un tribunal (juridique), devant sa conscience (morale) ou devant la collectivité (sociale), parce qu'il en est l'auteur libre.
\end{definition}

\begin{aretenir}[Le lien constitutif]
La responsabilité \textbf{présuppose} la liberté : on ne peut demander des comptes qu'à qui a pu choisir. \textbf{Pas de responsabilité sans liberté.} C'est la thèse centrale de l'unité : la liberté est le \emph{fondement} ontologique et moral de la responsabilité.
\end{aretenir}

\begin{exemplebox}[Situation d'entrée : l'élève et l'examen]
Un élève de Terminale technique à Yaoundé entre dans la salle d'examen du Probatoire. Il a le choix : réviser ses notes une dernière fois ou discuter avec son voisin. Il choisit de discuter. Résultat : il rate une question facile.
\begin{itemize}
\item \textbf{Liberté} : il a pu choisir (personne ne l'obligeait à parler). Il est l'auteur de son choix.
\item \textbf{Responsabilité} : il assume l'échec (morale) et subit la note (scolaire/juridique).
\item \textbf{Lien} : s'il avait été malade (contrainte) ou s'il avait eu un malaise (nécessité), il n'aurait pas été libre \emph{de se taire} et sa responsabilité serait atténuée.
\end{itemize}
\end{exemplebox}

\courssub{Les formes de la liberté}

\begin{definition}[Liberté physique (ou d'indifférence)]
Absence de contrainte \emph{extérieure} matérielle (chaînes, prison, enfermement, force physique). C'est la condition \textbf{nécessaire} mais non suffisante de la liberté morale. \emph{Exemple :} un prisonnier libéré recouvre sa liberté physique.
\end{definition}

\begin{definition}[Liberté politique (ou civile)]
Ensemble des droits garantis par la loi (Constitution, Déclarations des droits) : liberté d'expression, de circulation, de vote, de réunion, de conscience. Elle dépend de l'État de droit. \emph{Exemple :} l'article 1er de la Constitution camerounaise garantit les libertés fondamentales aux citoyens.
\end{definition}

\begin{definition}[Liberté morale (ou libre arbitre, autonomie)]
Pouvoir de se \emph{déterminer soi-même} par la raison et la volonté, en choisissant entre le bien et le mal, entre plusieurs possibles. C'est la liberté \textbf{proprement philosophique}.
\begin{itemize}
\item \textbf{Spinoza} : la liberté n'est pas le libre arbitre (illusion), mais la \cle{connaissance de la nécessité} : « Je suis libre quand je suis cause adéquate de mes actions » (connaissance des causes qui me déterminent).
\item \textbf{Kant} : la liberté est \cle{autonomie} : la volonté se donne sa propre loi (loi morale), indépendamment des inclinations (hétéonomie). « Devoir implique pouvoir ».
\item \textbf{Sartre} : l'homme est \cle{condamné à être libre} : « L'homme n'est rien d'autre que ce qu'il se fait. » Pas de nature humaine prédéfinie ; l'existence précède l'essence. Même ne pas choisir, c'est choisir.
\end{itemize}
\end{definition}

\begin{propriete}[Distinction fondamentale : Liberté \emph{de} / Liberté \emph{pour}]
\begin{itemize}
\item \textbf{Liberté de} (négative) : absence d'obstacles, non-ingérence (physique, politique). « Liberté \emph{de} circuler. »
\item \textbf{Liberté pour} (positive) : capacité effective d'agir, maîtrise de soi, autonomie, moyens réels d'accomplir sa volonté. « Liberté \emph{pour} réussir ses études. »
\end{itemize}
\textbf{Au Bac :} ne pas confondre les deux. Un chômeur a la liberté \emph{de} travailler (personne ne l'en empêche légalement) mais pas toujours la liberté \emph{pour} (manque de moyens, d'offres).
\end{propriete}

\courssub{Les formes de la responsabilité}

\begin{definition}[Responsabilité morale]
L'agent répond de ses actes devant \textbf{sa conscience} (for intérieur). Sanction intérieure : \cle{culpabilité} (si acte mauvais), \cle{remords}, ou \cle{satisfaction} (si acte bon). Elle concerne l'intention (le \emph{vouloir} faire) plus que le résultat. \emph{Exemple :} un mensonge non découvert pèse sur la conscience du menteur.
\end{definition}

\begin{definition}[Responsabilité juridique]
L'agent répond de ses actes devant \textbf{la loi} (tribunaux). Sanction extérieure : \cle{réparation} (dommages-intérêts) ou \cle{peine} (prison, amende). Elle suppose l'imputabilité (capacité juridique, discernement). \emph{Exemple :} le Code pénal camerounais (Loi n°2016/007) punit le vol (art. 318) et prévoit des causes d'irresponsabilité (art. 63-65 : démence, contrainte, minorité).
\end{definition}

\begin{definition}[Responsabilité sociale (ou professionnelle, collective)]
Devoirs attachés à une \cle{fonction}, un \cle{rôle} ou une \cle{mission}. L'agent répond devant la \textbf{collectivité} ou sa \cle{hiérarchie}. \emph{Exemples :} l'enseignant responsable de la sécurité de ses élèves ; le chef d'entreprise responsable de la santé de ses employés (Code du travail) ; le ministre responsable de sa politique devant le Parlement.
\end{definition}

\begin{aretenir}[Tableau récapitulatif]
\begin{center}
\begin{tabularx}{\linewidth}{|l|X|X|X|}\hline
\textbf{Forme} & \textbf{Tribunal} & \textbf{Sanction} & \textbf{Fondement} \\ \hline
Morale & Conscience & Culpabilité / Remords / Estime de soi & Intention libre \\ \hline
Juridique & Juge / Loi & Peine / Réparation & Imputabilité (discernement + liberté) \\ \hline
Sociale & Collectivité / Hiérarchie / Opinion & Révocation / Blâme / Perte de confiance & Contrat / Mandat / Rôle \\ \hline
\end{tabularx}
\end{center}
\end{aretenir}

\courssub{La liberté, fondement de la responsabilité}

\begin{propriete}[Thèse : La liberté est la condition \emph{sine qua non} de la responsabilité]
\begin{enumerate}
\item \textbf{Argument ontologique} : imputer un acte à quelqu'un, c'est affirmer qu'il en est la \cle{cause efficiente} par sa volonté. Si l'acte lui est imposé (contrainte, réflexe, destin), il n'en est que l'instrument, pas l'auteur.
\item \textbf{Argument moral (Kant)} : « Devoir implique pouvoir ». On ne peut exiger un devoir (être responsable) que si le sujet a la \cle{puissance} de l'accomplir (liberté/autonomie).
\item \textbf{Argument existentiel (Sartre)} : l'homme \emph{est} sa liberté. Fuir sa responsabilité (mauvaise foi), c'est nier sa liberté. « Nous sommes nos choix. »
\end{enumerate}
\end{propriete}

\begin{definition}[Causes d'irresponsabilité (ou d'atténuation)]
Situations où la liberté fait défaut, supprimant ou diminuant la responsabilité :
\begin{itemize}
\item \textbf{Contrainte physique absolue} (violence, force majeure) : l'agent est un pur instrument. \emph{Exemple :} conducteur poussé par un tiers sur un piéton.
\item \textbf{Contrainte morale irrésistible} (menace de mort immédiate sur soi ou ses proches) : la volonté est paralysée par la terreur.
\item \textbf{Trouble mental / Démence} (folie) : abolition du discernement (art. 63 Code pénal camerounais). Le sujet ne comprend pas la portée de ses actes.
\item \textbf{Minorité} : discernement incomplet selon l'âge (responsabilité atténuée, mesures éducatives privilégiées).
\item \textbf{Ignorance invincible} : l'agent ne \emph{peut pas} savoir (ex. : substance empoisonnée dans un verre qu'on lui tend). Distinguer de l'ignorance \emph{vincible} (négligence), qui n'excuse pas.
\end{itemize}
\end{definition}

\begin{aretenir}[Degrés de responsabilité]
La responsabilité n'est pas binaire (tout ou rien) mais \textbf{proportionnelle au degré de liberté} :
\begin{itemize}
\item \textbf{Responsabilité pleine} : agent libre, conscient, sans contrainte.
\item \textbf{Responsabilité atténuée} : liberté réduite (émotion violente, provocation, ignorance partiellement excusable, minorité proche de la majorité).
\item \textbf{Irresponsabilité} : liberté abolie (folie, contrainte absolue, enfance jeune).
\end{itemize}
\end{aretenir}

\courssub{Applications concrètes au Cameroun}

\begin{exemplebox}[Code de la route et corruption à Douala]
Un conducteur en infraction (feu grillé, permis non présenté) propose 2000 FCFA au policier pour « arranger ». Le policier accepte.
\begin{itemize}
\item \textbf{Conducteur} : libre de corrompre (personne ne le menace). Responsabilité morale (choix du mal) + juridique (corruption active, art. 134 Code pénal).
\item \textbf{Policier} : libre de refuser (devoir professionnel). Responsabilité morale (trahison de la mission) + juridique (corruption passive) + sociale (manquement à la fonction).
\item \textbf{Analyse} : l'argument « c'est le système qui oblige » est une \emph{mauvaise foi} (Sartre) : chacun garde une marge de liberté, même mince. La responsabilité sociale du policier aggrave sa faute.
\end{itemize}
\end{exemplebox}

\begin{exemplebox}[Héritage coutumier et droit moderne à Bafoussam]
Une veuve se voit spoliée de ses biens par la famille du défunt au nom de la coutume. Elle saisit le tribunal.
\begin{itemize}
\item \textbf{Famille} : invoque la tradition (contrainte sociale/coutumière). Mais la liberté de choisir le droit (loi n°2016/007 art. 1er : la coutume ne s'applique que si elle n'est pas contraire à la loi) existe. Responsabilité juridique (recel, abus de confiance).
\item \textbf{Veuve} : exerce sa liberté politique/juridique (accès à la justice). Responsabilité morale de défendre ses droits et ceux de ses enfants.
\item \textbf{État (Justice)} : responsabilité sociale de faire primer la loi sur la coutume discriminatoire.
\end{itemize}
\end{exemplebox}

\begin{exemplebox}[Échec au Probatoire et orientation]
Un élève rate le Probatoire technique. Le père accuse l'école ; l'élève accuse le programme ; l'enseignant accuse la paresse.
\begin{itemize}
\item \textbf{Analyse philosophique} : l'élève a-t-il eu la liberté \emph{pour} réussir (moyens, temps, santé) ? A-t-il usé de sa liberté \emph{de} travailler (choix quotidiens) ?
\item \textbf{Responsabilité partagée} : l'élève (auteur principal de son travail), la famille (cadre de vie), l'école (qualité pédagogique), l'État (moyens). Mais \emph{in fine}, c'est l'élève qui passe l'examen : sa liberté morale est le point nodal.
\end{itemize}
\end{exemplebox}

\courssub{Méthode : rédiger et justifier une démarche philosophique}

\begin{methode}[Analyser une situation concrète avec les notions de liberté et de responsabilité]
Pour appliquer les notions de l'unité à une situation de la vie courante :
\begin{enumerate}
\item \textbf{Lire et décrire} la situation : qui agit ? que fait-il ? dans quelles circonstances (contexte camerounais si possible) ?
\item \textbf{Identifier la liberté} : l'agent a-t-il pu choisir ? Avait-il conscience de ce qu'il faisait ? Était-il contraint (violence, ignorance, contrainte extérieure, trouble mental) ? Distinguer \emph{incitation} (pourboire, menace légère) et \emph{contrainte} (force physique, terreur panique).
\item \textbf{Identifier la responsabilité} : de quel type s'agit-il ? \cle{Morale} (jugement de la conscience), \cle{Juridique} (obligation de réparer ou sanction pénale), \cle{Sociale} (rôle, devoir professionnel, mandat) ?
\item \textbf{Relier les deux notions} : la responsabilité n'existe que si l'agent était libre. Vérifier si des \cle{causes d'irresponsabilité} (folie, contrainte absolue, minorité, ignorance invincible) diminuent ou suppriment la responsabilité.
\item \textbf{Conclure} par un jugement argumenté : l'agent est-il pleinement, partiellement ou pas responsable ? Justifier par le degré de liberté réel.
\end{enumerate}
\textbf{Réflexe Bac} : toujours se demander : « L'agent \emph{pouvait-il faire autrement} ? » Si oui, il est libre et donc responsable.
\end{methode}

\begin{exoresolu}[Le chauffeur de taxi à Douala]
\textbf{Énoncé} : À Douala, un chauffeur de taxi, pressé par un client qui lui promet un gros pourboire s'il arrive à l'heure, grille un feu rouge et renverse un piéton. Un passant affirme : « Ce n'est pas de sa faute, c'est le client qui le pressait. » Analysez cette situation à l'aide des notions de liberté et de responsabilité.
\tcblower
\textbf{Solution détaillée.}
\begin{enumerate}
\item \textbf{Description} : le chauffeur est l'agent ; il commet une infraction (feu rouge grillé) causant un dommage corporel (piéton renversé). La promesse de pourboire du client est la circonstance atténuante alléguée.
\item \textbf{Analyse de la liberté} : le chauffeur n'était ni fou, ni ignorant du code de la route, ni physiquement contraint (personne ne tenait le volant). La promesse d'un pourboire est une \emph{incitation} (mobile), non une \emph{contrainte} (cause efficiente nécessaire) : il \textbf{pouvait faire autrement} (refuser la course, rouler normalement, déposer le client). Il a agi librement et en connaissance de cause.
\item \textbf{Analyse de la responsabilité} :
\begin{itemize}
\item \cle{Juridique} : infraction au code de la route (mise en danger d'autrui, blessures involontaires), obligation de réparer le dommage (Code civil), sanction pénale possible (Code pénal).
\item \cle{Morale} : sa conscience peut lui reprocher d'avoir privilégié l'argent sur une vie humaine.
\item \cle{Sociale} : en tant que professionnel du transport, il a un devoir de sécurité accru envers les usagers.
\end{itemize}
\item \textbf{Lien liberté-responsabilité} : la responsabilité repose sur la liberté. Ici, aucune cause d'irresponsabilité ne s'applique : le chauffeur est \textbf{pleinement responsable}. Le client qui pressait porte une responsabilité \emph{morale indirecte} (complicité morale), mais il n'a pas agi à la place du chauffeur (pas de co-action physique).
\end{enumerate}
\textbf{Conclusion} : l'argument du passant est irrecevable juridiquement et moralement : une incitation n'abolit pas la liberté. Le chauffeur, libre de son choix, est pleinement responsable de ses actes devant la loi et devant sa conscience.
\end{exoresolu}

\begin{methode}[Définir un concept philosophique en trois temps (standard Bac)]
Pour définir correctement un concept (liberté, responsabilité) dans une copie :
\begin{enumerate}
\item Donner d'abord le \textbf{sens commun} (définition du dictionnaire, usage quotidien).
\item Donner ensuite la \textbf{définition philosophique} précise, en distinguant les acceptions (ex. liberté physique / politique / morale ; responsabilité morale / juridique / sociale).
\item Appuyer la définition sur une \textbf{référence d'auteur} brève (nom + thèse centrale, pas de longue citation).
\end{enumerate}
\textbf{Formule à retenir} : « Au sens courant, X désigne... ; philosophiquement, il faut distinguer... ; comme le montre Auteur, ... »
\end{methode}

\begin{exoresolu}[Définir la liberté et la responsabilité]
\textbf{Énoncé} : Définissez les concepts de liberté et de responsabilité en distinguant leurs différentes formes, puis illustrez chaque forme par un exemple camerounais.
\tcblower
\textbf{Solution détaillée.}

\textbf{1. La liberté.} Au sens courant, être libre, c'est ne pas être empêché de faire ce qu'on veut. Philosophiquement, on distingue :
\begin{itemize}
\item la \cle{liberté physique} : absence de contrainte extérieure (le prisonnier libéré à la prison centrale de Kondengui est physiquement libre) ;
\item la \cle{liberté politique} : possibilité reconnue par la loi d'exprimer ses opinions, de circuler, de voter (le citoyen camerounais jouit de libertés garanties par la Constitution du 18 janvier 1996, révisée) ;
\item la \cle{liberté morale} ou libre arbitre : pouvoir de se déterminer soi-même, de choisir entre plusieurs possibles en conscience. Sartre affirme que « l'homme est condamné à être libre » : même en refusant de choisir, nous choisissons.
\end{itemize}

\textbf{2. La responsabilité.} Au sens courant, être responsable, c'est devoir répondre de ses actes. Philosophiquement, on distingue :
\begin{itemize}
\item la \cle{responsabilité morale} : l'agent répond de ses actes devant sa conscience (le fonctionnaire qui refuse un pot-de-vin éprouve de la fierté ; celui qui accepte éprouve de la culpabilité) ;
\item la \cle{responsabilité juridique} : l'agent répond de ses actes devant la loi et doit réparer le dommage ou subir la sanction (le conducteur ivre condamné par le tribunal de grande instance) ;
\item la \cle{responsabilité sociale} : devoirs attachés à un rôle ou une fonction (le chef de chantier est responsable de la sécurité de ses ouvriers ; le délégué de classe est responsable de l'ordre).
\end{itemize}

\textbf{Conclusion} : liberté et responsabilité sont des concepts à plusieurs dimensions, qu'il faut toujours distinguer avant de les articuler.
\end{exoresolu}

\begin{methode}[Problématiser un sujet de dissertation (épreuve écrite)]
Pour construire une problématique à partir d'un sujet :
\begin{enumerate}
\item \textbf{Repérer les termes clés} du sujet et les définir rapidement (mentalement ou sur brouillon).
\item \textbf{Dégager la tension} : quel paradoxe ou quelle opposition le sujet contient-il ? (ex. si la liberté fonde la responsabilité, peut-on être responsable sans être libre ? La responsabilité n'est-elle pas une limite à la liberté ?)
\item \textbf{Formuler la question} en une phrase interrogative claire, qui appelle une réponse nuancée (pas un simple oui/non).
\item \textbf{Annoncer le plan} : deux ou trois moments logiques (thèse / antithèse / synthèse, ou analyse progressive : concept $\rightarrow$ lien $\rightarrow$ limites).
\end{enumerate}
\textbf{Réflexe} : une bonne problématique n'appelle jamais une réponse par « oui » ou « non » seulement.
\end{methode}

\begin{exoresolu}[Problématiser : « Sommes-nous responsables de tous nos actes ? »]
\textbf{Énoncé} : Pour le sujet « Sommes-nous responsables de tous nos actes ? », dégagez la problématique et proposez un plan détaillé en trois parties.
\tcblower
\textbf{Solution détaillée.}

\textbf{1. Termes clés} : « responsables » (devoir répondre de ses actes devant la conscience, la loi, les autres), « tous nos actes » (sans exception ? y compris réflexes, actes sous contrainte, rêves, oublis ?).

\textbf{2. Tension} : d'un côté, la responsabilité suppose la liberté, et l'homme semble libre de ses choix (Sartre) ; de l'autre, certains actes sont accomplis sous contrainte physique, dans l'ignorance, la folie, l'émotion violente ou l'enfance : peut-on alors en répondre pleinement ?

\textbf{3. Problématique} : « Si la responsabilité suppose la liberté comme condition nécessaire, dans quelle mesure pouvons-nous être tenus pour responsables de \emph{tous} nos actes, y compris ceux accomplis sans pleine liberté (contrainte, ignorance, pathologie) ? »

\textbf{4. Plan annoncé} :
\begin{itemize}
\item \textbf{I. La thèse de la responsabilité universelle :} la liberté comme fondement. L'agent libre et conscient répond de ses actes (Sartre : nous sommes « condamnés à être libres » ; exemple : l'élève qui triche librement au Probatoire).
\item \textbf{II. Les limites de la responsabilité : les causes d'irresponsabilité.} La contrainte absolue, la folie (démence), l'ignorance invincible, la minorité diminuent ou suppriment la responsabilité (exemple : l'enfant de 7 ans qui casse une vitre ; le malade mental devant le juge des tutelles).
\item \textbf{III. La responsabilité à degrés : une responsabilité proportionnée.} Il faut distinguer responsabilité pleine, atténuée, partagée. La société (droit camerounais) prévoit des régimes adaptés (irresponsabilité pénale, responsabilité civile du fait d'autrui, mesures éducatives). Être responsable, ce n'est pas être responsable de tout, mais répondre de ce qu'on a librement voulu.
\end{itemize}
\textbf{Conclusion attendue} : la problématique oppose la thèse de la liberté totale (responsabilité absolue) aux limites concrètes de la liberté (irresponsabilité), ce qui autorise une réponse nuancée : responsables de nos actes \emph{libres}, pas de nos actes \emph{subis}.
\end{exoresolu}

\begin{methode}[Rédiger et justifier un paragraphe argumenté (corps de dissertation)]
Pour rédiger un paragraphe argumenté au Bac (structure IDEA : Idée - Définition - Explication/Argument - Auteur - Exemple - Amorce/Bilan) :
\begin{enumerate}
\item \textbf{Idée (Thèse)} : Affirmer clairement l'idée directrice du paragraphe (une phrase).
\item \textbf{Définition} : Préciser les termes clés mobilisés dans l'idée.
\item \textbf{Argumentation} : Développer le raisonnement logique (cause, conséquence, analogie, principe).
\item \textbf{Autorité} : Appuyer sur une référence d'auteur citée brièvement (nom + thèse), expliquée et liée à l'argument.
\item \textbf{Exemple concret} : Illustrer par un exemple de la vie courante ou du contexte camerounais (obligatoire au Bac technique).
\item \textbf{Bilan (Transition)} : Conclure le paragraphe en rappelant ce qui est établi et en ouvrant sur le suivant.
\end{enumerate}
\textbf{Formule à retenir} : Idée $\rightarrow$ Définition $\rightarrow$ Argument $\rightarrow$ Auteur $\rightarrow$ Exemple $\rightarrow$ Bilan.
\end{methode}

\begin{exoresolu}[Rédiger un paragraphe : la liberté fonde la responsabilité]
\textbf{Énoncé} : Rédigez un paragraphe argumenté montrant que la liberté est le fondement de la responsabilité.
\tcblower
\textbf{Solution détaillée.}

\textit{[Idée]} On ne peut être responsable que de ce qu'on a librement accompli : la liberté est le fondement ontologique et moral de la responsabilité. \textit{[Définition]} Être libre, c'est pouvoir choisir entre plusieurs possibles sans contrainte déterminante ; être responsable, c'est devoir répondre de ses actes devant sa conscience ou devant la loi en tant qu'auteur. \textit{[Argument]} En effet, attribuer un acte à quelqu'un, c'est supposer qu'il en est la cause par sa volonté : si l'acte était inévitable (nécessité physique, contrainte), le reprocher serait absurde, comme reprocher à la pierre de tomber (Spinoza). \textit{[Auteur]} C'est pourquoi Sartre soutient que l'homme, « condamné à être libre », est responsable de tout ce qu'il fait, car chaque acte exprime un choix, fût-ce le choix de ne pas choisir. \textit{[Exemple]} Ainsi, un élève de Yaoundé qui triche à l'examen du Probatoire technique peut être sanctionné (exclusion, zéro), car il aurait pu choisir de réviser ; en revanche, on ne sanctionne pas celui qui, bousculé violemment par un camarade, déchire malgré lui la copie de ce dernier. \textit{[Bilan]} La responsabilité ne s'attache donc qu'aux actes libres : là où la liberté disparaît, la responsabilité disparaît avec elle.
\end{exoresolu}

\begin{methode}[Traiter une question de réflexion courte (type évaluation / QCM justifié)]
Pour répondre à une question directe en 10 à 15 lignes (environ 80-100 mots) :
\begin{enumerate}
\item \textbf{Reformuler} la question en une phrase pour montrer qu'elle est comprise.
\item \textbf{Définir} la ou les notions mobilisées (liberté, responsabilité, cause d'irresponsabilité).
\item \textbf{Répondre} en deux ou trois arguments maximum, chacun illustré d'un exemple concret.
\item \textbf{Conclure} par une phrase de synthèse qui répond explicitement à la question posée.
\end{enumerate}
\textbf{Réflexe} : ne jamais réciter tout le cours ; sélectionner uniquement ce qui répond à la question.
\end{methode}

\begin{exoresolu}[Question courte : peut-on être responsable sans être libre ?]
\textbf{Énoncé} : Peut-on être responsable sans être libre ? Répondez en une dizaine de lignes en vous appuyant sur un exemple.
\tcblower
\textbf{Solution détaillée.}

La question demande si la responsabilité peut exister lorsque la liberté fait défaut. Être responsable, c'est devoir répondre de ses actes ; être libre, c'est avoir pu choisir entre plusieurs possibles. Or, on ne peut raisonnablement demander des comptes qu'à celui qui aurait pu agir autrement. Ainsi, un homme poussé violemment dans la foule à la gare de Ngaoundéré et qui renverse un passant n'est pas responsable du dommage : la contrainte physique a supprimé sa liberté de mouvement. De même, le droit camerounais (Code pénal, art. 63-65) reconnaît des causes d'irresponsabilité : la démence, la contrainte physique irrésistible, la minorité atténuent ou effacent la responsabilité pénale. En revanche, dès que le choix est possible, la responsabilité revient tout entière : l'élève qui copie librement au Baccalauréat technique répond de sa fraude. \textbf{Conclusion} : on ne peut pas être pleinement responsable sans être libre ; la responsabilité suit la liberté comme l'ombre suit le corps.
\end{exoresolu}

\begin{methode}[Construire une conclusion de dissertation]
Pour rédiger la conclusion (dernier paragraphe, pas de saut de page) :
\begin{enumerate}
\item \textbf{Répondre} explicitement à la problématique posée dans l'introduction (reprendre les termes de la question).
\item \textbf{Bilan} : résumer en une ou deux phrases le chemin parcouru (les grandes étapes du devoir : I, II, III).
\item \textbf{Nuance} : rappeler la position défendue et ses limites éventuelles (éviter le dogmatisme).
\item \textbf{Ouverture} (facultative mais valorisée) : poser une question nouvelle qui prolonge la réflexion (ex. de la responsabilité individuelle vers la responsabilité collective, environnementale, ou technologique).
\end{enumerate}
\textbf{Attention} : la conclusion ne doit contenir \textbf{aucun argument nouveau} ni citation inédite.
\end{methode}

\begin{exoresolu}[Rédiger une conclusion]
\textbf{Énoncé} : Rédigez la conclusion d'une dissertation sur le sujet : « La liberté est-elle le fondement de la responsabilité ? »
\tcblower
\textbf{Solution détaillée.}

\textit{[Réponse à la problématique]} Nous pouvons désormais répondre : oui, la liberté est bien le fondement nécessaire de la responsabilité, car on ne peut répondre que des actes qu'on a librement choisis. \textit{[Bilan]} Notre réflexion a d'abord montré que tout acte libre engage son auteur devant sa conscience et devant la loi (I) ; elle a ensuite reconnu que la contrainte, la folie ou l'ignorance suspendent ou atténuent cette responsabilité (II) ; elle a enfin établi que la responsabilité admet des degrés, proportionnés à la liberté réelle de l'agent (III). \textit{[Nuance]} Dire que la liberté fonde la responsabilité ne signifie donc pas que nous sommes responsables de tout, mais que nous répondons de ce que nous avons voulu. \textit{[Ouverture]} Reste alors une question cruciale pour notre époque : à l'heure où les décisions individuelles (pollution, usage du numérique, votes) affectent la société entière et les générations futures, notre responsabilité ne devient-elle pas aussi \emph{collective} et \emph{planétaire} ?
\end{exoresolu}

\begin{attention}[Les 5 erreurs qui coûtent des points au Bac (Philosophie Technique)]
\begin{enumerate}
\item \textbf{Confondre liberté et absence de règles / libertinage} : la liberté philosophique n'est pas « faire ce qu'on veut » sans limites ; elle s'exerce dans un cadre moral et juridique (lois, devoirs). \textbf{Ne jamais écrire} : « être libre, c'est faire n'importe quoi. »
\item \textbf{Oublier les causes d'irresponsabilité} : affirmer que l'homme est responsable de \emph{tous} ses actes sans nuance est une erreur grave. Toujours mentionner la contrainte, la folie, l'ignorance invincible, la minorité comme limites de la responsabilité (droit et morale).
\item \textbf{Ne pas distinguer les formes de responsabilité} : mélanger responsabilité morale, juridique et sociale dans le même raisonnement sans les nommer brouille la copie. \textbf{Nommer précisément} la forme dont on parle à chaque paragraphe.
\item \textbf{Citer un auteur sans l'utiliser (name-dropping)} : écrire « Sartre dit que l'homme est condamné à être libre » sans expliquer la thèse (l'existence précède l'essence, pas de nature humaine) ni la relier à l'argument ne rapporte aucun point. Toute référence doit être \textbf{brièvement expliquée et intégrée} au raisonnement.
\item \textbf{Rédiger sans exemple concret (savoir-faire exigé)} : le programme exige l'application à des situations de la vie courante. Un devoir sans exemple (chauffeur, élève, citoyen, fonctionnaire, chef d'entreprise au Cameroun) reste abstrait et perd des points. \textbf{Mais} : l'exemple doit \emph{illustrer} l'argument, pas le remplacer.
\end{enumerate}
\end{attention}

\begin{savaistu}[Repères historiques et contextuels pour l'oral / l'écrit]
\begin{itemize}
\item \textbf{Antiquité} : Épicure (clinamen, atome qui dévie) vs Stoïciens (liberté = assentiment au destin). Aristote : responsabilité = acte volontaire (non forcé, non par ignorance).
\item \textbf{Moyen Âge} : Libre arbitre théologique (Dieu juge nos choix). Responsabilité devant Dieu.
\item \textbf{Classiques} : Descartes (liberté = indifférence / spontanéité), Spinoza (liberté = nécessité comprise), Leibniz (harmonie préétablie).
\item \textbf{Lumières} : Kant (autonomie = dignité, fondement de la morale et du droit), Rousseau (liberté civile = obéissance à la loi qu'on s'est prescrite).
\item \textbf{XIXe-XXe} : Marx (liberté = reconnaissance de la nécessité historique, conditions matérielles), Freud (inconscient limite la liberté consciente), Sartre (liberté radicale, angoisse, responsabilité absolue), Ricoeur (éthique de la responsabilité, imputabilité narrative).
\item \textbf{Droit camerounais} : Constitution 1996 (Préambule : libertés fondamentales), Code pénal 2016 (irresponsabilité art. 63-65, causes d'atténuation), Code civil (responsabilité civile art. 1382-1386), Code du travail (responsabilité employeur).
\end{itemize}
\end{savaistu}

\begin{exercice}[Entraînement 1 : Analyse de situation]
M. Njoya, chef de service dans une administration publique à Bertoua, signe un marché public surfacturé sous la menace explicite d'un supérieur hiérarchique : « Si tu ne signes pas, ta mutation dans le Nord est signée demain. » M. Njoya signe. Le marché est découvert. Le tribunal administratif le poursuit.
\textbf{Question :} Analysez la responsabilité de M. Njoya (morale, juridique, sociale) en distinguant le degré de sa liberté.
\end{exercice}

\begin{exercice}[Entraînement 2 : Dissertation]
\textbf{Sujet :} « La responsabilité est-elle la limite de la liberté ? »
\textbf{Consigne :} Problématisez, proposez un plan détaillé en trois parties (avec arguments, auteurs, exemples pour chaque sous-partie), rédigez l'introduction et la conclusion.
\end{exercice}

\begin{exercice}[Entraînement 3 : Question courte]
\textbf{Sujet :} « L'ignorance excuse-t-elle toujours ? » Répondez en 12 lignes maximum en distinguant ignorance invincible et ignorance vincible, avec un exemple pour chaque.
\end{exercice}

\begin{corrige}[Exercice 1 : Analyse de situation]
\textbf{Analyse :}
\begin{enumerate}
\item \textbf{Liberté} : M. Njoya subit une \emph{contrainte morale forte} (menace de mutation punitive). Sa liberté est \textbf{réduite mais non abolie} : il pouvait refuser de signer (démissionner, saisir l'inspection, alerter la CONAC). Ce n'est pas une contrainte physique absolue.
\item \textbf{Responsabilité juridique} : En droit camerounais, la contrainte morale n'est pas une cause d'irresponsabilité absolue (art. 64 CPP : « contrainte physique ou morale irrésistible »). La mutation n'est pas une menace de mort. Il reste \textbf{responsable juridiquement} (faux en écriture publique, corruption passive), mais la contrainte est une \textbf{circonstance atténuante} (peine réduite).
\item \textbf{Responsabilité morale} : Il a choisi le mal moindre (signer vs mutation). Sa conscience peut être en paix (cas de conscience) ou troublée. Responsabilité \textbf{atténuée}.
\item \textbf{Responsabilité sociale} : En tant que chef de service, il a un devoir de probité. Sa signature engage l'État. Responsabilité \textbf{pleine} vis-à-vis de sa fonction (faute de gestion), même si sa liberté était contrainte.
\end{enumerate}
\textbf{Conclusion} : Responsabilité juridique et sociale pleine (avec circonstances atténuantes pour la peine), responsabilité morale atténuée. Le supérieur est le principal responsable (auteur moral).
\end{corrige}

\begin{corrige}[Exercice 3 : Question courte]
\textbf{Réponse :} Non, l'ignorance n'excuse pas toujours. Seule l'\textbf{ignorance invincible} (que l'agent ne \emph{peut pas} surmonter malgré sa diligence) supprime la responsabilité. \emph{Exemple :} un pharmacien vend un médicament dont un lot défectueux (indétectable) tue un patient : irresponsable. En revanche, l'\textbf{ignorance vincible} (négligence, paresse à s'informer) n'excuse pas. \emph{Exemple :} un conducteur ignore la nouvelle limitation de vitesse à 50 km/h en ville (panneaux posés, campagne médiatique) : il est responsable de son excès de vitesse. Le droit (art. 1382 Code civil) et la morale sanctionnent la négligence.
\end{corrige}

\newpage

\coursec{Exercices}

\courssub{Je vérifie mes connaissances}

\begin{exercice}[QCM : les concepts fondamentaux]
Pour chaque question, une seule réponse est correcte. Recopie la lettre correspondante.
\begin{enumerate}
\item La \cle{responsabilité} se définit comme :
\begin{enumerate}
\item[a)] le pouvoir de faire ce qu'on veut sans limites ;
\item[b)] l'obligation de répondre de ses actes et d'en assumer les conséquences ;
\item[c)] la soumission aux lois de l'État.
\end{enumerate}
\item Le \cle{libre arbitre} désigne :
\begin{enumerate}
\item[a)] la capacité de choisir entre plusieurs possibilités sans être totalement déterminé ;
\item[b)] la liberté de penser sans agir ;
\item[c)] l'absence de toute loi morale.
\end{enumerate}
\item Selon Sartre, l'homme est « condamné à être libre » parce que :
\begin{enumerate}
\item[a)] il est puni par la société ;
\item[b)] il n'a pas choisi d'exister mais doit choisir ce qu'il fait de son existence ;
\item[c)] il est prisonnier de ses instincts.
\end{enumerate}
\item L'\cle{imputation} consiste à :
\begin{enumerate}
\item[a)] attribuer un acte à un auteur capable d'en répondre ;
\item[b)] calculer une amende ;
\item[c)] excuser un coupable.
\end{enumerate}
\end{enumerate}
\end{exercice}

\begin{exercice}[Vrai ou faux, en justifiant]
Indique si chaque affirmation est vraie ou fausse, puis justifie ta réponse en deux ou trois lignes.
\begin{enumerate}
\item « On ne peut être responsable que de ce qu'on a fait librement. »
\item « Le déterminisme affirme que nos actes n'ont aucune cause. »
\item « La responsabilité morale se confond toujours avec la responsabilité pénale. »
\item « Un mineur ou un fou peut être déclaré non responsable de ses actes. »
\end{enumerate}
\end{exercice}

\begin{exercice}[Définitions et emploi]
\begin{enumerate}
\item Définis en trois lignes maximum chacune des notions suivantes : \cle{liberté}, \cle{libre arbitre}, \cle{responsabilité}, \cle{imputation}.
\item Utilise ensuite chacune de ces quatre notions dans une phrase argumentée portant sur une situation de la vie courante (école, famille, circulation, travail).
\end{enumerate}
\end{exercice}

\begin{exercice}[Reconnaître les formes de responsabilité]
Pour chacune des situations suivantes, indique s'il s'agit d'une responsabilité \cle{civile}, \cle{pénale} ou \cle{morale} (plusieurs réponses possibles par situation) :
\begin{enumerate}
\item Un élève de Terminale technique casse la vitre de l'atelier et doit la rembourser.
\item Un chauffeur de taxi-brousse, ivre, provoque un accident mortel sur l'axe Yaoundé--Bafoussam.
\item Un père de famille abandonne ses enfants sans ressources, sans enfreindre aucune loi écrite.
\item Un commerçant de Mokolo vend des produits périmés et est condamné à une peine de prison.
\end{enumerate}
\end{exercice}

\courssub{Je m'entraîne}

\begin{exercice}[Liberté et contraintes]
Un élève dit : « Je ne suis pas libre, car mes parents m'obligent à aller à l'école et le règlement m'interdit beaucoup de choses. »
\begin{enumerate}
\item Distingue la \cle{liberté de faire} (liberté extérieure) et la \cle{liberté de vouloir} (liberté intérieure).
\item Montre, à partir de deux exemples scolaires précis, que l'existence de règles ne supprime pas nécessairement la liberté.
\item Rédige une réponse argumentée de dix à quinze lignes à cet élève.
\end{enumerate}
\end{exercice}

\begin{exercice}[Le déterminisme en question]
On appelle \cle{déterminisme} la thèse selon laquelle tout événement, y compris nos décisions, est l'effet nécessaire de causes antérieures (hérédité, milieu social, éducation).
\begin{enumerate}
\item Explique en quelques lignes pourquoi le déterminisme semble menacer la responsabilité.
\item À l'aide d'un exemple camerounais (par exemple un jeune délinquant issu d'un quartier défavorisé de Douala), montre qu'on peut à la fois reconnaître l'influence des causes et maintenir une part de responsabilité.
\item Rédige une conclusion personnelle justifiée.
\end{enumerate}
\end{exercice}

\begin{exercice}[Étude de cas : l'accident de la circulation]
À Douala, un conducteur de moto-taxi (\emph{bendskin}), pressé par un client, grille un feu rouge et renverse un piéton. Le piéton est blessé ; la moto, prêtée par un ami, n'est pas assurée.
\begin{enumerate}
\item Identifie les personnes impliquées et les actes commis.
\item Établis, en justifiant, la responsabilité \cle{civile}, la responsabilité \cle{pénale} et la responsabilité \cle{morale} du conducteur.
\item Le client pressé et l'ami propriétaire de la moto sont-ils responsables ? Argumente.
\item Quelles conditions (conscience, liberté, connaissance de la loi) rendent l'\cle{imputation} légitime dans ce cas ?
\end{enumerate}
\end{exercice}

\begin{exercice}[Débat écrit type Probatoire]
Sujet : « La responsabilité des gouvernants est-elle plus grande que celle des citoyens ? »
\begin{enumerate}
\item Rédige un paragraphe (huit à dix lignes) défendant la thèse : oui, car le gouvernant a plus de pouvoir, donc plus de responsabilité.
\item Rédige un paragraphe défendant la thèse opposée : non, car chaque citoyen est responsable de ses actes dans sa propre sphère.
\item Conclus en prenant position de manière justifiée, en t'appuyant sur un exemple de la vie publique camerounaise.
\end{enumerate}
\end{exercice}

\courssub{Je me prépare au Bac}

\begin{methode}[Rédiger une dissertation philosophique au Probatoire]
\begin{enumerate}
\item \textbf{Analyse du sujet :} Identifie les termes clés, les oppose ou les relie, dégage la tension (problème). Évite les définitions dictionnaires ; propose des définitions \emph{philosophiques} (enjeu conceptuel).
\item \textbf{Problématisation :} Transforme le sujet en question explicite : « Dans quelle mesure... ? », « Faut-il... ? » Annonce le plan (dialectique : Thèse / Antithèse / Synthèse).
\item \textbf{Développement :} 3 parties équilibrées. Chaque partie = 1 idée directrice + 1 argumentation (concepts, auteurs, exemples) + 1 transition. Cite les auteurs au programme (Sartre, Kant, Spinoza, etc.).
\item \textbf{Conclusion :} Réponse synthétique à la problématique. Ouverture sur une question voisine ou une application concrète.
\end{enumerate}
\end{methode}

\begin{exercice}[Dissertation : liberté et responsabilité]
\textbf{Sujet :} « Faut-il être libre pour être responsable ? »

Travail demandé :
\begin{enumerate}
\item Analyse le sujet : dégage le problème posé par la formulation « faut-il » et définis les termes \cle{liberté} et \cle{responsabilité} dans l'introduction.
\item Construis un plan dialectique en trois parties : (I) la responsabilité suppose la liberté (on ne répond que de ce qu'on a choisi) ; (II) mais nos choix sont influencés par des déterminismes (milieu, éducation, passions) ; (III) synthèse : la liberté n'est jamais absolue, mais une liberté partielle suffit à fonder la responsabilité.
\item Rédige la dissertation complète (introduction, développement, conclusion) en illustrant chaque partie d'exemples camerounais et en citant au moins deux auteurs étudiés (Sartre, Kant, Spinoza).
\end{enumerate}
\end{exercice}

\begin{methode}[Commentaire de texte philosophique (Sartre)]
\begin{enumerate}
\item \textbf{Thèse :} Résume l'idée centrale en une phrase affirmative.
\item \textbf{Paradoxe / Concepts clés :} Explique les termes forts (« condamné », « liberté », « responsabilité totale »).
\item \textbf{Argumentation :} Reconstitue le cheminement logique : Premisse 1 (facticité) $\rightarrow$ Premisse 2 (liberté ontologique) $\rightarrow$ Conséquence (responsabilité universelle).
\item \textbf{Portée :} Explique la portée éthique (universalisation du choix).
\item \textbf{Discussion critique :} Évalue la thèse (forces/limites) avec un contre-exemple concret.
\end{enumerate}
\end{methode}

\begin{exercice}[Commentaire de texte : Sartre]
\textbf{Texte :} « L'homme est condamné à être libre. Condamné, parce qu'il ne s'est pas créé lui-même ; et par ailleurs libre, parce qu'une fois jeté dans le monde, il est responsable de tout ce qu'il fait. [...] L'homme est responsable de ce qu'il est. [...] En choisissant, je choisis pour tous les hommes. » (J.-P. Sartre, \emph{L'Existentialisme est un humanisme}, 1946.)

Travail demandé :
\begin{enumerate}
\item Dégage la \cle{thèse} du texte et reformule-la en une phrase.
\item Explique le paradoxe de l'expression « condamné à être libre ».
\item Reconstitue l'argumentation : pourquoi, selon Sartre, la liberté entraîne-t-elle une responsabilité totale ?
\item Quelle est la portée morale de la formule « en choisissant, je choisis pour tous les hommes » ?
\item Discussion : cette conception d'une responsabilité sans excuse est-elle acceptable ? Appuie ta réponse sur un exemple concret (par exemple un élève qui triche à l'examen du Probatoire).
\end{enumerate}
\end{exercice}

\begin{exercice}[Question de réflexion argumentée : déterminisme et responsabilité]
\textbf{Sujet :} « Le déterminisme supprime-t-il la responsabilité ? »

Travail demandé :
\begin{enumerate}
\item Définis le \cle{déterminisme} et explique en quoi il semble exclure toute responsabilité (si tout acte a des causes nécessaires, personne ne « choisit » vraiment).
\item Première partie (thèse) : montre, avec Spinoza, que nous nous croyons libres parce que nous ignorons les causes qui nous déterminent ; illustre par un exemple camerounais (consommation d'alcool héritée du milieu familial, violence apprise dans l'enfance).
\item Deuxième partie (antithèse et synthèse) : montre que la société et le droit camerounais continuent d'imputer les actes à leurs auteurs, parce qu'une conscience, même influencée, conserve un pouvoir de choix ; illustre par le cas d'un accusé devant un tribunal de première instance.
\item Rédige une conclusion personnelle et justifiée d'une dizaine de lignes.
\end{enumerate}
\end{exercice}

\newpage

\coursec{Corrigés des exercices}

\coursec{Introduction : problématisation et définitions}

\begin{definition}[Liberté]
La \cle{liberté} est le pouvoir de se déterminer soi-même à agir, c'est-à-dire la capacité de choisir entre plusieurs possibles sans contrainte extérieure absolue ni nécessité intérieure fatale. Elle se décline en \emph{liberté d'agir} (absence d'obstacles physiques ou juridiques) et \emph{liberté de vouloir} (autonomie de la volonté, capacité de choisir ses fins).
\end{definition}

\begin{definition}[Responsabilité]
La \cle{responsabilité} est l'obligation pour un agent moral ou juridique de répondre de ses actes, d'en assumer les conséquences et de subir, le cas échéant, la sanction ou l'obligation de réparation. Elle suppose l'imputabilité : l'acte doit pouvoir être rattaché à un sujet conscient et libre.
\end{definition}

\begin{aretenir}[Problématique centrale]
Le couple liberté-responsabilité structure l'anthropologie philosophique et le droit. \textbf{Problème} : si la liberté est la condition de la responsabilité, que devient cette dernière face au déterminisme (hérédité, milieu, inconscient) ? Une liberté partielle suffit-elle à fonder une responsabilité pleine ? Le programme de Terminale technique exige de distinguer les formes de liberté, les formes de responsabilité, et d'établir le lien de fondement entre les deux, en l'appliquant à des situations concrètes camerounaises.
\end{aretenir}

\coursec{Les formes de la liberté}

\courssub{Liberté d'agir (ou liberté extérieure) et liberté de vouloir (ou liberté intérieure)}

\begin{exemplebox}[Distinction fondamentale]
La \cle{liberté d'agir} est le pouvoir d'exécuter physiquement ses décisions sans entrave matérielle (chaînes, enfermement, paralysie) ni contrainte juridique absolue. Elle est variable, soumise aux lois et aux circonstances. La \cle{liberté de vouloir} est le pouvoir de choisir ses intentions, ses jugements, ses consentements. Elle relève du for intérieur : même enchaîné, l'esclave stoïcien \textbf{Épictète} reste libre d'assentir ou non à sa condition. \emph{« Il n'y a qu'un chemin qui mène au bonheur : cesser de se soucier de ce qui est hors de notre pouvoir. »} (Épictète, \emph{Manuel}, I).
\end{exemplebox}

\courssub{Autonomie et hétéronomie (Kant)}

\begin{definition}[Autonomie / Hétéronomie]
L'\cle{autonomie} (loi que l'on se donne à soi-même) est la marque de la liberté morale : la volonté rationnelle se soumet à la loi morale par devoir, non par inclination. L'\cle{hétéronomie} (loi venue de l'extérieur : désirs, coutumes, autorité) est l'état de la volonté déterminée par des causes sensibles. Pour \textbf{Kant}, \emph{« Autonomie est le fondement de la dignité de la nature humaine et de toute nature raisonnable »} (\emph{Fondements de la métaphysique des mœurs}). Seule l'autonomie fonde une responsabilité véritable.
\end{definition}

\courssub{Liberté et déterminisme}

\begin{propriete}[Déterminisme]
Le \cle{déterminisme} affirme que tout événement, y compris les décisions humaines, est l'effet nécessaire de causes antérieures (physiques, biologiques, psychologiques, sociales). \textbf{Spinoza} (\emph{Éthique}, III, scolie de la prop. 2) montre que les hommes se croient libres par ignorance des causes qui les déterminent : \emph{« Les hommes se croient libres parce qu'ils ont conscience de leurs actions et ignorent les causes par lesquelles ils sont déterminés. »} Cette thèse menace l'imputabilité.
\end{propriete}

\begin{attention}[Piège à éviter]
Ne pas confondre \emph{déterminisme} (nécessité causale universelle) et \emph{fatalisme} (destin inéluctable indépendamment des causes). Le déterminisme n'exclut pas l'action, il en explique la genèse. La responsabilité compatibiliste (voir §4) repose sur cette nuance.
\end{attention}

\coursec{Les formes de la responsabilité}

\courssub{Responsabilité morale}
Elle engage la conscience devant soi-même et devant autrui. Elle suppose la \cle{liberté de vouloir} et la \cle{conscience morale} (discernement entre le bien et le mal). Sanction : remords, blâme, estime ou mépris social. Exemple : un élève qui triche au Probatoire porte une responsabilité morale même s'il n'est pas pris.

\courssub{Responsabilité juridique (Droit camerounais)}

\begin{center}
\begin{tabularx}{\linewidth}{|l|X|X|}
\hline
\rowcolor{popblueL} \textbf{Type} & \textbf{Fondement \& Conditions} & \textbf{Sanction / Réparation} \\
\hline
\cle{Pénale} & Infraction à la loi pénale (Code pénal). Auteur matériel ou complice. \textbf{Conditions :} conscience, liberté, imputabilité (majorité, discernement). & Peines : emprisonnement, amende, peine de mort (théorique), travaux d'intérêt général. But : défense sociale, réinsertion. \\
\hline
\cle{Civile} & Faute (art. 1382 Code civil : \emph{« Tout fait quelconque de l'homme, qui cause à autrui un dommage, oblige celui par la faute duquel il est arrivé à le réparer »}), dommage, lien de causalité. & Réparation intégrale (dommages-intérêts) : restitution, réparation en nature, compensation pécuniaire. \\
\hline
\cle{Administrative} & Faute de service public (dysfonctionnement de l'administration). & Réparation par l'État (Tribunal administratif). \\
\hline
\end{tabularx}
\end{center}

\courssub{Responsabilité collective et politique}
Responsabilité des gouvernants (reddition des comptes, gestion des deniers publics), responsabilité de l'État (protection des droits humains), responsabilité solidaire (ex. : dette écologique). Au Cameroun, la \cle{Constitution} (préambule) et la loi n°2018/011 sur le code de transparence et de bonne gouvernance encadrent la responsabilité des gestionnaires publics.

\coursec{La liberté, fondement de la responsabilité}

\courssub{Le lien logique : « Tu dois » implique « Tu peux » (Kant)}

\begin{propriete}[Condition de possibilité]
L'imputation (rattacher l'acte à l'agent) n'est légitime que si l'agent \emph{aurait pu faire autrement}. \textbf{Kant} : \emph{« Le devoir commande... et cette commande suppose la liberté. »} (\emph{Critique de la raison pratique}). Le droit camerounais l'applique : irresponsabilité pénale pour trouble mental abolissant le discernement (art. 72 Code pénal), responsabilité atténuée pour le mineur (art. 73).
\end{propriete}

\courssub{L'objection déterministe (Spinoza) et la réponse compatibiliste}

\begin{experience}[Débat guidé : Le jeune de New Bell]
\textbf{Situation :} Un jeune de New Bell (Douala), élevé dans la violence et la pauvreté, commet un vol à l'arraché.
\textbf{Analyse déterministe (Spinoza) :} Le milieu a déterminé ses désirs et ses habitudes. Il n'a pas « choisi » d'être né là. Punir serait punir une avalanche d'être tombée.
\textbf{Réponse compatibiliste (Droit \& Morale courante) :} D'autres jeunes du même quartier choisissent l'apprentissage ou l'école. Le milieu \emph{incline} sans \emph{nécéssiter} absolument. Le tribunal impute l'acte (liberté réelle, quoique située) et module la peine (circonstances atténuantes).
\textbf{Conclusion :} La liberté n'a pas besoin d'être absolue (indéterminisme métaphysique) pour fonder la responsabilité ; une \cle{liberté partielle, située, conquise} suffit.
\end{experience}

\courssub{Sartre : la responsabilité radicale}

\begin{aretenir}[Sartre, \emph{L'Existentialisme est un humanisme}]
\emph{« L'homme est condamné à être libre. »} \textbf{Facticité} (je n'ai pas choisi ma naissance, mon corps, mon époque) + \textbf{Liberté ontologique} (aucune essence ne me précède, je me fais en agissant) = \textbf{Responsabilité totale}. \emph{« En choisissant, je choisis pour tous les hommes. »} Mon choix érige une valeur universelle. Pas d'alibi possible (milieu, passions, ordres). \emph{Force :} dignité, lutte contre la mauvaise foi. \emph{Limite :} ignore les contraintes réelles (maladie mentale, contrainte physique irrésistible) que le droit prend en compte.
\end{aretenir}

\coursec{Applications concrètes au Cameroun}

\begin{corrige}[Exercice 5 — Liberté et contraintes]
\begin{enumerate}
\item \textbf{Distinction.} La \cle{liberté de faire} (ou liberté extérieure) est le pouvoir d'accomplir physiquement ses actes sans obstacle matériel ni coercition : elle peut être limitée par les lois, les règlements, l'enfermement. La \cle{liberté de vouloir} (ou liberté intérieure) est le pouvoir de choisir ses intentions, ses jugements, ses consentements : elle relève du for intérieur et demeure même sous la contrainte, comme le montre l'esclave stoïcien \textbf{Épictète}, libre dans sa pensée malgré les chaînes.
\item \textbf{Exemples scolaires.} (a) Le règlement interdit de sortir pendant les cours : il limite ma liberté de faire, mais ne supprime pas ma liberté de vouloir — je peux choisir librement d'approuver cette règle parce qu'elle protège mon travail, ou de la critiquer intérieurement. (b) L'obligation de passer le Probatoire ne m'empêche pas de choisir librement ma méthode de travail, mes horaires de révision, ni même d'assumer le refus de travailler (au prix d'en assumer les conséquences). La règle encadre l'exercice de la liberté ; elle ne l'abolit pas. \textbf{Kant} dirait même que la loi morale, que je me donne à moi-même (autonomie), est la condition d'une liberté véritable et non son ennemie.
\item \textbf{Réponse argumentée (modèle).} Ton raisonnement confond deux choses : subir des règles et être privé de liberté. Certes, tes parents et le règlement limitent ta liberté de faire : tu ne peux pas tout faire, n'importe quand. Mais aucune contrainte extérieure ne peut te forcer à vouloir : c'est toi qui choisis d'obéir par conviction, par intérêt ou par crainte, et ce choix t'appartient. D'ailleurs, les règles scolaires ne visent pas à t'asservir mais à rendre possible un bien plus grand : l'instruction, qui élargira demain tes possibilités réelles. Un élève qui comprend cela cesse de vivre l'école comme une prison : il transforme l'obligation subie en projet librement assumé. La vraie question n'est donc pas « suis-je libre ? » mais « que fais-je de la liberté qui me reste ? » — et celle-là, personne ne peut te la retirer.
\end{enumerate}
\textbf{Points de vigilance :} ne pas nier la réalité des contraintes (ce serait naïf) ; distinguer clairement les deux sens du mot « liberté » ; conclure sur la responsabilité (assumer ses choix dans le cadre donné).
\end{corrige}

\begin{corrige}[Exercice 6 — Le déterminisme en question]
\begin{enumerate}
\item \textbf{Pourquoi le déterminisme menace la responsabilité.} Si tout acte est l'effet nécessaire de causes antérieures (hérédité, milieu, éducation), alors l'agent n'a pas pu faire autrement : or on ne peut être tenu responsable que de ce qu'on aurait pu ne pas faire. Le déterminisme semble donc dissoudre l'imputabilité : punir reviendrait à punir une avalanche d'être tombée.
\item \textbf{Exemple camerounais.} Un jeune de New Bell (Douala), élevé dans la violence et la pauvreté, commet un vol à l'arraché. Reconnaître les causes, c'est comprendre : son milieu a pesé sur ses habitudes, ses tentations, son absence de perspectives — et le juge peut en tenir compte (circonstances atténuantes). Mais maintenir une part de responsabilité, c'est constater que d'autres jeunes du même quartier, dans les mêmes conditions, choisissent l'apprentissage d'un métier ou l'école : le milieu incline sans nécessiter absolument. Le tribunal de première instance de Douala imputera donc l'acte au jeune homme, tout en modulant la peine. La société qui excuserait totalement le délinquant le traiterait comme une chose, non comme un homme — ce serait le mépriser davantage encore.
\item \textbf{Conclusion personnelle (modèle).} Le déterminisme décrit des influences réelles que l'honnêteté interdit de nier ; mais l'influence n'est pas la nécessité absolue. Entre la cause et l'acte subsiste un espace de délibération, si étroit soit-il, où le sujet peut dire oui ou non. C'est cet espace — une liberté partielle, située, conquise — qui suffit à fonder la responsabilité. Reconnaître les déterminismes doit donc nous conduire non à excuser à tout prix, mais à éduquer, prévenir et punir avec mesure : la justice camerounaise fait exactement cela lorsqu'elle distingue les circonstances atténuantes de l'irresponsabilité totale.
\end{enumerate}
\textbf{Points de vigilance :} bien définir le déterminisme (nécessité des causes, non absence de causes) ; éviter deux excès symétriques — la liberté absolue qui ignore le milieu, et l'excuse totale qui nie le sujet.
\end{corrige}

\begin{corrige}[Exercice 7 — Étude de cas : l'accident de la circulation]
\begin{enumerate}
\item \textbf{Personnes et actes.} Personnes : le conducteur de moto-taxi, le piéton blessé (victime), le client pressé, l'ami propriétaire de la moto. Actes : franchissement d'un feu rouge (infraction au code de la route), blessures causées au piéton, circulation d'un véhicule non assuré.
\item \textbf{Responsabilités du conducteur.}
\begin{itemize}
\item \textbf{Pénale :} il commet une infraction (violation du feu rouge) et une blessure involontaire causée par imprudence ; le Code pénal camerounais punit les atteintes corporelles involontaires résultant d'une faute d'imprudence. Il encourt une peine d'emprisonnement et/ou d'amende.
\item \textbf{Civile :} il doit réparer intégralement le dommage causé à autrui (faute + dommage + lien de causalité) : frais médicaux du piéton, indemnisation du préjudice. L'absence d'assurance ne l'exonère pas : il paiera de son patrimoine.
\item \textbf{Morale :} devant sa conscience et la communauté, il a mis la vie d'autrui en danger pour un gain de temps ; il doit reconnaissance, réparation et repentir.
\end{itemize}
\item \textbf{Le client et l'ami.} Le client a exercé une pression, mais une pression verbale n'est pas une contrainte absolue : le conducteur gardait le libre arbitre de refuser de griller le feu. Le client porte donc une \emph{responsabilité morale} (incitation à l'imprudence), mais pas de responsabilité pénale directe, sauf si sa pression constituait une contrainte irrésistible ou une complicité. L'ami propriétaire a prêté un véhicule non assuré : il engage sa \emph{responsabilité civile} (le défaut d'assurance est une faute qui l'expose à garantir l'indemnisation) et éventuellement pénale si la loi sanctionne la mise en circulation d'un véhicule non assuré ; il n'est en revanche pas auteur de l'infraction de conduite.
\item \textbf{Conditions de l'imputation.} L'imputation du conducteur est légitime car : (a) \emph{conscience} — il savait ce qu'il faisait (griller un feu) ; (b) \emph{liberté} — aucune contrainte physique irrésistible ne l'y forçait, la pression du client laissait subsister un choix ; (c) \emph{connaissance de la loi} — nul n'est censé ignorer la loi, et tout conducteur connaît la signification d'un feu rouge. Ces trois conditions réunies, l'acte peut lui être attribué comme à une cause libre.
\end{enumerate}
\textbf{Points de vigilance :} articuler les trois éléments de la responsabilité civile (faute, dommage, causalité) ; distinguer soigneusement l'auteur principal, l'instigateur moral et le garant civil ; rappeler que l'absence d'assurance aggrave la situation au lieu d'excuser.
\end{corrige}

\begin{corrige}[Exercice 8 — Débat écrit type Probatoire]
\begin{enumerate}
\item \textbf{Thèse (oui).} Le gouvernant détient un pouvoir dont les effets portent sur des millions de personnes : une décision budgétaire, un décret, un choix d'infrastructure engagent la vie de tout un pays. Or la responsabilité croît avec le pouvoir : celui qui peut davantage doit répondre de davantage. Le ministre qui détourne les fonds destinés à un lycée technique prive des milliers d'élèves de formation ; sa faute est incommensurable à celle d'un citoyen ordinaire. De plus, le gouvernant a librement sollicité ou accepté sa charge : il a donc contracté, devant la nation, une obligation renforcée de rendre compte. Enfin, il dispose de moyens d'information et de conseillers qui rendent ses choix plus éclairés, donc plus imputables. À pouvoir accru, responsabilité accrue : telle est la logique même de l'imputation.
\item \textbf{Antithèse (non).} Pourtant, la responsabilité morale se mesure d'abord à la liberté de l'agent, et celle-ci est égale en tous : chaque citoyen choisit ses actes dans sa sphère et doit en répondre pleinement. Le commerçant qui vend des produits périmés, le chauffeur ivre, l'électeur qui vend sa voix : aucun ne peut se réfugier derrière la faute des gouvernants. Dire que le citoyen est « moins responsable » serait l'infantiliser et lui refuser la dignité d'agent moral. De plus, les gouvernants sortent du peuple : une cité corrompue produit des dirigeants corrompus. La responsabilité est donc égale en nature, même si elle diffère en étendue.
\item \textbf{Conclusion (modèle).} Il faut distinguer : en \emph{nature}, la responsabilité est égale, car elle repose sur la liberté, identique en chaque homme ; en \emph{étendue}, elle est proportionnelle au pouvoir, car les conséquences des actes du gouvernant sont plus vastes. Lorsqu'au Cameroun l'opinion réclame la reddition des comptes (\emph{rendre compte} de la gestion publique), elle ne dispense pas pour autant chaque citoyen de payer ses impôts, de respecter le code de la route ou de voter selon sa conscience. La cité juste est celle où chacun répond de ses actes à la mesure de sa liberté et de sa fonction : le gouvernant devant la nation, le citoyen devant sa conscience et la loi.
\end{enumerate}
\textbf{Barème indicatif (sur 10) :} thèse argumentée (3 pts), antithèse argumentée (3 pts), conclusion personnelle avec distinction conceptuelle et exemple camerounais (3 pts), clarté et langue (1 pt).
\textbf{Points de vigilance :} ne pas opposer les thèses de façon caricaturale ; fonder l'égalité sur la liberté et la différence sur l'étendue du pouvoir ; exemple précis et non polémique.
\end{corrige}

\coursec{Méthode : rédiger et justifier une démarche philosophique}

\begin{methode}[Rédiger une dissertation philosophique (Type Probatoire / Bac Tech)]
\begin{enumerate}
\item \textbf{Analyser le sujet} (15 min) : définir les termes clés, identifier la tension (problème), formuler la problématique (question précise), annoncer le plan dialectique (Thèse / Antithèse / Synthèse).
\item \textbf{Construire le plan détaillé} (15 min) : pour chaque partie, 2 à 3 arguments, 1 auteur/référence précis, 1 exemple concret (idéalement camerounais). Vérifier la progression logique.
\item \textbf{Rédiger l'introduction} (10 min) : accroche (exemple, citation, fait d'actualité) $\rightarrow$ définitions $\rightarrow$ problématique $\rightarrow$ annonce du plan.
\item \textbf{Rédiger le développement} (1h15) : un paragraphe = un argument. Structure : \textbf{Thèse partielle} $\rightarrow$ \textbf{Argumentation} (raisonnement, citation, concept) $\rightarrow$ \textbf{Exemple/Illustration} $\rightarrow$ \textbf{Mini-transition}. Respecter l'équilibre des parties.
\item \textbf{Rédiger la conclusion} (10 min) : réponse synthétique à la problématique $\rightarrow$ ouverture (question nouvelle, prolongement éthique/politique).
\item \textbf{Relire} (10 min) : orthographe, syntaxe, cohérence, respect du sujet.
\end{enumerate}
\end{methode}

\begin{methode}[Commentaire de texte philosophique]
\begin{enumerate}
\item \textbf{Lecture active} (20 min) : identifier la thèse principale, le vocabulaire technique, la structure du raisonnement (prémisses $\rightarrow$ conclusion), le contexte auteur/époque.
\item \textbf{Introduction} : présenter l'auteur, l'œuvre, le thème, la thèse, le plan du commentaire (suivre le texte ou thématique).
\item \textbf{Explication linéaire ou thématique} : expliciter \emph{pas à pas} le sens des phrases, définir les concepts, reconstituer l'argumentation, montrer les enjeux.
\item \textbf{Discussion critique} : évaluer la force de la thèse (cohérence, portée) ET ses limites (objections, contre-exemples, distinctions nécessaires). Ne pas se contenter d'approuver.
\item \textbf{Conclusion} : résumer l'apport du texte, situer sa portée philosophique.
\end{enumerate}
\end{methode}

\begin{corrige}[Exercice 9 — Dissertation : « Faut-il être libre pour être responsable ? »]
\textbf{1. Analyse du sujet.} Le verbe « faut-il » interroge sur une \emph{condition nécessaire} : la liberté est-elle requise pour que la responsabilité soit légitime ? Le problème : si oui, tout déterminisme (milieu, hérédité, passions) menace de dissoudre la responsabilité ; si non, on pourrait punir des actes que l'agent n'a pas choisis, ce qui semble injuste. Définitions : la \cle{liberté} est le pouvoir de se déterminer soi-même à agir ; la \cle{responsabilité} est l'obligation de répondre de ses actes devant la conscience, la société ou la loi.

\textbf{2. Plan dialectique.} I. La responsabilité suppose la liberté (thèse). II. Objection : nos choix sont déterminés (antithèse). III. Synthèse : une liberté partielle, située, suffit à fonder la responsabilité.

\textbf{3. Dissertation rédigée (modèle).}

\emph{Introduction.} Un tribunal camerounais qui juge un accusé présuppose toujours une chose : que cet homme aurait pu agir autrement. On ne reproche à personne d'être tombé quand on l'a poussé. La responsabilité — obligation de répondre de ses actes — semble donc inséparable de la liberté, pouvoir de se déterminer soi-même. Mais nos actes ne sont-ils pas, comme le veut le déterminisme, l'effet nécessaire de causes qui nous échappent ? Faut-il alors être libre pour être responsable, et de quelle liberté s'agit-il ?

\emph{I. La responsabilité suppose la liberté.} L'imputation — rattacher un acte à son auteur — n'est légitime que si l'agent était conscient et libre. C'est pourquoi le droit camerounais déclare irresponsable le fou au moment des faits et atténue la responsabilité du mineur : là où le discernement ou la liberté manquent, le blâme perd son sens. \textbf{Kant} le formule rigoureusement : le « tu dois » moral implique un « tu peux » ; on ne peut exiger d'un être que ce dont il est capable. Et \textbf{Sartre} radicalise le lien : l'homme, « condamné à être libre », est « responsable de tout ce qu'il fait », car aucune essence ne décide à sa place. L'élève qui triche au Probatoire ne peut invoquer ni la tentation ni l'exemple des autres : c'est lui qui a choisi, et c'est à lui que la sanction s'adresse.

\emph{II. Mais nos choix sont déterminés.} Cette belle construction se heurte au déterminisme. \textbf{Spinoza} affirme que les hommes se croient libres parce qu'ils ont conscience de leurs désirs tout en ignorant les causes qui les déterminent : comme la pierre lancée qui se croirait libre de sa chute. Le jeune délinquant de Douala n'a pas choisi son quartier, sa famille désunie, la violence qui fut son école ; l'alcoolique n'a pas choisi l'hérédité ni le milieu où l'on boit dès l'adolescence. Si tout acte a des causes nécessaires, « avoir pu faire autrement » serait une illusion, et punir ressemblerait à punir la pluie de tomber. La responsabilité semble alors se dissoudre dans l'enchaînement universel des causes.

\emph{III. Synthèse : une liberté partielle suffit.} Il faut pourtant distinguer influence et nécessité absolue. Les causes pèsent sur nos choix sans les supprimer : du même quartier défavorisé sortent aussi des apprentis, des commerçants honnêtes, des bacheliers. Entre la cause et l'acte demeure un espace de délibération, si étroit soit-il, où le sujet peut résister. C'est cette liberté partielle et située — non une liberté absolue hors du monde — que le droit et la morale présupposent raisonnablement. Le juge camerounais le reconnaît d'ailleurs en pratique : il impute l'acte à l'accusé, mais module la peine selon les circonstances atténuantes. Connaître les déterminismes ne doit pas excuser à tout prix ; il doit conduire à éduquer, prévenir et juger avec mesure. La liberté n'a pas besoin d'être totale pour que la responsabilité soit réelle : il suffit qu'elle soit possible.

\emph{Conclusion.} Il faut donc être libre pour être responsable — mais d'une liberté qui n'est jamais absolue : un pouvoir réel, quoique limité, de choisir au sein des conditions qui nous sont faites. La responsabilité n'exige pas que nous soyons sans causes ; elle exige que nous ne soyons pas que des effets. Reste alors une question pratique : comment l'éducation peut-elle élargir, chez chaque jeune, cet espace de liberté où sa responsabilité prend racine ?

\textbf{Barème indicatif (sur 20) :} introduction (analyse des termes, problème, annonce du plan : 4 pts) ; développement (3 parties équilibrées, arguments, auteurs cités correctement, exemples camerounais : 10 pts) ; conclusion (réponse à la problématique, ouverture : 3 pts) ; langue, organisation, orthographe (3 pts).
\textbf{Points de vigilance :} citer Spinoza et Sartre (ou Kant) avec exactitude ; ne pas plaider la liberté absolue (hors sujet et intenable) ; la synthèse doit être une position, non un compromis mou ; chaque partie doit contenir au moins un exemple concret.
\end{corrige}

\begin{corrige}[Exercice 10 — Commentaire de texte : Sartre]
\begin{enumerate}
\item \textbf{Thèse.} L'existence humaine est une liberté sans excuse : n'ayant pas choisi d'exister, l'homme est néanmoins l'auteur total de tout ce qu'il fait et de ce qu'il est, et chacun de ses choix engage l'humanité entière.
\item \textbf{Le paradoxe « condamné à être libre ».} Une condamnation est une peine subie ; la liberté semble être le contraire d'une peine. Le paradoxe se dissout si l'on distingue deux moments : \emph{condamné}, parce que je n'ai pas choisi ma naissance, mon corps, mon époque (ce que Sartre appelle la \emph{facticité}) ; \emph{libre}, parce que, dès que j'existe, aucune nature humaine préétablie ne décide à ma place — je suis obligé de choisir, et même refuser de choisir est encore un choix. La liberté n'est donc pas un don agréable mais un fardeau inévitable : on ne peut pas ne pas être libre.
\item \textbf{Argumentation.} Le raisonnement se reconstitue en trois temps : (a) \emph{premisse 1} — l'homme ne s'est pas créé lui-même : il est « jeté » dans le monde sans l'avoir voulu ; (b) \emph{premisse 2} — mais il n'existe aucune essence humaine, aucun Dieu ni nature qui fixerait d'avance ce qu'il doit être : l'existence précède l'essence, donc l'homme se fait en agissant ; (c) \emph{conséquence} — puisque rien ne le détermine, il est la cause entière de ses actes : sa responsabilité est totale, sans alibi possible (ni les passions, ni le milieu, ni les ordres reçus). Liberté ontologique et responsabilité universelle sont les deux faces d'une même réalité.
\item \textbf{Portée morale de « en choisissant, je choisis pour tous les hommes ».} Quand je choisis, j'affirme par mon acte que telle conduite vaut d'être choisie : j'érige mon choix en exemple, je dessine une image de l'homme que je propose à tous. Celui qui triche affirme silencieusement que tricher est acceptable pour tout élève ; celui qui est honnête affirme que l'honnêteté vaut pour tous. Cette formule, proche de l'impératif kantien (« agis de telle sorte que la maxime de ton action puisse être érigée en loi universelle »), donne à chaque acte une dimension universelle : nul ne peut se dire « ce n'est que mon affaire ».
\item \textbf{Discussion.} Cette conception a une force indéniable : elle rend à l'homme sa dignité en refusant les excuses faciles. L'élève qui triche au Probatoire ne peut invoquer valablement la tentation, la pression familiale ou l'exemple des autres : c'est lui qui a signé l'acte, et Sartre a raison de le lui rappeler — assumer sa tricherie est le premier pas vers la probité. Pourtant, la thèse atteint ses limites lorsqu'elle ignore les contraintes réelles : peut-on dire qu'un enfant soldat, qu'un malade mental, qu'un affamé sont « totalement responsables » de leurs actes ? Le droit camerounais, plus nuancé, tient compte du discernement et des circonstances. La responsabilité sartrienne est donc un idéal régulateur exigeant — utile contre la mauvaise foi — mais qu'il faut tempérer par l'attention aux conditions concrètes de la liberté.
\end{enumerate}
\textbf{Barème indicatif (sur 20) :} thèse dégagée (3 pts) ; paradoxe expliqué (4 pts) ; argumentation reconstituée (5 pts) ; portée morale (3 pts) ; discussion critique avec exemple (4 pts) ; présentation (1 pt).
\textbf{Points de vigilance :} ne pas confondre « condamné » avec « puni par la société » ; mobiliser la notion de facticité ; la discussion doit être un vrai débat (force + limite), non une simple approbation.
\end{corrige}

\begin{corrige}[Exercice 11 — Déterminisme et responsabilité]
\begin{enumerate}
\item \textbf{Définition et problème.} Le \cle{déterminisme} est la thèse selon laquelle tout événement, y compris toute décision humaine, est l'effet nécessaire de causes antérieures : hérédité, milieu social, éducation, passions. Il semble exclure la responsabilité : si mon acte était nécessaire, je n'aurais pas pu faire autrement ; or on ne répond que de ce qu'on aurait pu ne pas faire. Le déterminisme transformerait le blâme en absurdité : autant blâmer une machine.
\item \textbf{Thèse (Spinoza).} \textbf{Spinoza} soutient que la liberté humaine est une illusion née de l'ignorance : les hommes « ont conscience de leurs désirs et ignorent les causes qui les déterminent à désirer ». Comme la pierre lancée dans les airs qui, si elle pensait, se croirait libre de sa trajectoire, nous confondons le sentiment de choisir avec le pouvoir de choisir. Illustrations camerounaises : le jeune qui consomme de l'alcool reproduit souvent un modèle familial et festif qu'il n'a pas choisi ; l'adulte violent a fréquemment été un enfant battu, pour qui la violence est la seule langue apprise. Dans les deux cas, l'acte semble « choisi » alors qu'il prolonge une chaîne causale antérieure. Si Spinoza a raison, la responsabilité au sens de la culpabilité s'effondre : il ne reste que des causes à comprendre et à traiter.
\item \textbf{Antithèse et synthèse.} Pourtant, ni la morale ni le droit camerounais ne tirent cette conséquence. Devant le tribunal de première instance, l'accusé est interrogé, il se défend, il plaide : toute la procédure le traite comme un être capable de comprendre, de délibérer et de répondre — c'est-à-dire comme un sujet libre. Et cette pratique n'est pas une simple fiction : une conscience, même influencée, conserve un pouvoir réel d'accepter ou de refuser. L'alcoolique peut entrer en cure ; l'enfant battu peut refuser de battre à son tour — beaucoup le font. Le déterminisme décrit des \emph{penchants}, non des \emph{destins}. La synthèse s'impose alors : reconnaître les causes (pour comprendre, prévenir, atténuer) tout en maintenant l'imputation (pour respecter la dignité de l'agent et protéger la société). Le juge camerounais réalise cette synthèse lorsqu'il condamne tout en retenant les circonstances atténuantes.
\item \textbf{Conclusion personnelle (modèle).} Le déterminisme ne supprime donc pas la responsabilité ; il la déplace et l'éclaire. Il nous apprend l'humilité — nul n'est pur auteur de soi — et l'exigence — nul n'est pur produit du monde. Supprimer la responsabilité au nom des causes serait traiter l'homme comme une chose et lui refuser le respect qu'on doit à un être capable de choisir ; ignorer les causes au nom de la liberté serait une cruauté aveugle. La sagesse, qui est aussi celle de notre droit, consiste à punir l'acte sans haïr l'homme, à comprendre les causes sans excuser le crime, et surtout à agir sur ces causes — par l'école, la famille, la justice sociale — pour que chacun dispose d'une liberté plus grande, donc d'une responsabilité plus pleine.
\end{enumerate}
\textbf{Barème indicatif (sur 20) :} définitions et problème (4 pts) ; thèse déterministe avec Spinoza et exemples (5 pts) ; antithèse-synthèse avec le cas judiciaire (6 pts) ; conclusion personnelle justifiée (4 pts) ; présentation (1 pt).
\textbf{Points de vigilance :} attribuer correctement à Spinoza la thèse de l'illusion de la liberté (et non l'indéterminisme) ; éviter l'écueil de l'excuse totale ; la conclusion doit trancher (le déterminisme \emph{modifie} la responsabilité sans la supprimer) et rester dans les bornes du programme.
\end{corrige}

\newpage

\coursec{Fiche bilan}

\coursec{La liberté et la responsabilité}

\courssub{Carte mentale de la leçon}
\begin{popfigure}
\begin{tikzpicture}[
  mindmap,
  concept color=popblue,
  level 1/.style={level distance=4.5cm, sibling angle=360/6},
  level 2/.style={level distance=3.2cm, sibling angle=360/4},
  every node/.style={concept, execute at begin node=\small\bfseries},
  text width=2.8cm, align=center,
  font=\small
]
\node[concept color=popblue, align=center]{Liberté et\\ Responsabilité}
  child[concept color=poporange] {node[align=center]{1. Définitions\\ clés}
    child {node[align=center]{Liberté : pouvoir\\ de choisir}}
    child {node[align=center]{Responsabilité :\\ répondre de ses actes}}
    child {node[align=center]{Autonomie vs\\ Hétéronomie}}
    child {node {Imputabilité}}}
  child[concept color=popgreen] {node[align=center]{2. Formes de\\ la liberté}
    child {node[align=center]{Liberté\\ d'indifférence}}
    child {node[align=center]{Liberté pour\\ l'autonomie (Kant)}}
    child {node[align=center]{Liberté\\ politique}}
    child {node[align=center]{Libre arbitre\\ (théologique)}}}
  child[concept color=poppurple] {node[align=center]{3. Formes de\\ la responsabilité}
    child {node[align=center]{Responsabilité\\ morale}}
    child {node[align=center]{Responsabilité\\ juridique}}
    child {node {Civile / Pénale}}
    child {node[align=center]{Responsabilité\\ collective}}}
  child[concept color=poppink] {node[align=center]{4. Liberté,\\ fondement de la\\ responsabilité}
    child {node[align=center]{Condition\\ nécessaire}}
    child {node[align=center]{Déterminisme\\ vs Liberté}}
    child {node[align=center]{Sans liberté $\to$\\ pas de sanction juste}}
    child {node[align=center]{Ricœur :\\ imputabilité}}}
  child[concept color=popgold] {node[align=center]{5. Applications\\ Cameroun}
    child {node[align=center]{Code du travail\\ (faute lourde)}}
    child {node[align=center]{Code pénal\\ (irresponsabilité)}}
    child {node[align=center]{Élections\\ (liberté politique)}}
    child {node[align=center]{Coutume vs\\ Droit moderne}}}
  child[concept color=popdark] {node[align=center]{6. Méthode\\ Dissertation /\ Commentaire}
    child {node[align=center]{Problématiser\\ le sujet}}
    child {node[align=center]{Thèses / Antithèses\\ (auteurs)}}
    child {node[align=center]{Exemples concrets\\ (Cameroun)}}
    child {node[align=center]{Synthèse\\ nuancée}}};
\end{tikzpicture}
\legende{Figure 1 -- Vue d'ensemble de la leçon : concepts, formes, lien fondamental, ancrage camerounais et méthode d'examen.}
\end{popfigure}

\courssub{Bilan révisable}
\begin{aretenir}[Bilan révisable -- Liberté et Responsabilité]

\textbf{Définitions clés à connaître par cœur}
\begin{itemize}
  \item \cle{Liberté} : Capacité d'agir ou de ne pas agir selon sa propre volonté ; pouvoir de choisir entre plusieurs possibles. Distinction : \emph{liberté d'indifférence} (choix indifférent au bien/mal, Descartes) vs \emph{liberté pour l'autonomie} (obéissance à sa propre loi rationnelle, Kant).
  \item \cle{Responsabilité} : Obligation de répondre de ses actes devant autrui ou devant sa conscience ; imputabilité d'une action à un sujet moral. Distinction : \emph{responsabilité morale} (for intérieur, conscience) vs \emph{responsabilité juridique} (sanction légale, civile ou pénale).
  \item \cle{Imputabilité} : Qualité de ce qui peut être attribué à un auteur libre et conscient. Conditions : liberté + conscience + absence de contrainte irrésistible.
  \item \cle{Autonomie} (Kant) : Se donner sa propre loi morale ; \cle{Hétéronomie} : Recevoir la loi de l'extérieur (désirs, coutumes, autorité).
  \item \cle{Mauvaise foi} (Sartre) : Fuite devant sa liberté et sa responsabilité en se faisant passer pour un objet déterminé.
\end{itemize}

\textbf{Thèses majeures et auteurs de référence (à citer avec précision)}
\begin{center}
\begin{tabularx}{\linewidth}{|l|X|X|}\hline
\textbf{Auteur / Courant} & \textbf{Thèse centrale} & \textbf{Concept clé} \\ \hline
\rowcolor{popblueL}
Aristote & Action volontaire vs involontaire ; la responsabilité suppose que le principe de l'action est en soi (\emph{archè}). & Volontaire / Involontaire / Ignorance \\ \hline
Descartes & Liberté d'indifférence : la volonté est infinie, l'entendement fini $\to$ possibilité d'erreur et de choix sans raison déterminante. & Libre arbitre / Indifférence \\ \hline
\rowcolor{poporangeL}
Spinoza & Liberté = nécessité comprise ; déterminisme universel. « Libre est ce qui existe par la seule nécessité de sa nature » (\emph{Éthique}, I, Déf. 7). & Déterminisme / Conatus / Servitude \\ \hline
Kant & Autonomie de la volonté : se donner la loi morale à soi-même (impératif catégorique). Liberté = postulat pratique nécessaire à la morale. & Autonomie / Devoir / Postulat \\ \hline
\rowcolor{popgreenL}
Sartre & « L'homme est condamné à être libre » (\emph{L'Être et le Néant}) ; responsabilité absolue, angoisse, fuite en mauvaise foi. & Existence précède essence / Angoisse \\ \hline
Ricœur & Responsabilité comme \emph{imputabilité} : répondre de ses actes comme auteur ; identité narrative et attestation. & Imputabilité / Attestation / Narrativité \\ \hline
\end{tabularx}
\end{center}

\textbf{Méthodes express (Bac -- Dissertation, Commentaire, Application)}
\begin{enumerate}
  \item \textbf{Dissertation} : 
    \begin{enumerate}
      \item Analyser les termes du sujet (définir, distinguer, opposer).
      \item Problématiser (ex: \og La liberté est-elle la condition \emph{unique} de la responsabilité ? \fg{} ou \og Peut-on être responsable sans être libre ? \fg{}).
      \item Construire un plan dialectique : 
        \begin{itemize}
          \item Thèse : La liberté est le fondement nécessaire de la responsabilité (Aristote, Kant, Droit).
          \item Antithèse : Limites de la liberté (déterminismes : Spinoza, psychanalyse, sociologie, neurosciences) ; responsabilité sans liberté pleine (responsabilité civile sans faute, responsabilité collective, Ricœur).
          \item Synthèse : La responsabilité comme attribution progressive (Ricœur) ; distinction responsabilité morale / juridique ; liberté comme horizon de la responsabilité.
        \end{itemize}
      \item Illustrer \textbf{obligatoirement} par des exemples camerounais (Code du travail, Code pénal, vie citoyenne).
    \end{enumerate}
  \item \textbf{Commentaire de texte} : 
    \begin{enumerate}
      \item Introduction : Présenter l'auteur, l'œuvre, la thèse, l'intérêt philosophique, le plan.
      \item Analyse structurée : Suivre le fil du texte, expliquer les arguments, définir les concepts techniques.
      \item Discussion critique : Portée et limites de la thèse ; confrontation avec un autre auteur (ex: Kant vs Spinoza, Sartre vs Ricœur).
      \item Conclusion : Bilan de la discussion, actualité de la thèse, ouverture.
    \end{enumerate}
  \item \textbf{Application concrète (Cas pratique)} : 
    \begin{enumerate}
      \item Identifier la situation de fait (contexte, acteurs, acte).
      \item Qualifier l'acte : Volontaire / Non volontaire / Contraint (force majeure, trouble mental, ordre légitime).
      \item Déterminer le type de responsabilité engagé : Morale (conscience, culpabilité) et/ou Juridique (Civile : réparation ; Pénale : sanction).
      \item Vérifier les causes d'irresponsabilité ou d'atténuation (Code pénal Art. 39-41, minorité).
      \item Conclure sur la sanction, la réparation ou l'acquittement justifié.
    \end{enumerate}
\end{enumerate}

\textbf{Exemples camerounais types à maîtriser (références légales exactes)}
\begin{itemize}
  \item \textbf{Droit du travail (Loi n°92/007 du 14 août 1992 modifiée par Loi n°2024/016)} : 
    \begin{itemize}
      \item Faute lourde (vol, ivresse habituelle, indiscipline grave, absence injustifiée $> 4$ jours) $\to$ Licenciement sans préavis ni indemnité (Art. 39).
      \item Responsabilité de l'employeur : Obligation de sécurité, de paiement des salaires, de formation.
    \end{itemize}
  \item \textbf{Droit pénal (Loi n°2016/007 du 12 juillet 2016 - Code pénal)} : 
    \begin{itemize}
      \item Causes d'irresponsabilité pénale (Art. 39-41) : Trouble mental abolissant le discernement (Art. 39) ; Contrainte physique irrésistible (Art. 40) ; Ordre de la loi / commandement de l'autorité légitime (Art. 41).
      \item Minorité : $< 10$ ans = irresponsabilité absolue (Art. 81) ; $10-18$ ans = responsabilité atténuée, mesures de protection/éducation (Art. 82-83).
      \item Distinction Responsabilité Civile (réparation dommage, Art. 1382 Code Civil inspiration) vs Pénale (sanction trouble à l'ordre public).
    \end{itemize}
  \item \textbf{Vie civique et citoyenneté} : 
    \begin{itemize}
      \item Vote : Liberté politique (choix) + Devoir civique (Art. 1 Constitution, Code électoral).
      \item Code de la route : Liberté de circuler vs Responsabilité de sécurité (permis, assurance, respect signalisation).
      \item Gestion des biens communs : Liberté d'usage vs Responsabilité de préservation (environnement, patrimoine).
      \item Conflit Coutume / Droit moderne : Ex. Héritage, mariage coutumier, chefferies traditionnelles vs Code civil / Code des personnes et de la famille.
    \end{itemize}
\end{itemize}

\end{aretenir}

\courssub{Checklist avant le Bac}
\begin{aretenir}[Checklist avant le Bac -- Liberté et Responsabilité]
\begin{itemize}
  \item $\square$ Je sais définir \cle{liberté} (pouvoir choisir entre plusieurs possibles) et \cle{responsabilité} (obligation de répondre de ses actes) sans les confondre.
  \item $\square$ Je distingue \cle{liberté d'indifférence} (Descartes : choix sans motif déterminant) et \cle{liberté pour l'autonomie} (Kant : obéissance à la loi morale rationnelle).
  \item $\square$ Je distingue \cle{responsabilité morale} (for intérieur, culpabilité, conscience) et \cle{responsabilité juridique} (sanction légale : civile = réparation / pénale = peine).
  \item $\square$ J'explique pourquoi la \cle{liberté est le fondement nécessaire} de la responsabilité (pas de sanction juste sans auteur libre ; \og Nul n'est responsable que de ses actes volontaires \fg{}).
  \item $\square$ Je cite au moins \textbf{deux auteurs} avec leur thèse précise (ex: Kant \emph{autonomie}, Sartre \emph{condamné à être libre}, Ricœur \emph{imputabilité}, Spinoza \emph{déterminisme}, Aristote \emph{volontaire/involontaire}).
  \item $\square$ Je maîtrise le vocabulaire technique : \cle{imputabilité}, \cle{autonomie}, \cle{hétéronomie}, \cle{libre arbitre}, \cle{déterminisme}, \cle{mauvaise foi}, \cle{force majeure}, \cle{trouble mental}.
  \item $\square$ Je donne un \textbf{exemple concret camerounais} (Code du travail : faute lourde ; Code pénal : irresponsabilité ; Vie citoyenne : vote, route) pour illustrer chaque grande notion.
  \item $\square$ Je sais structurer une \textbf{dissertation} (Intro : accroche, définitions, problème, annonce de plan / Développement : 2-3 parties argumentées avec auteurs et exemples / Conclusion : bilan synthétique, ouverture).
  \item $\square$ Je sais structurer un \textbf{commentaire} (Intro : présentation, thèse, intérêt, plan / Analyse : explication linéaire ou thématique / Discussion : portée/limites, confrontation / Conclusion : bilan, actualité).
  \item $\square$ J'évite les pièges classiques : 
    \begin{itemize}
      \item Confondre \emph{liberté} (pouvoir choisir) et \emph{licence} (faire n'importe quoi sans règle).
      \item Nier toute liberté (déterminisme dur non nuancé : Spinoza mal compris, neurosciences réductrices).
      \item Oublier les causes d'irresponsabilité juridique (force majeure, trouble mental, ordre légitime, minorité).
      \item Confondre responsabilité civile (réparation) et pénale (sanction).
    \end{itemize}
  \item $\square$ Je gère mon temps d'examen : 3h pour la dissertation, 3h pour le commentaire (épreuves type Bac) ; je réserve 15 min pour la relecture (orthographe, syntaxe, noms d'auteurs, numéros d'articles).
\end{itemize}
\end{aretenir}

\findecours

\end{document}
